\input{preamble}

% OK, start here.
%
\begin{document}

\title{Cohomologie des champs algébriques}

\maketitle

\phantomsection
\label{section-phantom}

\tableofcontents




\section{Introduction}
\label{section-introduction}

\noindent
Dans ce chapitre, nous traitons de la cohomologie des champs algébriques.
Il s'agit notamment de la cohomologie des faisceaux quasi-cohérents, c'est-à-dire que
nous démontrons des analogues des résultats des chapitres intitulés
``Cohomologie des schémas'' et ``Cohomologie des espaces algébriques''.
Les résultats de ce chapitre diffèrent de ceux
de \cite{LM-B} principalement parce que nous utilisons systématiquement les
« grands sites ». Avant de lire ce chapitre, nous invitons le lecteur à parcourir
le chapitre ``Faisceaux sur les champs algébriques'' afin de se
familiariser avec la terminologie qui y est introduite ; voir
Faisceaux sur les champs, section \ref{stacks-sheaves-section-introduction}.



\section{Conventions et abus de langage}
\label{section-conventions}

\noindent
Nous continuons d'employer les conventions et abus de langage
introduits dans
Propriétés des champs, section \ref{stacks-properties-section-conventions}.











\section{Notations}
\label{section-notation}

\noindent
Différentes topologies. Si nous désignons un champ algébrique par une lettre
calligraphique, telle que $\mathcal{X}, \mathcal{Y}, \mathcal{Z}$, alors les notations
$\mathcal{X}_{Zar}, \mathcal{X}_\etale, \mathcal{X}_{smooth},
\mathcal{X}_{syntomic}, \mathcal{X}_{fppf}$ désignent les sites introduits
dans
Faisceaux sur les champs, définition
\ref{stacks-sheaves-definition-inherited-topologies}.
(Il faut penser à un « grand site ».) Par conséquent, le faisceau structural de
$\mathcal{X}$ est un faisceau sur $\mathcal{X}_{fppf}$.
En revanche, les espaces algébriques et les schémas
sont habituellement désignés par des majuscules romaines, telles que $X, Y, Z$, et, dans ce cas,
$X_\etale$ désigne le petit site étale de $X$ (tel qu'il est
défini dans
Topologies, définition
\ref{topologies-definition-big-small-etale}
ou
Propriétés des espaces, définition
\ref{spaces-properties-definition-etale-site}).
Cette distinction paraît suffisamment claire.

\medskip\noindent
La topologie sous-entendue est la topologie fppf. Ainsi, nous dirons parfois
« faisceau sur $\mathcal{X}$ » ou « faisceau de $\mathcal{O}_\mathcal{X}$-modules »
pour désigner un faisceau sur $\mathcal{X}_{fppf}$ ou un objet de
$\textit{Mod}(\mathcal{X}_{fppf}, \mathcal{O}_\mathcal{X})$.

\medskip\noindent
Si $f : \mathcal{X} \to \mathcal{Y}$ est un morphisme de champs
algébriques, alors les foncteurs $f_*$ et $f^{-1}$ définis sur les préfaisceaux
préservent les faisceaux pour chacune des topologies mentionnées ci-dessus. En particulier,
lorsque nous parlons de l'image directe ou de l'image inverse d'un faisceau, il n'est pas nécessaire
de préciser la topologie considérée. Il n'en va pas de même
lorsque nous calculons les groupes de cohomologie ou les images directes supérieures. Dans ce
cas, nous préciserons toujours la topologie considérée.

\medskip\noindent
Supposons que $f : X \to \mathcal{Y}$ soit un morphisme d'un espace
algébrique $X$ vers un champ algébrique $\mathcal{Y}$. Soit $\mathcal{G}$
un faisceau sur $\mathcal{Y}_\tau$ pour une topologie $\tau$. Alors
$f^{-1}\mathcal{G}$ est un faisceau pour la topologie $\tau$ sur $\mathcal{S}_X$
(le champ algébrique associé à $X$), car, suivant nos conventions, $f$
est en fait un $1$-morphisme $f : \mathcal{S}_X \to \mathcal{Y}$.
Si $\tau = \etale$ ou si la topologie est plus fine, nous écrivons
$f^{-1}\mathcal{G}|_{X_\etale}$
pour désigner la restriction au site étale de $X$ ; voir
Faisceaux sur les champs, section \ref{stacks-sheaves-section-compare}.
Si $\mathcal{G}$ est un $\mathcal{O}_\mathcal{X}$-module, nous écrivons parfois
plutôt
$f^*\mathcal{G}$ et $f^*\mathcal{G}|_{X_\etale}$.




\section{Image inverse des modules quasi-cohérents}
\label{section-pullback}

\noindent
Soit $f : \mathcal{X} \to \mathcal{Y}$ un morphisme de champs algébriques.
Il est très général que les modules quasi-cohérents sur les topos annelés
soient compatibles aux images inverses. En particulier, l'image inverse $f^*$ préserve
les modules quasi-cohérents et nous obtenons un foncteur
$$
f^* :
\QCoh(\mathcal{O}_\mathcal{Y})
\longrightarrow
\QCoh(\mathcal{O}_\mathcal{X}),
$$
voir Faisceaux sur les champs, lemme
\ref{stacks-sheaves-lemma-pullback-quasi-coherent}.
En général, ce foncteur n'est pas exact, mais il l'est si $f$ est plat.

\begin{lemma}
\label{lemma-flat-pullback-quasi-coherent}
Si $f : \mathcal{X} \to \mathcal{Y}$ est un morphisme plat de champs algébriques,
alors $f^* : \QCoh(\mathcal{O}_\mathcal{Y}) \to
\QCoh(\mathcal{O}_\mathcal{X})$ est un foncteur exact.
\end{lemma}

\begin{proof}
Choisissons un schéma $V$ et un morphisme lisse surjectif $V \to \mathcal{Y}$.
Choisissons un schéma $U$ et un morphisme lisse surjectif
$U \to V \times_\mathcal{Y} \mathcal{X}$. Alors $U \to \mathcal{X}$ est
encore lisse et surjectif, comme composé de deux tels morphismes.
Du diagramme commutatif
$$
\xymatrix{
U \ar[d] \ar[r]_{f'} & V \ar[d] \\
\mathcal{X} \ar[r]^f & \mathcal{Y}
}
$$
nous déduisons un diagramme commutatif
$$
\xymatrix{
\QCoh(\mathcal{O}_U) & \QCoh(\mathcal{O}_V) \ar[l] \\
\QCoh(\mathcal{O}_\mathcal{X}) \ar[u] &
\QCoh(\mathcal{O}_\mathcal{Y}) \ar[l] \ar[u]
}
$$
de catégories abéliennes. La démonstration du fait que les deux catégories inférieures de ce
diagramme sont abéliennes a montré que les foncteurs verticaux sont des foncteurs
exacts et fidèles (voir la démonstration de
Faisceaux sur les champs, lemme
\ref{stacks-sheaves-lemma-quasi-coherent-algebraic-stack}).
Puisque $f'$ est un morphisme plat de schémas (d'après notre définition des
morphismes plats de champs algébriques), nous voyons que $(f')^*$ est un
foncteur exact sur les faisceaux quasi-cohérents de $V$. Cela conclut.
\end{proof}

\begin{lemma}
\label{lemma-quasi-coherent-check-exact}
Soit $\mathcal{X}$ un champ algébrique. Soit $I$ un ensemble et,
pour $i \in I$, soit $x_i : U_i \to \mathcal{X}$ un objet
de $\mathcal{X}$. Supposons que $x_i$ soit plat et que
$\coprod x_i : \coprod U_i \to \mathcal{X}$ soit surjectif.
Soit $\varphi : \mathcal{F} \to \mathcal{G}$ une flèche de
$\QCoh(\mathcal{O}_\mathcal{X})$. Notons $\varphi_i$
la restriction de $\varphi$ à $(U_i)_\etale$.
Alors $\varphi$ est injective, resp.\ surjective, resp.\ un isomorphisme
si et seulement si chaque $\varphi_i$ l'est.
\end{lemma}

\begin{proof}
Choisissons un schéma $U$ et un morphisme lisse surjectif
$x : U \to \mathcal{X}$. Nous pouvons considérer, et considérons, $x$
comme un objet de $\mathcal{X}$.
On obtient ainsi une présentation $\mathcal{X} = [U/R]$ par un groupoïde
en espaces $(U, R, s, t, c)$ et, par conséquent, une équivalence
$$
\QCoh(\mathcal{O}_\mathcal{X}) = \QCoh(U, R, s, t, c)
$$
Voir la discussion dans Faisceaux sur les champs, section
\ref{stacks-sheaves-section-quasi-coherent-algebraic-stacks}.
La structure de catégorie abélienne du membre de droite est telle que $\varphi$
est injective, resp.\ surjective, resp.\ un isomorphisme
si et seulement si la restriction $\varphi|_{U_\etale}$ l'est ; voir
Groupoïdes dans les espaces, lemme \ref{spaces-groupoids-lemma-abelian}.

\medskip\noindent
Pour chaque $i$, choisissons un recouvrement étale
$\{W_{i, j} \to V \times_\mathcal{X} U_i\}_{j \in J_i}$
par des schémas. Notons $g_{i, j} : W_{i, j} \to V$ et
$h_{i, j} : W_{i, j} \to U_i$ les flèches évidentes.
Chacun des morphismes de schémas $g_{i, j} : W_{i, j} \to U$ est plat,
et ils sont conjointement surjectifs.
De même, pour chaque $i$ fixé, les morphismes de schémas
$h_{i, j} : W_{i, j} \to U_i$ sont plats et conjointement surjectifs.
D'après Faisceaux sur les champs, lemme \ref{stacks-sheaves-lemma-quasi-coherent},
l'image inverse par $(g_{i, j})_{small}$ de la restriction
$\varphi|_{U_\etale}$ est la restriction $\varphi|_{(W_{i, j})_\etale}$,
et l'image inverse par $(h_{i, j})_{small}$ de la restriction
$\varphi|_{(U_i)_\etale}$ est la restriction $\varphi|_{(W_{i, j})_\etale}$.
L'image inverse des modules quasi-cohérents par un morphisme plat de schémas est exacte,
et l'image inverse par une famille conjointement surjective de morphismes plats de schémas
reflète les applications injectives, resp.\ surjectives, resp.\ bijectives
de modules quasi-cohérents (en fait, cela vaut pour tous les modules,
car on peut vérifier l'exactitude sur les germes). Ainsi,
$$
\varphi|_{U_\etale} \text{ injective}
\Leftrightarrow
\varphi|_{(W_{i, j})_\etale} \text{ injective pour tous }i, j
\Leftrightarrow
\varphi|_{(U_i)_\etale} \text{ injective pour tout }i
$$
Ceci achève la démonstration.
\end{proof}




\section{Images directes supérieures de certaines classes de modules}
\label{section-key}

\noindent
Le lemme suivant fonde notre étude des
images directes supérieures de certaines classes de faisceaux de modules.
Il en existe deux versions : l'une pour la topologie étale et
l'autre pour la topologie fppf.

\begin{lemma}
\label{lemma-general-pushforward}
Soit $\mathcal{M}$ une règle qui associe à tout champ algébrique
$\mathcal{X}$ une sous-catégorie $\mathcal{M}_\mathcal{X}$ de
$\textit{Mod}(\mathcal{X}_\etale, \mathcal{O}_\mathcal{X})$
telle que
\begin{enumerate}
\item $\mathcal{M}_\mathcal{X}$ soit une sous-catégorie faible de Serre
de $\textit{Mod}(\mathcal{X}_\etale, \mathcal{O}_\mathcal{X})$
(voir Homologie, définition \ref{homology-definition-serre-subcategory})
pour tout champ algébrique $\mathcal{X}$,
\item pour tout morphisme lisse de champs algébriques
$f : \mathcal{Y} \to \mathcal{X}$, le foncteur $f^*$ envoie
$\mathcal{M}_\mathcal{X}$ dans $\mathcal{M}_\mathcal{Y}$,
\item si $f_i : \mathcal{X}_i \to \mathcal{X}$ est une famille de morphismes
lisses de champs algébriques telle que
$|\mathcal{X}| = \bigcup |f_i|(|\mathcal{X}_i|)$, alors un objet
$\mathcal{F}$ de
$\textit{Mod}(\mathcal{X}_\etale, \mathcal{O}_\mathcal{X})$
appartient à $\mathcal{M}_\mathcal{X}$ si et seulement si
$f_i^*\mathcal{F}$ appartient à $\mathcal{M}_{\mathcal{X}_i}$ pour tout $i$, et
\item si $f : \mathcal{Y} \to \mathcal{X}$ est un morphisme de champs
algébriques tel que $\mathcal{X}$ et $\mathcal{Y}$ soient représentables
par des schémas affines, alors $R^if_*$ envoie $\mathcal{M}_\mathcal{Y}$
dans $\mathcal{M}_\mathcal{X}$.
\end{enumerate}
Alors, pour tout morphisme quasi-compact et quasi-séparé
$f : \mathcal{Y} \to \mathcal{X}$ de champs algébriques,
$R^if_*$ envoie $\mathcal{M}_\mathcal{Y}$
dans $\mathcal{M}_\mathcal{X}$. (Les images directes supérieures sont calculées pour la topologie
étale.)
\end{lemma}

\begin{proof}
Soit $f : \mathcal{Y} \to \mathcal{X}$ un morphisme quasi-compact et quasi-séparé
de champs algébriques, et soit $\mathcal{F}$ un objet de
$\mathcal{M}_\mathcal{Y}$. Choisissons un morphisme lisse surjectif
$\mathcal{U} \to \mathcal{X}$, où $\mathcal{U}$ est représentable par
un schéma. D'après
Faisceaux sur les champs, lemme
\ref{stacks-sheaves-lemma-base-change-higher-direct-images}
la prise des images directes supérieures commute au changement de base.
L'hypothèse (2) montre que l'image inverse de $\mathcal{F}$
sur $\mathcal{U} \times_\mathcal{X} \mathcal{Y}$ appartient à
$\mathcal{M}_{\mathcal{U} \times_\mathcal{X} \mathcal{Y}}$
car la projection
$\mathcal{U} \times_\mathcal{X} \mathcal{Y} \to \mathcal{Y}$
est lisse comme changement de base d'un morphisme lisse. Ainsi, (3) montre que nous pouvons
remplacer $\mathcal{Y} \to \mathcal{X}$ par la projection
$\mathcal{U} \times_\mathcal{X} \mathcal{Y} \to \mathcal{U}$.
Autrement dit, nous pouvons supposer que $\mathcal{X}$
est représentable par un schéma.
En appliquant encore (3), nous voyons que la question est locale sur
$\mathcal{X}$ pour la topologie de Zariski ; nous pouvons donc supposer que $\mathcal{X}$ est représentable par
un schéma affine. Comme $f$ est quasi-compact, il s'ensuit que
$\mathcal{Y}$ est également quasi-compacte. Nous pouvons donc choisir un morphisme lisse surjectif
$g : \mathcal{V} \to \mathcal{Y}$, où $\mathcal{V}$ est représentable
par un schéma affine.

\medskip\noindent
Dans cette situation, nous disposons de la suite spectrale
$$
E_2^{p, q} = R^q(f \circ g_p)_*g_p^*\mathcal{F}
\Rightarrow
R^{p + q}f_*\mathcal{F}
$$
de
Faisceaux sur les champs, proposition
\ref{stacks-sheaves-proposition-smooth-covering-compute-direct-image}.
Rappelons qu'il s'agit d'une suite spectrale du premier quadrant ; nous pouvons donc
appliquer la dernière partie de Homologie, lemme \ref{homology-lemma-first-quadrant-ss}.
Notons que les morphismes
$$
g_p : \mathcal{V}_p =
\mathcal{V} \times_\mathcal{Y} \ldots \times_\mathcal{Y} \mathcal{V}
\longrightarrow
\mathcal{Y}
$$
sont lisses comme composés de changements de base du morphisme lisse $g$.
Ainsi, les faisceaux $g_p^*\mathcal{F}$ appartiennent à
$\mathcal{M}_{\mathcal{V}_p}$ d'après (2). Il suffit donc de démontrer que les
images directes supérieures des objets de $\mathcal{M}_{\mathcal{V}_p}$ par
les morphismes
$$
\mathcal{V}_p =
\mathcal{V} \times_\mathcal{Y} \ldots \times_\mathcal{Y} \mathcal{V}
\longrightarrow
\mathcal{X}
$$
appartiennent à $\mathcal{M}_\mathcal{X}$. Les champs algébriques $\mathcal{V}_p$
sont quasi-compacts et quasi-séparés d'après
Morphismes de champs, lemme
\ref{stacks-morphisms-lemma-quasi-compact-quasi-separated-permanence}.
Bien entendu, chaque $\mathcal{V}_p$ est représentable par un espace algébrique
(la diagonale du champ algébrique $\mathcal{Y}$ est représentable
par des espaces algébriques). Nous sommes donc ramenés au cas où
$\mathcal{Y}$ est représentable par un espace algébrique et $\mathcal{X}$
par un schéma affine.

\medskip\noindent
Lorsque $\mathcal{Y}$ est représentable par un espace
algébrique et $\mathcal{X}$ par un schéma affine, choisissons
de nouveau un morphisme lisse surjectif $\mathcal{V} \to \mathcal{Y}$, où
$\mathcal{V}$ est représentable par un schéma affine. En reprenant
l'argument précédent, nous sommes de nouveau ramenés aux morphismes
$\mathcal{V}_p \to \mathcal{X}$. Mais, dans la situation présente, les champs
algébriques $\mathcal{V}_p$ sont représentables par des schémas quasi-compacts et quasi-séparés
(car la diagonale d'un espace algébrique est représentable par des
schémas).

\medskip\noindent
Nous pouvons donc supposer que $\mathcal{Y}$ est représentable par un schéma et
$\mathcal{X}$ par un schéma affine. Choisissons (encore)
un morphisme lisse surjectif $\mathcal{V} \to \mathcal{Y}$, où
$\mathcal{V}$ est représentable par un schéma affine. Dans ce cas, tous
les champs algébriques $\mathcal{V}_p$ sont représentables par des schémas
séparés (car la diagonale d'un schéma est séparée).

\medskip\noindent
Nous pouvons donc supposer que $\mathcal{Y}$ est représentable par un schéma séparé et
$\mathcal{X}$ par un schéma affine. Choisissons (une fois encore)
un morphisme lisse surjectif $\mathcal{V} \to \mathcal{Y}$, où
$\mathcal{V}$ est représentable par un schéma affine. Dans ce cas, tous
les champs algébriques $\mathcal{V}_p$ sont représentables par des schémas affines
(car la diagonale d'un schéma séparé est une immersion fermée, donc affine),
et ce cas est couvert par l'hypothèse (4).
Ceci achève la démonstration.
\end{proof}

\noindent
Voici la version pour la topologie fppf.

\begin{lemma}
\label{lemma-general-pushforward-fppf}
Soit $\mathcal{M}$ une règle qui associe à tout champ algébrique
$\mathcal{X}$ une sous-catégorie $\mathcal{M}_\mathcal{X}$ de
$\textit{Mod}(\mathcal{O}_\mathcal{X})$
telle que
\begin{enumerate}
\item $\mathcal{O}_\mathcal{X}$ soit une sous-catégorie faible de Serre
de $\textit{Mod}(\mathcal{O}_\mathcal{X})$
pour tout champ algébrique $\mathcal{X}$,
\item pour tout morphisme lisse de champs algébriques
$f : \mathcal{Y} \to \mathcal{X}$, le foncteur $f^*$ envoie
$\mathcal{M}_\mathcal{X}$ dans $\mathcal{M}_\mathcal{Y}$,
\item si $f_i : \mathcal{X}_i \to \mathcal{X}$ est une famille de morphismes
lisses de champs algébriques telle que
$|\mathcal{X}| = \bigcup |f_i|(|\mathcal{X}_i|)$, alors un objet
$\mathcal{F}$ de $\textit{Mod}(\mathcal{O}_\mathcal{X})$
appartient à $\mathcal{M}_\mathcal{X}$ si et seulement si
$f_i^*\mathcal{F}$ appartient à $\mathcal{M}_{\mathcal{X}_i}$ pour tout $i$, et
\item si $f : \mathcal{Y} \to \mathcal{X}$ est un morphisme de champs
algébriques tel que $\mathcal{X}$ et $\mathcal{Y}$ soient représentables
par des schémas affines, alors $R^if_*$ envoie $\mathcal{M}_\mathcal{Y}$
dans $\mathcal{M}_\mathcal{X}$.
\end{enumerate}
Alors, pour tout morphisme quasi-compact et quasi-séparé
$f : \mathcal{Y} \to \mathcal{X}$ de champs algébriques,
$R^if_*$ envoie $\mathcal{M}_\mathcal{Y}$
dans $\mathcal{M}_\mathcal{X}$. (Les images directes supérieures sont calculées pour la topologie
fppf.)
\end{lemma}

\begin{proof}
Identique à la démonstration du lemme \ref{lemma-general-pushforward}.
\end{proof}


\section{Modules localement quasi-cohérents}
\label{section-locally-quasi-coherent}

\noindent
Soit $\mathcal{X}$ un champ algébrique. Soit $\mathcal{F}$ un préfaisceau
de $\mathcal{O}_\mathcal{X}$-modules. On peut se demander si $\mathcal{F}$
est {\it localement quasi-cohérent} ; voir
Faisceaux sur les champs, définition
\ref{stacks-sheaves-definition-locally-quasi-coherent}.
En bref, cela signifie que $\mathcal{F}$ est un $\mathcal{O}_\mathcal{X}$-module
pour la topologie étale tel que, pour tout morphisme $f : U \to \mathcal{X}$,
la restriction $f^*\mathcal{F}|_{U_\etale}$ soit quasi-cohérente
sur $U_\etale$. (La définition proprement dite est légèrement différente, mais
équivalente.) Un fait utile est que
$$
\textit{LQCoh}(\mathcal{O}_\mathcal{X}) \subset
\textit{Mod}(\mathcal{X}_\etale, \mathcal{O}_\mathcal{X})
$$
est une sous-catégorie faible de Serre ; voir
Faisceaux sur les champs, lemme \ref{stacks-sheaves-lemma-lqc-colimits}.

\begin{lemma}
\label{lemma-check-lqc-on-etale-covering}
Soit $\mathcal{X}$ un champ algébrique. Soit
$f_j : \mathcal{X}_j \to \mathcal{X}$ une famille de morphismes
lisses de champs algébriques telle que
$|\mathcal{X}| =\bigcup |f_j|(|\mathcal{X}_j|)$.
Soit $\mathcal{F}$ un faisceau de $\mathcal{O}_\mathcal{X}$-modules
sur $\mathcal{X}_\etale$. Si chaque $f_j^{-1}\mathcal{F}$
est localement quasi-cohérent, alors $\mathcal{F}$ l'est aussi.
\end{lemma}

\begin{proof}
Nous pouvons remplacer chacun des champs algébriques $\mathcal{X}_j$ par
un schéma $U_j$ (en utilisant le fait que tout champ algébrique possède un recouvrement lisse par
un schéma et que les composés de morphismes lisses sont lisses ; voir
Morphismes de champs, lemme \ref{stacks-morphisms-lemma-composition-smooth}).
L'image inverse de $\mathcal{F}$ sur $(\Sch/U_j)_\etale$ est encore
localement quasi-cohérente ; voir
Faisceaux sur les champs, lemme \ref{stacks-sheaves-lemma-pullback-lqc}.
Alors $f = \coprod f_j : U = \coprod U_j \to \mathcal{X}$ est un morphisme
lisse surjectif. Soit $x$ un objet de $\mathcal{X}$. D'après
Faisceaux sur les champs, lemme
\ref{stacks-sheaves-lemma-surjective-flat-locally-finite-presentation}
il existe un recouvrement étale $\{x_i \to x\}_{i \in I}$
tel que chaque $x_i$ se relève en un objet $u_i$ de $(\Sch/U)_\etale$.
Cela signifie simplement que $x$, $x_i$ sont au-dessus de schémas $V$, $V_i$, que
$\{V_i \to V\}$ est un recouvrement étale et que $x_i$ provient
d'un morphisme $u_i : V_i \to U$. La restriction
$x_i^*\mathcal{F}|_{V_{i, \etale}}$ est égale à la restriction
de $f^*\mathcal{F}$ à $V_{i, \etale}$ ; voir
Faisceaux sur les champs, lemme \ref{stacks-sheaves-lemma-comparison}.
Ainsi, $x^*\mathcal{F}|_{V_\etale}$
est un faisceau sur le petit site étale de $V$ dont la restriction
à $V_{i, \etale}$ est quasi-cohérente pour tout $i$.
Il est donc quasi-cohérent, comme souhaité, par exemple d'après
Propriétés des espaces, lemme
\ref{spaces-properties-lemma-characterize-quasi-coherent}.
\end{proof}

\begin{lemma}
\label{lemma-pushforward-locally-quasi-coherent}
Soit $f : \mathcal{X} \to \mathcal{Y}$ un morphisme quasi-compact et
quasi-séparé de champs algébriques. Soit
$\mathcal{F}$ un $\mathcal{O}_\mathcal{X}$-module localement quasi-cohérent
sur $\mathcal{X}_\etale$.
Alors $R^if_*\mathcal{F}$ (calculé pour la topologie étale) est
localement quasi-cohérent sur $\mathcal{Y}_\etale$.
\end{lemma}

\begin{proof}
Nous allons utiliser le
lemme \ref{lemma-general-pushforward}
pour le démontrer. Vérifions ses hypothèses (1) -- (4).
Les parties (1) et (2) résultent de
Faisceaux sur les champs, lemme \ref{stacks-sheaves-lemma-lqc-colimits}.
La partie (3) résulte du
lemme \ref{lemma-check-lqc-on-etale-covering}.
Il suffit donc de démontrer (4).

\medskip\noindent
Supposons que $f : \mathcal{X} \to \mathcal{Y}$ soit un morphisme de champs algébriques
tel que $\mathcal{X}$ et $\mathcal{Y}$ soient représentables par des schémas
affines $X$ et $Y$. Choisissons un objet quelconque $y$ de $\mathcal{Y}$ au-dessus d'un
schéma $V$. Pour plus de clarté, notons $\mathcal{V} = (\Sch/V)_{fppf}$ le
champ algébrique correspondant à $V$. Considérons le diagramme cartésien
$$
\xymatrix{
\mathcal{Z} \ar[d] \ar[r]_g \ar[d]_{f'} & \mathcal{X} \ar[d]^f \\
\mathcal{V} \ar[r]^y & \mathcal{Y}
}
$$
Ainsi, $\mathcal{Z}$ est représentable par le schéma $Z = V \times_Y X$
et $f'$ est quasi-compact et séparé (même affine). D'après
Faisceaux sur les champs, lemme
\ref{stacks-sheaves-lemma-compare-representable-morphism-cohomology}
nous avons
$$
R^if_*\mathcal{F}|_{V_\etale} =
R^if'_{small, *}\big(g^*\mathcal{F}|_{Z_\etale}\big)
$$
Le membre de droite est un faisceau quasi-cohérent sur $V_\etale$ d'après
Cohomologie des espaces, lemme
\ref{spaces-cohomology-lemma-higher-direct-image}.
Il s'ensuit que le membre de gauche est quasi-cohérent, ce qu'il
fallait démontrer.
\end{proof}

\begin{lemma}
\label{lemma-check-lqc-on-flat-covering}
Soit $\mathcal{X}$ un champ algébrique. Soit
$f_j : \mathcal{X}_j \to \mathcal{X}$ une famille de morphismes plats
et localement de présentation finie de champs algébriques telle que
$|\mathcal{X}| =\bigcup |f_j|(|\mathcal{X}_j|)$.
Soit $\mathcal{F}$ un faisceau de $\mathcal{O}_\mathcal{X}$-modules
sur $\mathcal{X}_{fppf}$. Si chaque $f_j^{-1}\mathcal{F}$
est localement quasi-cohérent, alors $\mathcal{F}$ l'est aussi.
\end{lemma}

\begin{proof}
Supposons d'abord qu'il existe un morphisme $a : \mathcal{U} \to \mathcal{X}$
surjectif, plat, localement de présentation finie, quasi-compact
et quasi-séparé, tel que $a^*\mathcal{F}$ soit localement quasi-cohérent.
Il existe alors une suite exacte
$$
0 \to \mathcal{F} \to a_*a^*\mathcal{F} \to b_*b^*\mathcal{F}
$$
où $b$ est le morphisme
$b : \mathcal{U} \times_\mathcal{X} \mathcal{U} \to \mathcal{X}$ ; voir
Faisceaux sur les champs, proposition
\ref{stacks-sheaves-proposition-exactness-cech-complex} et
lemme \ref{stacks-sheaves-lemma-surjective-flat-locally-finite-presentation}.
De plus, l'image inverse $b^*\mathcal{F}$ est l'image inverse de $a^*\mathcal{F}$
par l'une des projections ; elle est donc localement quasi-cohérente
(Faisceaux sur les champs, lemme \ref{stacks-sheaves-lemma-pullback-lqc}).
Les modules $a_*a^*\mathcal{F}$ et $b_*b^*\mathcal{F}$ sont localement
quasi-cohérents d'après le lemme \ref{lemma-pushforward-locally-quasi-coherent}.
(Notons que $a_*$ et $b_*$ ne dépendent pas de la topologie
utilisée pour les calculer.)
Nous en concluons que $\mathcal{F}$ est localement quasi-cohérent ; voir
Faisceaux sur les champs, lemme \ref{stacks-sheaves-lemma-lqc-colimits}.

\medskip\noindent
Nous allons ramener la démonstration du cas général à la
situation du premier paragraphe. Soit $x$ un objet de $\mathcal{X}$
au-dessus du schéma $U$. Nous devons montrer que
$\mathcal{F}|_{U_\etale}$ est un $\mathcal{O}_U$-module quasi-cohérent.
Il suffit de le faire localement sur $U$ pour la topologie de Zariski ; nous pouvons donc
supposer $U$ affine. D'après
Morphismes de champs, lemme
\ref{stacks-morphisms-lemma-surjective-family-flat-locally-finite-presentation}
il existe un recouvrement fppf $\{a_i : U_i \to U\}$ tel que
chaque $x \circ a_i$ se factorise par un certain $f_j$. Ainsi, $a_i^*\mathcal{F}$
est localement quasi-cohérent sur $(\Sch/U_i)_{fppf}$. En raffinant
le recouvrement, nous pouvons supposer que $\{U_i \to U\}_{i = 1, \ldots, n}$
est un recouvrement fppf standard. Alors $x^*\mathcal{F}$ est un
module fppf sur $(\Sch/U)_{fppf}$ dont l'image inverse par le morphisme
$a : U_1 \amalg \ldots \amalg U_n \to U$ est localement quasi-cohérente.
Le premier paragraphe montre donc que $x^*\mathcal{F}$ est localement
quasi-cohérent, ce qui implique assurément que $\mathcal{F}|_{U_\etale}$
est quasi-cohérent.
\end{proof}






\section{Morphismes de comparaison plats}
\label{section-flat-comparison}

\noindent
Soit $\mathcal{X}$ un champ algébrique et soit $\mathcal{F}$ un objet
de $\textit{Mod}(\mathcal{X}_\etale, \mathcal{O}_\mathcal{X})$.
Étant donné un objet $x$ de $\mathcal{X}$ au-dessus du schéma $U$, la
restriction $\mathcal{F}|_{U_\etale}$ est la restriction de
$x^{-1}\mathcal{F}$ au petit site étale de $U$ ; voir
Faisceaux sur les champs, définition \ref{stacks-sheaves-definition-pullback}.
Soit ensuite $\varphi : x \to x'$ un morphisme de $\mathcal{X}$ au-dessus
d'un morphisme de schémas $f : U \to U'$. Nous avons donc un diagramme $2$-commutatif
$$
\xymatrix{
U \ar[rd]_x \ar[rr]_f & & U' \ar[ld]^{x'} \\
& \mathcal{X}
}
$$
À $\varphi$ est associé un morphisme de comparaison entre les restrictions
\begin{equation}
\label{equation-comparison-modules}
c_\varphi :
f_{small}^*(\mathcal{F}|_{U'_\etale})
\longrightarrow
\mathcal{F}|_{U_\etale}
\end{equation}
voir Faisceaux sur les champs, équation
(\ref{stacks-sheaves-equation-comparison-modules}).
Dans cette situation, nous pouvons considérer la propriété suivante
de $\mathcal{F}$.

\begin{definition}
\label{definition-flat-base-change}
Soit $\mathcal{X}$ un champ algébrique et soit $\mathcal{F}$ dans
$\textit{Mod}(\mathcal{X}_\etale, \mathcal{O}_\mathcal{X})$.
Nous disons que $\mathcal{F}$ possède la {\it propriété de changement de base plat}\footnote{Cette
terminologie n'est peut-être pas standard.}
si et seulement si $c_\varphi$ est un isomorphisme dès que $f$ est plat.
\end{definition}

\noindent
Voici un lemme donnant quelques propriétés de cette notion.

\begin{lemma}
\label{lemma-check-flat-comparison-on-etale-covering}
Soit $\mathcal{X}$ un champ algébrique. Soit $\mathcal{F}$
un $\mathcal{O}_\mathcal{X}$-module sur $\mathcal{X}_\etale$.
\begin{enumerate}
\item Si $\mathcal{F}$ possède la propriété de changement de base plat, alors, pour tout morphisme
$g : \mathcal{Y} \to \mathcal{X}$ de champs algébriques, son
image inverse $g^*\mathcal{F}$ la possède aussi.
\item La sous-catégorie pleine de
$\textit{Mod}(\mathcal{X}_\etale, \mathcal{O}_\mathcal{X})$
formée des modules possédant la propriété de changement de base plat
est une sous-catégorie faible de Serre.
\item Soit $f_i : \mathcal{X}_i \to \mathcal{X}$ une famille de morphismes
lisses de champs algébriques telle que
$|\mathcal{X}| = \bigcup_i |f_i|(|\mathcal{X}_i|)$. Si chaque
$f_i^*\mathcal{F}$ possède la propriété de changement de base plat, alors
$\mathcal{F}$ la possède aussi.
\item La catégorie des $\mathcal{O}_\mathcal{X}$-modules
sur $\mathcal{X}_\etale$ possédant la propriété de changement de base plat
admet des limites inductives, qui coïncident avec celles de
$\textit{Mod}(\mathcal{X}_\etale, \mathcal{O}_\mathcal{X})$.
\item Si $\mathcal{F}$ et $\mathcal{G}$ sont dans
$\textit{Mod}(\mathcal{X}_\etale, \mathcal{O}_\mathcal{X})$
et possèdent la propriété de changement de base plat, alors le produit tensoriel
$\mathcal{F} \otimes_{\mathcal{O}_\mathcal{X}} \mathcal{G}$
possède cette propriété.
\item Si $\mathcal{F}$ et $\mathcal{G}$ sont dans
$\textit{Mod}(\mathcal{X}_\etale, \mathcal{O}_\mathcal{X})$
avec $\mathcal{F}$ de présentation finie et $\mathcal{G}$ possédant
la propriété de changement de base plat, alors le faisceau
$\SheafHom_{\mathcal{O}_\mathcal{X}}(\mathcal{F}, \mathcal{G})$
possède la propriété de changement de base plat.
\end{enumerate}
\end{lemma}

\begin{proof}
Soit $g : \mathcal{Y} \to \mathcal{X}$ comme en (1).
Soit $y$ un objet de $\mathcal{Y}$ au-dessus d'un schéma $V$. D'après
Faisceaux sur les champs, lemme \ref{stacks-sheaves-lemma-comparison}
nous avons
$(g^*\mathcal{F})|_{V_\etale} = \mathcal{F}|_{V_\etale}$.
De plus, un morphisme de comparaison pour le faisceau $g^*\mathcal{F}$ sur $\mathcal{Y}$
est un cas particulier d'un morphisme de comparaison pour le faisceau $\mathcal{F}$ sur
$\mathcal{X}$ ; voir
Faisceaux sur les champs, lemme \ref{stacks-sheaves-lemma-comparison}.
Ainsi, (1) est clair.

\medskip\noindent
Démonstration de (2). Nous utilisons la caractérisation des sous-catégories faibles de Serre de
Homologie, lemme \ref{homology-lemma-characterize-weak-serre-subcategory}.
Les noyaux et conoyaux des
morphismes entre faisceaux possédant la propriété de changement de base plat
possèdent eux aussi cette propriété. C'est clair, car
$f_{small}^*$ est exact pour un morphisme plat de schémas et les
foncteurs de restriction $(-)|_{U_\etale}$ sont exacts (puisque nous
travaillons avec la topologie étale). Enfin, si
$0 \to \mathcal{F}_1 \to \mathcal{F}_2 \to \mathcal{F}_3 \to 0$
est une suite exacte courte dans
$\textit{Mod}(\mathcal{X}_\etale, \mathcal{O}_\mathcal{X})$
et que les deux faisceaux extrêmes possèdent la propriété de changement de base plat, alors
celui du milieu la possède aussi, toujours par l'exactitude de
$f_{small}^*$ et des foncteurs de restriction (ainsi que par le lemme des cinq).

\medskip\noindent
Démonstration de (3).
Soit $f_i : \mathcal{X}_i \to \mathcal{X}$ une famille conjointement surjective de
morphismes lisses de champs algébriques et supposons que chaque $f_i^*\mathcal{F}$
possède la propriété de changement de base plat. D'après la partie (1), la définition
d'un champ algébrique et le fait que les composés de morphismes lisses
sont lisses (voir
Morphismes de champs, lemme \ref{stacks-morphisms-lemma-composition-smooth})
nous pouvons supposer que chaque $\mathcal{X}_i$ est représentable par un schéma.
Soit $\varphi : x \to x'$ un morphisme de $\mathcal{X}$ au-dessus
d'un morphisme plat de schémas $a : U \to U'$. D'après
Faisceaux sur les champs, lemme
\ref{stacks-sheaves-lemma-surjective-flat-locally-finite-presentation}
il existe une famille conjointement surjective de morphismes étales
$U'_i \to U'$ telle que $U'_i \to U' \to \mathcal{X}$ se factorise par
$\mathcal{X}_i$. Nous obtenons ainsi des diagrammes commutatifs
$$
\xymatrix{
U_i = U \times_{U'} U_i' \ar[r]_-{a_i} \ar[d] &
U_i' \ar[r]_{x_i'} \ar[d] & \mathcal{X}_i \ar[d]^{f_i} \\
U \ar[r]^a & U' \ar[r]^{x'} & \mathcal{X}
}
$$
Notons que chaque $a_i$ est un morphisme plat de schémas, comme changement de base de $a$.
Notons $\psi_i : x_i \to x'_i$ le morphisme de $\mathcal{X}_i$ au-dessus de
$a_i$ et de but $x_i'$. Par hypothèse, le morphisme de comparaison
$c_{\psi_i} :
(a_i)_{small}^*\big(f_i^*\mathcal{F}|_{(U'_i)_\etale}\big)
\to f_i^*\mathcal{F}|_{(U_i)_\etale}$ est un isomorphisme.
Comme les flèches verticales $U_i' \to U'$ et $U_i \to U$ sont étales,
les faisceaux $f_i^*\mathcal{F}|_{(U_i')_\etale}$ et
$f_i^*\mathcal{F}|_{(U_i)_\etale}$ sont les restrictions de
$\mathcal{F}|_{U'_\etale}$ et $\mathcal{F}|_{U_\etale}$,
et le morphisme $c_{\psi_i}$ est la restriction de $c_\varphi$ à
$(U_i)_\etale$ ; voir
Faisceaux sur les champs, lemme \ref{stacks-sheaves-lemma-comparison}.
Puisque $\{U_i \to U\}$ est un recouvrement étale, il s'ensuit
que le morphisme de comparaison $c_\varphi$ est un isomorphisme, ce qu'il
fallait démontrer.

\medskip\noindent
Démonstration de (4). Soit $\mathcal{I} \to
\textit{Mod}(\mathcal{X}_\etale, \mathcal{O}_\mathcal{X})$,
$i \mapsto \mathcal{F}_i$ un diagramme, et supposons que chaque $\mathcal{F}_i$
possède la propriété de changement de base plat. Soit $\varphi : x \to x'$ un morphisme
de $\mathcal{X}$ au-dessus du morphisme plat de schémas $f : U \to U'$.
Rappelons que $\colim_i \mathcal{F}_i$ est le faisceau associé à la limite inductive des préfaisceaux.
Comme nous utilisons la topologie étale, il est clair que
$$
(\colim_i \mathcal{F}_i)|_{U_\etale} =
\colim_i {\mathcal{F}_i}|_{U_\etale}
$$
et de même pour la restriction à $U'_\etale$. Ainsi,
\begin{align*}
f_{small}^*((\colim_i \mathcal{F}_i)|_{U'_\etale})
& =
f_{small}^*(\colim_i {\mathcal{F}_i}|_{U'_\etale}) \\
& =
\colim_i f_{small}^*({\mathcal{F}_i}|_{U'_\etale}) \\
& \xrightarrow{\colim c_\varphi}
\colim_i \mathcal{F}_i|_{U_\etale} \\
& =
(\colim_i \mathcal{F}_i)|_{U_\etale}
\end{align*}
Pour la deuxième égalité, nous avons utilisé que $f_{small}^*$ commute aux limites inductives
(en tant qu'adjoint à gauche). La flèche est un isomorphisme puisque chaque
$\mathcal{F}_i$ possède la propriété de changement de base plat. La limite inductive
possède donc cette propriété, et (4) est vraie.

\medskip\noindent
La partie (5) vaut parce que les produits tensoriels commutent aux images inverses ; voir
Modules sur les sites, lemme \ref{sites-modules-lemma-tensor-product-pullback}.
Nous omettons les détails.

\medskip\noindent
Soient $\mathcal{F}$ et $\mathcal{G}$ comme en (6). Puisque $\mathcal{F}$
est quasi-cohérent, il possède la propriété de changement de base plat d'après
Faisceaux sur les champs, lemme \ref{stacks-sheaves-lemma-quasi-coherent}.
Soit $\varphi : x \to x'$ un morphisme de $\mathcal{X}$ au-dessus
du morphisme plat de schémas $f : U \to U'$.
Comme nous utilisons la topologie étale, nous avons
$$
\SheafHom_{\mathcal{O}_\mathcal{X}}(\mathcal{F}, \mathcal{G})|_{U_\etale} =
\SheafHom_{\mathcal{O}_U}(\mathcal{F}|_{U_\etale}, \mathcal{G}|_{U_\etale})
$$
et de même pour la restriction à $U'_\etale$ (nous omettons les détails). Ainsi,
\begin{align*}
f_{small}^*(
\SheafHom_{\mathcal{O}_\mathcal{X}}(\mathcal{F}, \mathcal{G})|_{U'_\etale})
& =
f_{small}^*(
\SheafHom_{\mathcal{O}_{U'}}(
\mathcal{F}|_{U'_\etale}, \mathcal{G}|_{U'_\etale})) \\
& =
\SheafHom_{\mathcal{O}_{U'}}(
f_{small}^*(\mathcal{F}|_{U'_\etale}),
f_{small}^*(\mathcal{G}|_{U'_\etale})) \\
& \xrightarrow{c_\varphi}
\SheafHom_{\mathcal{O}_U}(\mathcal{F}|_{U_\etale}, \mathcal{G}|_{U_\etale}) \\
& =
\SheafHom_{\mathcal{O}_\mathcal{X}}(\mathcal{F}, \mathcal{G})|_{U_\etale}
\end{align*}
Ici, la deuxième égalité est donnée par
Modules sur les sites, lemme \ref{sites-modules-lemma-pullback-internal-hom}
qui utilise que $f : U \to U'$ est plat et que, par conséquent, le morphisme
de sites annelés $f_{small}$ est plat lui aussi.
La flèche est un isomorphisme puisque $\mathcal{F}$ et $\mathcal{G}$
possèdent toutes deux la propriété de changement de base plat. Ainsi, notre $\SheafHom$
possède lui aussi cette propriété, comme souhaité.
\end{proof}

\begin{lemma}
\label{lemma-flat-comparison}
Soit $f : \mathcal{X} \to \mathcal{Y}$ un morphisme quasi-compact et
quasi-séparé de champs algébriques. Soit
$\mathcal{F}$ un objet de
$\textit{Mod}(\mathcal{X}_\etale, \mathcal{O}_\mathcal{X})$
qui soit localement quasi-cohérent et possède la propriété de changement de base plat.
Alors chaque $R^if_*\mathcal{F}$ (calculé pour la topologie étale)
possède la propriété de changement de base plat.
\end{lemma}

\begin{proof}
Nous allons utiliser le
lemme \ref{lemma-general-pushforward}
pour le démontrer. Pour tout champ algébrique $\mathcal{X}$, notons
$\textit{LQCoh}^{fbc}(\mathcal{O}_\mathcal{X})$ la sous-catégorie pleine de
$\textit{Mod}(\mathcal{X}_\etale, \mathcal{O}_\mathcal{X})$
formée des faisceaux localement quasi-cohérents possédant la propriété de
changement de base plat. Une fois vérifiées les conditions (1) -- (4) du
lemme \ref{lemma-general-pushforward}
le lemme en résultera. Les propriétés (1), (2) et (3) résultent de
Faisceaux sur les champs, lemmes \ref{stacks-sheaves-lemma-pullback-lqc} et
\ref{stacks-sheaves-lemma-lqc-colimits}
ainsi que des
lemmes \ref{lemma-check-lqc-on-etale-covering} et
\ref{lemma-check-flat-comparison-on-etale-covering}.
Il suffit donc de démontrer la partie (4).

\medskip\noindent
Supposons que $f : \mathcal{X} \to \mathcal{Y}$ soit un morphisme de champs algébriques
tel que $\mathcal{X}$ et $\mathcal{Y}$ soient représentables par des schémas
affines $X$ et $Y$. Dans ce cas, supposons que
$\psi : y \to y'$ soit un morphisme de $\mathcal{Y}$ au-dessus
d'un morphisme plat de schémas $b : V \to V'$. Pour plus de clarté, notons
$\mathcal{V} = (\Sch/V)_{fppf}$ et $\mathcal{V}' = (\Sch/V')_{fppf}$
les champs algébriques correspondants. Considérons le diagramme
de champs algébriques
$$
\xymatrix{
\mathcal{Z} \ar[d]_{f''} \ar[r]_a &
\mathcal{Z}' \ar[r]_{x'} \ar[d]_{f'} & \mathcal{X} \ar[d]^f \\
\mathcal{V} \ar[r]^b & \mathcal{V}' \ar[r]^{y'} & \mathcal{Y}
}
$$
dont les deux carrés sont cartésiens. Comme $f$ est représentable par des schémas
(et quasi-compact et séparé -- même affine), nous voyons que $\mathcal{Z}$ et
$\mathcal{Z}'$ sont représentables par des schémas $Z$ et $Z'$ et qu'en
fait $Z = V \times_{V'} Z'$. Puisque $\mathcal{F}$ possède la propriété de
changement de base plat, nous voyons que
$$
a_{small}^*\big(\mathcal{F}|_{Z'_\etale}\big)
\longrightarrow
\mathcal{F}|_{Z_\etale}
$$
est un isomorphisme. De plus,
$$
R^if_*\mathcal{F}|_{V'_\etale} =
R^i(f')_{small, *}\big(\mathcal{F}|_{Z'_\etale}\big)
$$
et
$$
R^if_*\mathcal{F}|_{V_\etale} =
R^i(f'')_{small, *}\big(\mathcal{F}|_{Z_\etale}\big)
$$
d'après
Faisceaux sur les champs, lemme
\ref{stacks-sheaves-lemma-compare-representable-morphism-cohomology}.
Nous voyons donc que le morphisme de comparaison
$$
c_\psi :
b_{small}^*(R^if_*\mathcal{F}|_{V'_\etale})
\longrightarrow
R^if_*\mathcal{F}|_{V_\etale}
$$
est un isomorphisme d'après
Cohomologie des espaces, lemme
\ref{spaces-cohomology-lemma-flat-base-change-cohomology}.
Ainsi, $R^if_*\mathcal{F}$ possède la propriété de changement de base plat.
Puisque $R^if_*\mathcal{F}$ est localement quasi-cohérent d'après le
lemme \ref{lemma-pushforward-locally-quasi-coherent}
nous avons terminé.
\end{proof}





\section[Modules quasi-cohérents et changement de base plat]{Modules localement quasi-cohérents possédant la propriété de changement de base plat}
\label{section-loc-qcoh-flat-base-change}

\noindent
Soit $\mathcal{X}$ un champ algébrique. Nous noterons\footnote{Nous nous excusons
de cette notation affreuse.}
$$
\textit{LQCoh}^{fbc}(\mathcal{O}_\mathcal{X})
\subset
\textit{Mod}(\mathcal{X}_\etale, \mathcal{O}_\mathcal{X})
$$
la sous-catégorie pleine dont les objets sont les
$\mathcal{O}_\mathcal{X}$-modules étales $\mathcal{F}$
qui sont à la fois localement quasi-cohérents
(section \ref{section-locally-quasi-coherent})
et possèdent la propriété de changement de base plat
(section \ref{section-flat-comparison}).
Nous avons
$$
\QCoh(\mathcal{O}_\mathcal{X}) \subset
\textit{LQCoh}^{fbc}(\mathcal{O}_\mathcal{X})
$$
d'après Faisceaux sur les champs, lemme \ref{stacks-sheaves-lemma-quasi-coherent}.

\begin{proposition}
\label{proposition-loc-qcoh-flat-base-change}
Résumé des résultats sur les modules localement quasi-cohérents possédant la propriété de
changement de base plat.
\begin{enumerate}
\item Soit $\mathcal{X}$ un champ algébrique.
Si $\mathcal{F}$ appartient à $\textit{LQCoh}^{fbc}(\mathcal{O}_\mathcal{X})$,
alors $\mathcal{F}$ est un faisceau pour la topologie fppf, c'est-à-dire
un objet de $\textit{Mod}(\mathcal{O}_\mathcal{X})$.
\item La catégorie $\textit{LQCoh}^{fbc}(\mathcal{O}_\mathcal{X})$
est une sous-catégorie faible de Serre à la fois de $\textit{Mod}(\mathcal{O}_\mathcal{X})$
et de $\textit{Mod}(\mathcal{X}_\etale, \mathcal{O}_\mathcal{X})$.
\item L'image inverse $f^*$ par tout morphisme de champs algébriques
$f : \mathcal{X} \to \mathcal{Y}$ induit un foncteur
$f^* : \textit{LQCoh}^{fbc}(\mathcal{O}_\mathcal{Y}) \to
\textit{LQCoh}^{fbc}(\mathcal{O}_\mathcal{X})$.
\item Si $f : \mathcal{X} \to \mathcal{Y}$ est un morphisme
quasi-compact et quasi-séparé de champs algébriques
et si $\mathcal{F}$ est un objet de
$\textit{LQCoh}^{fbc}(\mathcal{O}_\mathcal{X})$, alors
\begin{enumerate}
\item l'image directe totale $Rf_*\mathcal{F}$ et les images directes
supérieures $R^if_*\mathcal{F}$ peuvent être calculées indifféremment pour la topologie étale ou
la topologie fppf, avec le même résultat, et
\item chaque $R^if_*\mathcal{F}$ est un objet de
$\textit{LQCoh}^{fbc}(\mathcal{O}_\mathcal{Y})$.
\end{enumerate}
\item La catégorie $\textit{LQCoh}^{fbc}(\mathcal{O}_\mathcal{X})$ admet des
limites inductives, qui coïncident avec celles de
$\textit{Mod}(\mathcal{X}_\etale, \mathcal{O}_\mathcal{X})$
ainsi qu'avec celles de $\textit{Mod}(\mathcal{O}_\mathcal{X})$.
\item Si $\mathcal{F}$ et $\mathcal{G}$ appartiennent à
$\textit{LQCoh}^{fbc}(\mathcal{O}_\mathcal{X})$, alors le produit tensoriel
$\mathcal{F} \otimes_{\mathcal{O}_\mathcal{X}} \mathcal{G}$
appartient à $\textit{LQCoh}^{fbc}(\mathcal{O}_\mathcal{X})$.
\item Si $\mathcal{F}$ est de présentation finie et si $\mathcal{G}$ appartient à
$\textit{LQCoh}^{fbc}(\mathcal{O}_\mathcal{X})$, alors
$\SheafHom_{\mathcal{O}_\mathcal{X}}(\mathcal{F}, \mathcal{G})$
appartient à $\textit{LQCoh}^{fbc}(\mathcal{O}_\mathcal{X})$.
\end{enumerate}
\end{proposition}

\begin{proof}
La partie (1) est
Faisceaux sur les champs, lemme
\ref{stacks-sheaves-lemma-lqc-flat-base-change-fppf-sheaf}.

\medskip\noindent
La partie (2) pour l'inclusion
$\textit{LQCoh}^{fbc}(\mathcal{O}_\mathcal{X}) \subset
\textit{Mod}(\mathcal{X}_\etale, \mathcal{O}_\mathcal{X})$
a été établie dans la démonstration du
lemme \ref{lemma-flat-comparison}.
Démontrons (2) pour l'inclusion
$\textit{LQCoh}^{fbc}(\mathcal{O}_\mathcal{X}) \subset
\textit{Mod}(\mathcal{O}_\mathcal{X})$.
Soit $\varphi : \mathcal{F} \to \mathcal{G}$ un morphisme entre des
objets de $\textit{LQCoh}^{fbc}(\mathcal{O}_\mathcal{X})$. Puisque
$\Ker(\varphi)$
est le même qu'il soit calculé pour la topologie étale ou pour la topologie
fppf, nous voyons que $\Ker(\varphi)$ appartient à
$\textit{LQCoh}^{fbc}(\mathcal{O}_\mathcal{X})$
d'après le cas étale. D'autre part,
le conoyau calculé pour la topologie fppf est le faisceau fppf associé
au conoyau calculé pour la topologie étale. Or ce
conoyau étale appartient à $\textit{LQCoh}^{fbc}(\mathcal{O}_\mathcal{X})$,
donc il est un faisceau fppf
d'après (1), et nous voyons que le conoyau appartient à
$\textit{LQCoh}^{fbc}(\mathcal{O}_\mathcal{X})$.
Enfin, supposons que
$$
0 \to \mathcal{F}_1 \to \mathcal{F}_2 \to \mathcal{F}_3 \to 0
$$
soit une suite exacte dans $\textit{Mod}(\mathcal{O}_\mathcal{X})$
(c'est-à-dire pour la topologie fppf), avec $\mathcal{F}_1$, $\mathcal{F}_2$
dans $\textit{LQCoh}^{fbc}(\mathcal{O}_\mathcal{X})$. Pour montrer que
$\mathcal{F}_2$ est un objet de $\textit{LQCoh}^{fbc}(\mathcal{O}_\mathcal{X})$,
il suffit de montrer que
la suite est aussi exacte pour la topologie étale. Pour cela, il
suffit de montrer que tout élément de $H^1_{fppf}(x, \mathcal{F}_1)$
devient nul sur les membres d'un recouvrement étale de $x$ (pour tout
objet $x$ de $\mathcal{X}$). Cela est vrai parce que
$H^1_{fppf}(x, \mathcal{F}_1) = H^1_\etale(x, \mathcal{F}_1)$ d'après
Faisceaux sur les champs, lemme \ref{stacks-sheaves-lemma-compare-fppf-etale}
et par la localité de la cohomologie ; voir
Cohomologie sur les sites, lemme
\ref{sites-cohomology-lemma-kill-cohomology-class-on-covering}.
Ceci démontre (2).

\medskip\noindent
La partie (3) résulte du
lemme \ref{lemma-check-flat-comparison-on-etale-covering}
et de
Faisceaux sur les champs, lemme \ref{stacks-sheaves-lemma-pullback-lqc}.

\medskip\noindent
La partie (4)(b), pour $R^if_*\mathcal{F}$ calculé en cohomologie étale,
résulte du lemme \ref{lemma-flat-comparison}.
La partie (4)(a) résulte alors de
Faisceaux sur les champs, lemme \ref{stacks-sheaves-lemma-compare-fppf-etale}
combiné avec (1) ci-dessus.

\medskip\noindent
La partie (5), pour la topologie étale, résulte de
Faisceaux sur les champs, lemme \ref{stacks-sheaves-lemma-lqc-colimits} et du
lemme \ref{lemma-check-flat-comparison-on-etale-covering}.
La version fppf en résulte, puisque la limite inductive pour la topologie
étale est déjà un faisceau fppf d'après la partie (1).

\medskip\noindent
Les parties (6) et (7) résultent des parties correspondantes du
lemme \ref{lemma-check-flat-comparison-on-etale-covering} et de
Faisceaux sur les champs, lemme \ref{stacks-sheaves-lemma-lqc-colimits}.
\end{proof}

\begin{lemma}
\label{lemma-check-lqc-fbc-on-covering}
Soit $\mathcal{X}$ un champ algébrique.
\begin{enumerate}
\item Soit $f_j : \mathcal{X}_j \to \mathcal{X}$ une famille de morphismes
lisses de champs algébriques telle que
$|\mathcal{X}| =\bigcup |f_j|(|\mathcal{X}_j|)$.
Soit $\mathcal{F}$ un faisceau de $\mathcal{O}_\mathcal{X}$-modules
sur $\mathcal{X}_\etale$. Si chaque $f_j^{-1}\mathcal{F}$
appartient à $\textit{LQCoh}^{fpc}(\mathcal{O}_{\mathcal{X}_i})$, alors
$\mathcal{F}$ appartient à
$\textit{LQCoh}^{fbc}(\mathcal{O}_\mathcal{X})$.
\item Soit $f_j : \mathcal{X}_j \to \mathcal{X}$ une famille de morphismes plats
et localement de présentation finie de champs algébriques telle que
$|\mathcal{X}| =\bigcup |f_j|(|\mathcal{X}_j|)$.
Soit $\mathcal{F}$ un faisceau de $\mathcal{O}_\mathcal{X}$-modules
sur $\mathcal{X}_{fppf}$. Si chaque $f_j^{-1}\mathcal{F}$
appartient à $\textit{LQCoh}^{fbc}(\mathcal{O}_{\mathcal{X}_i})$, alors
$\mathcal{F}$ appartient à $\textit{LQCoh}^{fbc}(\mathcal{O}_\mathcal{X})$.
\end{enumerate}
\end{lemma}

\begin{proof}
La partie (1) résulte de la combinaison des
lemmes \ref{lemma-check-lqc-on-etale-covering} et
\ref{lemma-check-flat-comparison-on-etale-covering}.
La démonstration de (2) est analogue à celle du
lemme \ref{lemma-check-lqc-on-flat-covering}.
Soit $\mathcal{F}$ un faisceau de $\mathcal{O}_\mathcal{X}$-modules
sur $\mathcal{X}_{fppf}$.

\medskip\noindent
Supposons d'abord qu'il existe un morphisme $a : \mathcal{U} \to \mathcal{X}$
surjectif, plat, localement de présentation finie, quasi-compact
et quasi-séparé, tel que $a^*\mathcal{F}$ soit localement quasi-cohérent
et possède la propriété de changement de base plat.
Il existe alors une suite exacte
$$
0 \to \mathcal{F} \to a_*a^*\mathcal{F} \to b_*b^*\mathcal{F}
$$
où $b$ est le morphisme
$b : \mathcal{U} \times_\mathcal{X} \mathcal{U} \to \mathcal{X}$ ; voir
Faisceaux sur les champs, proposition
\ref{stacks-sheaves-proposition-exactness-cech-complex} et
lemme \ref{stacks-sheaves-lemma-surjective-flat-locally-finite-presentation}.
De plus, l'image inverse $b^*\mathcal{F}$ est l'image inverse de $a^*\mathcal{F}$
par l'une des projections ; elle est donc localement quasi-cohérente
et possède la propriété de changement de base plat ; voir
proposition \ref{proposition-loc-qcoh-flat-base-change}.
Les modules $a_*a^*\mathcal{F}$ et $b_*b^*\mathcal{F}$ sont localement
quasi-cohérents et possèdent la propriété de changement de base plat d'après la
proposition \ref{proposition-loc-qcoh-flat-base-change}.
Nous en concluons que $\mathcal{F}$ est localement quasi-cohérent et
possède la propriété de changement de base plat d'après la
proposition \ref{proposition-loc-qcoh-flat-base-change}.

\medskip\noindent
Choisissons un schéma $U$ et un morphisme lisse surjectif $x : U \to \mathcal{X}$.
D'après la partie (1), il suffit de montrer que $x^*\mathcal{F}$ est localement
quasi-cohérent et possède la propriété de changement de base plat.
Toujours d'après la partie (1), il suffit de le faire localement sur $U$ pour la topologie de Zariski ;
nous pouvons donc supposer que $U$ est affine. D'après
Morphismes de champs, lemme
\ref{stacks-morphisms-lemma-surjective-family-flat-locally-finite-presentation}
il existe un recouvrement fppf $\{a_i : U_i \to U\}$ tel que
chaque $x \circ a_i$ se factorise par un certain $f_j$. Ainsi, le module
$a_i^*\mathcal{F}$ sur $(\Sch/U_i)_{fppf}$
est localement quasi-cohérent et possède la propriété de changement de base plat.
Après avoir raffiné le recouvrement, nous pouvons supposer que $\{U_i \to U\}_{i = 1, \ldots, n}$
est un recouvrement fppf standard. Alors $x^*\mathcal{F}$ est un
module fppf sur $(\Sch/U)_{fppf}$ dont l'image inverse par le morphisme
$a : U_1 \amalg \ldots \amalg U_n \to U$ est localement quasi-cohérente
et possède la propriété de changement de base plat.
Le paragraphe précédent montre donc que $x^*\mathcal{F}$ est localement
quasi-cohérent et possède la propriété de changement de base plat, comme souhaité.
\end{proof}

\begin{lemma}
\label{lemma-loc-qcoh-fbc-compute-higher-direct-image}
Soit $f : \mathcal{X} \to \mathcal{Y}$ un morphisme de champs algébriques
qui soit quasi-compact, quasi-séparé et
représentable par des espaces algébriques. Soit $\mathcal{F}$ dans
$\textit{LQCoh}^{fbc}(\mathcal{O}_\mathcal{X})$.
Alors, pour tout objet $y : V \to \mathcal{Y}$ de $\mathcal{Y}$, nous avons
$$
(R^if_*\mathcal{F})|_{V_\etale} = R^if'_{small, *}(\mathcal{F}|_{U_\etale})
$$
où $f' : U = V \times_\mathcal{Y} \mathcal{X} \to V$ est le changement
de base de $f$.
\end{lemma}

\begin{proof}
D'après Faisceaux sur les champs, lemme
\ref{stacks-sheaves-lemma-base-change-higher-direct-images}
nous pouvons nous ramener au cas où $\mathcal{X}$
est représenté par $U$ et $\mathcal{Y}$
par $V$. Cela utilise bien entendu aussi que l'image inverse de
$\mathcal{F}$ sur $U$ appartient à $\textit{LQCoh}^{fbc}(\mathcal{O}_U)$
d'après la proposition \ref{proposition-loc-qcoh-flat-base-change}.
Le résultat découle alors de
Faisceaux sur les champs, lemme
\ref{stacks-sheaves-lemma-compare-morphism-cohomology}
et du fait que $R^if_*$ peut être calculé pour la
topologie étale d'après la
proposition \ref{proposition-loc-qcoh-flat-base-change}.
\end{proof}

\begin{lemma}
\label{lemma-loc-qcoh-fbc-affine-direct-image}
Soit $f : \mathcal{X} \to \mathcal{Y}$ un morphisme affine de champs
algébriques. Le foncteur $f_* : \textit{LQCoh}^{fbc}(\mathcal{O}_\mathcal{X}) \to
\textit{LQCoh}^{fbc}(\mathcal{O}_\mathcal{Y})$ est exact et commute
aux sommes directes. Les foncteurs $R^if_*$ pour $i > 0$ s'annulent sur
$\textit{LQCoh}^{fbc}(\mathcal{O}_\mathcal{X})$.
\end{lemma}

\begin{proof}
Les foncteurs existent d'après la proposition \ref{proposition-loc-qcoh-flat-base-change}.
D'après le lemme \ref{lemma-loc-qcoh-fbc-compute-higher-direct-image},
on se ramène au cas d'un morphisme affine d'espaces algébriques,
les images directes supérieures étant prises dans le cadre des modules
quasi-cohérents sur les espaces algébriques. D'après la discussion de
Cohomologie des espaces, section
\ref{spaces-cohomology-section-higher-direct-image}
nous nous ramenons au cas d'un morphisme affine de schémas.
Pour les morphismes affines de schémas, l'annulation des
images directes supérieures sur les modules quasi-cohérents résulte de
Cohomologie des schémas, lemme \ref{coherent-lemma-relative-affine-vanishing}.
L'annulation de $R^1f_*$ implique l'exactitude de $f_*$.
La commutation aux sommes directes résulte de
Morphismes, lemme \ref{morphisms-lemma-affine-equivalence-modules}
par exemple.
\end{proof}










\section{Modules parasites}
\label{section-parasitic}

\noindent
La définition suivante est compatible avec
Descente, définition \ref{descent-definition-parasitic}.

\begin{definition}
\label{definition-parasitic}
Soit $\mathcal{X}$ un champ algébrique.
Un préfaisceau de $\mathcal{O}_\mathcal{X}$-modules $\mathcal{F}$ est
{\it parasite} si $\mathcal{F}(x) = 0$ pour tout objet $x$
de $\mathcal{X}$ au-dessus d'un schéma $U$ tel que le morphisme
correspondant $x : U \to \mathcal{X}$ soit plat.
\end{definition}

\noindent
Voici un lemme donnant quelques propriétés de cette notion.

\begin{lemma}
\label{lemma-parasitic}
Soit $\mathcal{X}$ un champ algébrique. Soit $\mathcal{F}$
un préfaisceau de $\mathcal{O}_\mathcal{X}$-modules.
\begin{enumerate}
\item Si $\mathcal{F}$ est parasite et si
$g : \mathcal{Y} \to \mathcal{X}$ est un morphisme plat de champs algébriques,
alors $g^*\mathcal{F}$ est parasite.
\item Pour $\tau \in \{Zariski, \etale, smooth, syntomic, fppf\}$,
nous avons :
\begin{enumerate}
\item le faisceau $\tau$ associé à un préfaisceau parasite de modules est
parasite, et
\item la sous-catégorie pleine de
$\textit{Mod}(\mathcal{X}_\tau, \mathcal{O}_\mathcal{X})$
formée des modules parasites est une sous-catégorie de Serre.
\end{enumerate}
\item Supposons que $\mathcal{F}$ soit un faisceau pour la topologie étale.
Soit $f_i : \mathcal{X}_i \to \mathcal{X}$ une famille de
morphismes lisses de champs algébriques telle que
$|\mathcal{X}| = \bigcup_i |f_i|(|\mathcal{X}_i|)$. Si chaque
$f_i^*\mathcal{F}$ est parasite, alors $\mathcal{F}$ l'est aussi.
\item Supposons que $\mathcal{F}$ soit un faisceau pour la topologie fppf.
Soit $f_i : \mathcal{X}_i \to \mathcal{X}$ une famille de
morphismes plats et localement de présentation finie de champs algébriques telle que
$|\mathcal{X}| = \bigcup_i |f_i|(|\mathcal{X}_i|)$. Si chaque
$f_i^*\mathcal{F}$ est parasite, alors $\mathcal{F}$ l'est aussi.
\end{enumerate}
\end{lemma}

\begin{proof}
Pour démontrer la partie (1), soit $y$ un objet de $\mathcal{Y}$ au-dessus
d'un schéma $V$ tel que le morphisme correspondant $y : V \to \mathcal{Y}$
soit plat. Alors $g(y) : V \to \mathcal{Y} \to \mathcal{X}$ est plat
comme composé de morphismes plats (voir
Morphismes de champs, lemme \ref{stacks-morphisms-lemma-composition-flat})
et donc $\mathcal{F}(g(y))$ est nul par hypothèse. Puisque
$g^*\mathcal{F} = g^{-1}\mathcal{F}(y) = \mathcal{F}(g(y))$, nous en concluons que
$g^*\mathcal{F}$ est parasite.

\medskip\noindent
Pour démontrer la partie (2)(a), notons que si $\{x_i \to x\}$ est un recouvrement pour la topologie $\tau$
de $\mathcal{X}$, alors chacun des morphismes $x_i \to x$ est au-dessus
d'un morphisme plat de schémas. Ainsi, si $x$ est au-dessus d'un schéma
$U$ tel que $x : U \to \mathcal{X}$ soit plat, il en va de même de tous les
objets $x_i$. Par conséquent, le préfaisceau $\mathcal{F}^+$ (voir
Sites, section \ref{sites-section-sheafification})
est parasite si le préfaisceau $\mathcal{F}$ est
parasite. Cela démontre (2)(a), puisque le faisceau associé à $\mathcal{F}$
est $(\mathcal{F}^+)^+$.

\medskip\noindent
Soit $\mathcal{F}$ un $\tau$-module parasite. Il résulte immédiatement des
définitions que tout sous-module de $\mathcal{F}$ est parasite. D'autre
part, si $\mathcal{F}' \subset \mathcal{F}$ est un sous-module, il est
également clair que le préfaisceau
$x \mapsto \mathcal{F}(x)/\mathcal{F}'(x)$
est parasite. Ainsi, le quotient $\mathcal{F}/\mathcal{F}'$ est un module
parasite d'après (2)(a). Enfin, il faut montrer que, pour toute suite exacte courte
$0 \to \mathcal{F}_1 \to \mathcal{F}_2 \to \mathcal{F}_3 \to 0$
dont $\mathcal{F}_1$ et $\mathcal{F}_3$ sont parasites, $\mathcal{F}_2$
est parasite. Cela résulte immédiatement de l'évaluation en $x$, lorsque cet objet est
au-dessus d'un schéma plat sur $\mathcal{X}$. Ceci démontre (2)(b) ; voir
Homologie, lemme \ref{homology-lemma-characterize-serre-subcategory}.

\medskip\noindent
Soit $f_i : \mathcal{X}_i \to \mathcal{X}$ une famille conjointement surjective de
morphismes lisses de champs algébriques, et supposons que chaque $f_i^*\mathcal{F}$
soit parasite. Soit $x$ un objet de $\mathcal{X}$ au-dessus d'un
schéma $U$ tel que $x : U \to \mathcal{X}$ soit plat. Considérons un recouvrement
lisse surjectif $W_i \to U \times_{x, \mathcal{X}} \mathcal{X}_i$.
Notons $y_i : W_i \to \mathcal{X}_i$ la projection. Il s'ensuit
que $\{f_i(y_i) \to x\}$ est un recouvrement pour la topologie lisse
sur $\mathcal{X}$. Comme un composé de morphismes plats est plat, nous voyons que
$f_i^*\mathcal{F}(y_i) = 0$. D'autre part, comme nous l'avons vu dans la démonstration
de (1), nous avons $f_i^*\mathcal{F}(y_i) = \mathcal{F}(f_i(y_i))$.
Ainsi, pour un certain recouvrement lisse $\{x_i \to x\}_{i \in I}$
de $\mathcal{X}$, nous avons $\mathcal{F}(x_i) = 0$. Cela implique
$\mathcal{F}(x) = 0$, car la topologie lisse est la même
que la topologie étale ; voir
Compléments sur les morphismes, lemme \ref{more-morphisms-lemma-etale-dominates-smooth}.
En effet, $\{x_i \to x\}_{i \in I}$ est au-dessus d'un recouvrement lisse
$\{U_i \to U\}_{i \in I}$ de schémas. D'après le lemme qui vient d'être cité,
il existe un recouvrement étale $\{V_j \to U\}_{j \in J}$ qui
raffine $\{U_i \to U\}_{i \in I}$. Notons $x'_j = x|_{V_j}$.
Alors $\{x'_j \to x\}$ est un recouvrement étale dans $\mathcal{X}$
raffinant $\{x_i \to x\}_{i \in I}$. Cela signifie que l'application
$\mathcal{F}(x) \to \prod_{j \in J} \mathcal{F}(x'_j)$, qui est
injective puisque $\mathcal{F}$ est un faisceau pour la topologie étale,
se factorise par $\mathcal{F}(x) \to \prod_{i \in I} \mathcal{F}(x_i)$,
qui est nulle. Ainsi, $\mathcal{F}(x) = 0$, comme souhaité.

\medskip\noindent
Démonstration de (4) : omise. Indication : elle est semblable, mais plus simple, à celle de (3).
\end{proof}

\noindent
Les modules parasites sont préservés par toute image directe, sans aucune restriction.

\begin{lemma}
\label{lemma-pushforward-parasitic}
Soit $\tau \in \{\etale, fppf\}$.
Soit $\mathcal{X}$ un champ algébrique.
Soit $\mathcal{F}$ un objet parasite de
$\textit{Mod}(\mathcal{X}_\tau, \mathcal{O}_\mathcal{X})$.
\begin{enumerate}
\item $H^i_\tau(\mathcal{X}, \mathcal{F}) = 0$ pour tout $i$.
\item Soit $f : \mathcal{X} \to \mathcal{Y}$ un morphisme de champs algébriques.
Alors $R^if_*\mathcal{F}$ (calculé pour la topologie $\tau$) est un
objet parasite de $\textit{Mod}(\mathcal{Y}_\tau, \mathcal{O}_\mathcal{Y})$.
\end{enumerate}
\end{lemma}

\begin{proof}
Nous ramenons d'abord (2) à (1).
D'après Faisceaux sur les champs, lemme \ref{stacks-sheaves-lemma-pushforward-restriction},
nous voyons que $R^if_*\mathcal{F}$ est le faisceau associé au préfaisceau
$$
y \longmapsto
H^i_\tau\Big(V \times_{y, \mathcal{Y}} \mathcal{X},
\ \text{pr}^{-1}\mathcal{F}\Big)
$$
Ici, $y$ est un objet type de $\mathcal{Y}$ au-dessus du schéma $V$.
D'après le lemme \ref{lemma-parasitic}, il suffit de montrer que
ces groupes de cohomologie sont nuls lorsque $y : V \to \mathcal{Y}$ est plat.
Notons que $\text{pr} : V \times_{y, \mathcal{Y}} \mathcal{X} \to \mathcal{X}$
est plat comme changement de base de $y$. Ainsi, d'après le
lemme \ref{lemma-parasitic}, nous voyons que $\text{pr}^{-1}\mathcal{F}$
est parasite. Il suffit donc de démontrer (1).

\medskip\noindent
Pour démontrer (1), nous pouvons utiliser la suite spectrale de
Faisceaux sur les champs, proposition
\ref{stacks-sheaves-proposition-smooth-covering-compute-cohomology}
pour nous ramener au cas où $\mathcal{X}$
est un champ algébrique représentable par un espace algébrique.
Notons que, dans la suite spectrale, chaque
$f_p^{-1}\mathcal{F} = f_p^*\mathcal{F}$ est un module parasite d'après le
lemme \ref{lemma-parasitic}, car les morphismes
$f_p : \mathcal{U}_p =
\mathcal{U} \times_\mathcal{X} \ldots
\times_\mathcal{X} \mathcal{U} \to \mathcal{X}$ sont plats.
En réutilisant encore une fois cette suite spectrale (comme dans la
démonstration du lemme \ref{lemma-general-pushforward}),
nous nous ramenons au cas où le
champ algébrique $\mathcal{X}$ est représentable par un schéma $X$.
Alors $H^i_\tau(\mathcal{X}, \mathcal{F}) = H^i((\Sch/X)_\tau, \mathcal{F})$.
Dans ce cas, l'annulation résulte facilement d'un argument
avec des recouvrements de {\v C}ech ; voir
Descente, lemme \ref{descent-lemma-cohomology-parasitic}.
\end{proof}

\noindent
Le lemme suivant donne l'une des principales raisons de nous intéresser aux
modules parasites. Pour en comprendre l'énoncé, rappelons que les
foncteurs
$\QCoh(\mathcal{O}_\mathcal{X}) \to
\textit{Mod}(\mathcal{X}_\etale, \mathcal{O}_\mathcal{X})$
et
$\QCoh(\mathcal{O}_\mathcal{X}) \to
\textit{Mod}(\mathcal{O}_\mathcal{X})$
ne sont pas exacts en général.

\begin{lemma}
\label{lemma-exact-sequence-quasi-coherent-parasitic-cohomology}
Soit $\mathcal{X}$ un champ algébrique. Soient
$\alpha : \mathcal{F} \to \mathcal{G}$ et
$\beta : \mathcal{G} \to \mathcal{H}$
des morphismes dans $\QCoh(\mathcal{O}_\mathcal{X})$ tels que
$\beta \circ \alpha = 0$. Les assertions suivantes sont équivalentes :
\begin{enumerate}
\item dans la catégorie abélienne $\QCoh(\mathcal{O}_\mathcal{X})$,
le complexe $\mathcal{F} \to \mathcal{G} \to \mathcal{H}$
est exact en $\mathcal{G}$ ;
\item $\Ker(\beta)/\Im(\alpha)$, calculé soit dans
$\textit{Mod}(\mathcal{X}_\etale, \mathcal{O}_\mathcal{X})$, soit dans
$\textit{Mod}(\mathcal{X}_{fppf}, \mathcal{O}_\mathcal{X})$
est parasite.
\end{enumerate}
\end{lemma}

\begin{proof}
Nous avons $\QCoh(\mathcal{O}_\mathcal{X}) \subset
\textit{LQCoh}^{fbc}(\mathcal{O}_\mathcal{X})$ ; voir
section \ref{section-loc-qcoh-flat-base-change}.
Ainsi, les objets $\Ker(\beta)/\Im(\alpha)$ calculés dans
$\textit{Mod}(\mathcal{X}_\etale, \mathcal{O}_\mathcal{X})$ ou dans
$\textit{Mod}(\mathcal{X}_{fppf}, \mathcal{O}_\mathcal{X})$ coïncident ; voir
proposition \ref{proposition-loc-qcoh-flat-base-change}.
Désormais, nous utiliserons la topologie étale sur $\mathcal{X}$.

\medskip\noindent
Soit $\mathcal{E}$ la cohomologie de
$\mathcal{F} \to \mathcal{G} \to \mathcal{H}$
calculée dans la catégorie abélienne $\QCoh(\mathcal{O}_\mathcal{X})$.
Soit $x : U \to \mathcal{X}$ un morphisme plat, où $U$ est un schéma.
Comme nous utilisons la topologie étale, le foncteur de restriction
$\textit{Mod}(\mathcal{X}_\etale, \mathcal{O}_\mathcal{X})
\to \textit{Mod}(U_\etale, \mathcal{O}_U)$ est exact.
D'autre part, d'après le lemme \ref{lemma-flat-pullback-quasi-coherent} et
Faisceaux sur les champs, lemme \ref{stacks-sheaves-lemma-compare-quasi-coherent}
le foncteur de restriction
$$
\QCoh(\mathcal{O}_\mathcal{X}) \xrightarrow{x^*}
\QCoh((\Sch/U)_\etale, \mathcal{O}) \xrightarrow{{-}|_{U_\etale}}
\QCoh(U_\etale, \mathcal{O}_U)
$$
est lui aussi exact. Nous en concluons que
$\mathcal{E}|_{U_\etale} = (\Ker(\beta)/\Im(\alpha))|_{U_\etale}$.

\medskip\noindent
Si (1) vaut, alors $\mathcal{E} = 0$ ; ainsi, $\Ker(\beta)/\Im(\alpha)$
se restreint à zéro sur $U_\etale$ pour tout $U$ plat sur $\mathcal{X}$, et c'est
la définition d'un module parasite. Si (2) vaut, alors
$\Ker(\beta)/\Im(\alpha)$ se restreint à zéro sur $U_\etale$ pour tout $U$
plat sur $\mathcal{X}$ ; ainsi, $\mathcal{E}$ se restreint à zéro sur
$U_\etale$ pour tout $U$ plat sur $\mathcal{X}$. Cela implique certainement
que le module quasi-cohérent $\mathcal{E}$ est nul ; on peut par exemple appliquer le
lemme \ref{lemma-quasi-coherent-check-exact}
au morphisme $0 \to \mathcal{E}$.
\end{proof}





\section{Modules quasi-cohérents}
\label{section-quasi-coherent}

\noindent
Nous avons vu que la catégorie des modules quasi-cohérents sur un champ
algébrique est équivalente à la catégorie des modules quasi-cohérents sur une
présentation ; voir
Faisceaux sur les champs, section
\ref{stacks-sheaves-section-quasi-coherent-algebraic-stacks}.
Ce fait est à la base de ce qui suit.

\begin{lemma}
\label{lemma-adjoint}
Soit $\mathcal{X}$ un champ algébrique. Soit
$\textit{LQCoh}^{fbc}(\mathcal{O}_\mathcal{X})$
la catégorie des modules localement quasi-cohérents possédant la
propriété de changement de base plat ; voir
section \ref{section-loc-qcoh-flat-base-change}.
Le foncteur d'inclusion
$i : \QCoh(\mathcal{O}_\mathcal{X}) \to
\textit{LQCoh}^{fbc}(\mathcal{O}_\mathcal{X})$
admet un adjoint à droite
$$
Q : \textit{LQCoh}^{fbc}(\mathcal{O}_\mathcal{X}) \to
\QCoh(\mathcal{O}_\mathcal{X})
$$
tel que $Q \circ i$ soit le foncteur identité.
\end{lemma}

\begin{proof}
Choisissons un schéma $U$ et un morphisme lisse surjectif $f : U \to \mathcal{X}$.
Posons $R = U \times_\mathcal{X} U$ ; nous obtenons ainsi un groupoïde lisse
$(U, R, s, t, c)$ en espaces algébriques tel que
$\mathcal{X} = [U/R]$ ; voir
Champs algébriques, lemme \ref{algebraic-lemma-stack-presentation}.
Nous pouvons remplacer, et remplaçons, $\mathcal{X}$ par $[U/R]$. D'après
Faisceaux sur les champs, proposition \ref{stacks-sheaves-proposition-quasi-coherent}
il existe une équivalence
$$
q_1 :
\QCoh(U, R, s, t, c)
\longrightarrow
\QCoh(\mathcal{O}_\mathcal{X})
$$
Construisons un foncteur
$$
q_2 :
\textit{LQCoh}^{fbc}(\mathcal{O}_\mathcal{X})
\longrightarrow
\QCoh(U, R, s, t, c)
$$
par la règle suivante : si $\mathcal{F}$ est un objet de
$\textit{LQCoh}^{fbc}(\mathcal{O}_\mathcal{X})$, nous posons
$$
q_2(\mathcal{F}) = (f^*\mathcal{F}|_{U_\etale}, \alpha)
$$
où $\alpha$ est l'isomorphisme
$$
t_{small}^*(f^*\mathcal{F}|_{U_\etale})
\to
t^*f^*\mathcal{F}|_{R_\etale} \to
s^*f^*\mathcal{F}|_{R_\etale} \to
s_{small}^*(f^*\mathcal{F}|_{U_\etale})
$$
où les deux morphismes extrêmes sont les morphismes de comparaison. Notons que
$q_2(\mathcal{F})$ est quasi-cohérent précisément parce que $\mathcal{F}$ est
localement quasi-cohérent, et que nous avons utilisé (et eu besoin de)
la propriété de changement de base plat dans la construction de
la donnée de descente $\alpha$. Nous omettons la
vérification de la condition de cocycle (voir
Groupoïdes dans les espaces, définition
\ref{spaces-groupoids-definition-groupoid-module}).
En examinant la démonstration de
Faisceaux sur les champs, proposition \ref{stacks-sheaves-proposition-quasi-coherent}
nous voyons que $q_2 \circ i$ est un quasi-inverse de $q_1$.
Définissons $Q = q_1 \circ q_2$.
Soit $\mathcal{F}$ un objet de
$\textit{LQCoh}^{fbc}(\mathcal{O}_\mathcal{X})$ et
soit $\mathcal{G}$ un objet de $\QCoh(\mathcal{O}_\mathcal{X})$.
Nous avons
\begin{align*}
\Mor_{\textit{LQCoh}^{fbc}(\mathcal{O}_\mathcal{X})}
(i(\mathcal{G}), \mathcal{F})
& =
\Mor_{\QCoh(U, R, s, t, c)}(q_2(i(\mathcal{G})), q_2(\mathcal{F})) \\
& =
\Mor_{\QCoh(\mathcal{O}_\mathcal{X})}(\mathcal{G}, Q(\mathcal{F}))
\end{align*}
où la première égalité résulte de
Faisceaux sur les champs, lemme \ref{stacks-sheaves-lemma-map-from-quasi-coherent}
et la deuxième vaut parce que $q_1 \circ i$ et $q_2$ sont des équivalences de catégories
quasi-inverses. L'assertion $Q \circ i \cong \text{id}$
est une conséquence formelle du fait que $i$ est pleinement fidèle.
\end{proof}

\begin{lemma}
\label{lemma-adjoint-kernel-parasitic}
Soit $\mathcal{X}$ un champ algébrique.
Soit $Q : \textit{LQCoh}^{fbc}(\mathcal{O}_\mathcal{X}) \to
\QCoh(\mathcal{O}_\mathcal{X})$
le foncteur construit dans le lemme \ref{lemma-adjoint}.
\begin{enumerate}
\item Le noyau de $Q$ est exactement la collection des objets parasites
de $\textit{LQCoh}^{fbc}(\mathcal{O}_\mathcal{X})$.
\item Pour tout objet $\mathcal{F}$
de $\textit{LQCoh}^{fbc}(\mathcal{O}_\mathcal{X})$, le noyau et le
conoyau du
morphisme d'adjonction $Q(\mathcal{F}) \to \mathcal{F}$ sont parasites.
\item Le foncteur $Q$ est exact et commute à toutes les limites projectives et inductives.
\end{enumerate}
\end{lemma}

\begin{proof}
Écrivons $\mathcal{X} = [U/R]$ comme dans la démonstration du lemme \ref{lemma-adjoint}.
Soit $\mathcal{F}$ un objet de
$\textit{LQCoh}^{fbc}(\mathcal{O}_\mathcal{X})$.
Il ressort de la démonstration du lemme \ref{lemma-adjoint}
que $\mathcal{F}$ appartient au noyau de $Q$ si et seulement si
$\mathcal{F}|_{U_\etale} = 0$.
En particulier, si $\mathcal{F}$ est parasite, alors $\mathcal{F}$ appartient au
noyau. Soit ensuite $x : V \to \mathcal{X}$ un morphisme plat, où
$V$ est un schéma. Posons $W = V \times_\mathcal{X} U$ et considérons le diagramme
$$
\xymatrix{
W \ar[d]_p \ar[r]_q & V \ar[d] \\
U \ar[r] & \mathcal{X}
}
$$
Notons que la projection $p : W \to U$ est plate et que la projection
$q : W \to V$ est lisse et surjective. Il s'ensuit que $q_{small}^*$
est un foncteur fidèle sur les modules quasi-cohérents. Par hypothèse, $\mathcal{F}$
possède la propriété de changement de base plat, de sorte que nous obtenons
$p_{small}^*\mathcal{F}|_{U_\etale} \cong
q_{small}^*\mathcal{F}|_{V_\etale}$. Ainsi, si $\mathcal{F}$
appartient au noyau de $Q$, alors $\mathcal{F}|_{V_\etale} = 0$,
ce qui achève la démonstration de (1).

\medskip\noindent
La partie (2) résulte de la discussion précédente et du fait
que le morphisme $Q(\mathcal{F}) \to \mathcal{F}$ devient un isomorphisme après
restriction à $U_\etale$.

\medskip\noindent
Pour démontrer la partie (3), notons que $Q$ est exact à gauche comme adjoint à droite.
Soit $0 \to \mathcal{F} \to \mathcal{G} \to \mathcal{H} \to 0$
une suite exacte courte dans $\textit{LQCoh}^{fbc}(\mathcal{O}_\mathcal{X})$.
Considérons le diagramme commutatif suivant
$$
\xymatrix{
0 \ar[r] &
Q(\mathcal{F}) \ar[r] \ar[d]_a &
Q(\mathcal{G}) \ar[r] \ar[d]_b &
Q(\mathcal{H}) \ar[r] \ar[d]_c & 0 \\
0 \ar[r] &
\mathcal{F} \ar[r] &
\mathcal{G} \ar[r] &
\mathcal{H} \ar[r] & 0
}
$$
Puisque les noyaux et conoyaux de $a$, $b$ et $c$ sont parasites d'après
la partie (2), et puisque la ligne inférieure est une suite exacte courte, nous voyons que
la ligne supérieure, considérée comme complexe de $\mathcal{O}_\mathcal{X}$-modules, a des faisceaux de
cohomologie parasites (nous omettons les détails ; on utilise ici que la catégorie des
modules parasites est une sous-catégorie de Serre de la catégorie de tous les modules).
Par l'exactitude à gauche de $Q$, seule l'exactitude en $Q(\mathcal{H})$
reste à établir. Or le conoyau $\mathcal{Q}$ de
$Q(\mathcal{G}) \to Q(\mathcal{H}))$
peut être calculé soit dans $\textit{Mod}(\mathcal{O}_\mathcal{X})$,
soit dans $\QCoh(\mathcal{O}_\mathcal{X})$, avec le même résultat, car
le foncteur d'inclusion $\QCoh(\mathcal{O}_\mathcal{X}) \to
\textit{LQCoh}^{fbc}(\mathcal{O}_\mathcal{X})$
est un adjoint à gauche, donc est exact à droite. Ainsi, $\mathcal{Q} = Q(\mathcal{Q})$
est à la fois quasi-cohérent et parasite, donc $0$ d'après la partie (1), comme souhaité.

\medskip\noindent
Comme adjoint à droite, $Q$ commute à toutes les limites projectives. Puisque $Q$ est exact, pour
montrer que $Q$ commute à toutes les limites inductives, il suffit de montrer que $Q$
commute aux sommes directes ; voir
Catégories, lemme \ref{categories-lemma-colimits-coproducts-coequalizers}.
Soit $\mathcal{F}_i$, $i \in I$, une famille d'objets de
$\textit{LQCoh}^{fbc}(\mathcal{O}_\mathcal{X})$. Pour voir que
$Q(\bigoplus \mathcal{F}_i)$ est égal à $\bigoplus Q(\mathcal{F}_i)$,
considérons la construction de $Q$ dans la démonstration du
lemme \ref{lemma-adjoint}. Elle utilise une présentation $\mathcal{X} = [U/R]$,
où $U$ est un schéma. Alors $Q(\mathcal{F})$
se calcule en prenant d'abord la paire
$(\mathcal{F}|_{U_\etale}, \alpha)$ dans $\QCoh(U, R, s, t, c)$,
puis en utilisant l'équivalence
$\QCoh(U, R, s, t, c) \cong \QCoh(\mathcal{O}_\mathcal{X})$.
Puisque le foncteur de restriction $\textit{Mod}(\mathcal{O}_\mathcal{X}) \to
\textit{Mod}(\mathcal{O}_{U_\etale})$,
$\mathcal{F} \mapsto \mathcal{F}|_{U_\etale}$ commute aux
sommes directes, l'égalité voulue est claire.
\end{proof}

\begin{lemma}
\label{lemma-flat-morphism-and-Q}
Soit $f : \mathcal{X} \to \mathcal{Y}$ un morphisme plat de champs algébriques.
Alors $Q_\mathcal{X} \circ f^* = f^* \circ Q_\mathcal{Y}$, où
$Q_\mathcal{X}$ et $Q_\mathcal{Y}$ sont comme dans le lemme \ref{lemma-adjoint}.
\end{lemma}

\begin{proof}
Observons que $f^*$ préserve à la fois $\QCoh$ et $\textit{LQCoh}^{fbc}$ ;
voir Faisceaux sur les champs, lemme
\ref{stacks-sheaves-lemma-pullback-quasi-coherent} et
proposition \ref{proposition-loc-qcoh-flat-base-change}.
Si $\mathcal{F}$ appartient à $\textit{LQCoh}^{fbc}(\mathcal{O}_\mathcal{Y})$,
alors $Q_\mathcal{Y}(\mathcal{F}) \to \mathcal{F}$
a un noyau et un conoyau parasites
d'après le lemme \ref{lemma-adjoint-kernel-parasitic}.
Comme $f$ est plat, nous obtenons que $f^*Q_\mathcal{Y}(\mathcal{F}) \to f^*\mathcal{F}$
a un noyau et un conoyau parasites d'après le
lemme \ref{lemma-parasitic}.
Ainsi, le morphisme induit
$f^*Q_\mathcal{Y}(\mathcal{F}) \to Q_\mathcal{X}(f^*\mathcal{F})$
a un noyau et un conoyau parasites et est donc un isomorphisme,
par exemple d'après le
lemme \ref{lemma-exact-sequence-quasi-coherent-parasitic-cohomology}.
\end{proof}

\begin{lemma}
\label{lemma-flat-object-and-Q}
Soit $\mathcal{X}$ un champ algébrique. Soit $x$ un objet de $\mathcal{X}$
au-dessus du schéma $U$ tel que $x : U \to \mathcal{X}$ soit plat.
Alors, pour $\mathcal{F}$ dans $\QCoh^{fbc}(\mathcal{O}_\mathcal{X})$,
nous avons $Q(\mathcal{F})|_{U_\etale} = \mathcal{F}|_{U_\etale}$.
\end{lemma}

\begin{proof}
C'est vrai parce que le noyau et le conoyau de $Q(\mathcal{F}) \to \mathcal{F}$
sont parasites ; voir le lemme \ref{lemma-adjoint-kernel-parasitic}.
\end{proof}

\begin{remark}
\label{remark-QCoh-abelian}
Soit $\mathcal{X}$ un champ algébrique. La catégorie
$\QCoh(\mathcal{O}_\mathcal{X})$ est abélienne, le foncteur d'inclusion
$\QCoh(\mathcal{O}_\mathcal{X}) \to \textit{Mod}(\mathcal{O}_\mathcal{X})$
est exact à droite, mais n'est pas exact en général ; voir Faisceaux sur les champs, lemme
\ref{stacks-sheaves-lemma-quasi-coherent-algebraic-stack}.
Nous pouvons utiliser le foncteur $Q$ des
lemmes \ref{lemma-adjoint} et \ref{lemma-adjoint-kernel-parasitic}
pour comprendre ceci. Plus précisément, soit $\varphi : \mathcal{F} \to \mathcal{G}$
un morphisme de $\mathcal{O}_\mathcal{X}$-modules quasi-cohérents. Alors
\begin{enumerate}
\item le conoyau $\Coker(\varphi)$ calculé dans
$\textit{Mod}(\mathcal{O}_\mathcal{X})$ est quasi-cohérent et
est le conoyau de $\varphi$ dans $\QCoh(\mathcal{O}_\mathcal{X})$ ;
\item l'image $\Im(\varphi)$ calculée dans
$\textit{Mod}(\mathcal{O}_\mathcal{X})$ est quasi-cohérente et
est l'image de $\varphi$ dans $\QCoh(\mathcal{O}_\mathcal{X})$ ; et
\item le noyau $\Ker(\varphi)$ calculé dans
$\textit{Mod}(\mathcal{O}_\mathcal{X})$
appartient à $\textit{LQCoh}^{fbc}(\mathcal{O}_\mathcal{X})$
d'après la proposition \ref{proposition-loc-qcoh-flat-base-change}, et
$Q(\Ker(\varphi))$ est le noyau dans $\QCoh(\mathcal{O}_\mathcal{X})$.
\end{enumerate}
Cela résulte des références indiquées.
\end{remark}

\begin{remark}
\label{remark-QCoh-tensor}
Soit $\mathcal{X}$ un champ algébrique. Étant donnés deux
$\mathcal{O}_\mathcal{X}$-modules quasi-cohérents $\mathcal{F}$ et $\mathcal{G}$,
le module produit tensoriel
$\mathcal{F} \otimes_{\mathcal{O}_\mathcal{X}} \mathcal{G}$
est quasi-cohérent ; voir Faisceaux sur les champs, lemme
\ref{stacks-sheaves-lemma-quasi-coherent-algebraic-stack}, partie (5).
De même, étant donnés deux modules localement quasi-cohérents possédant
la propriété de changement de base plat, leur produit tensoriel possède la
même propriété ; voir la proposition \ref{proposition-loc-qcoh-flat-base-change}.
Ainsi, les foncteurs d'inclusion
$$
\QCoh(\mathcal{O}_\mathcal{X}) \to
\textit{LQCoh}^{fbc}(\mathcal{O}_\mathcal{X}) \to
\textit{Mod}(\mathcal{O}_\mathcal{X})
$$
sont des foncteurs de catégories monoïdales symétriques. Plus intéressant encore,
le foncteur
$$
Q :
\textit{LQCoh}^{fbc}(\mathcal{O}_\mathcal{X})
\longrightarrow
\QCoh(\mathcal{O}_\mathcal{X})
$$
est lui aussi un foncteur de catégories monoïdales symétriques. En effet, étant donnés
$\mathcal{F}$ et $\mathcal{G}$ dans
$\textit{LQCoh}^{fbc}(\mathcal{O}_\mathcal{X})$, nous obtenons
$$
\xymatrix{
Q(\mathcal{F})
\otimes_{\mathcal{O}_\mathcal{X}}
Q(\mathcal{G}) \ar[rr] \ar[rd] & &
\mathcal{F}
\otimes_{\mathcal{O}_\mathcal{X}}
\mathcal{G} \\
&
Q(\mathcal{F} \otimes_{\mathcal{O}_\mathcal{X}} \mathcal{G}) \ar[ru]
}
$$
où la flèche sud-ouest provient de la propriété universelle
de la flèche nord-ouest (et du fait déjà mentionné que
l'objet du coin supérieur gauche est quasi-cohérent).
Si nous restreignons ce diagramme à $U_\etale$ pour $U \to \mathcal{X}$
plat, alors les trois flèches deviennent des isomorphismes (voir les
lemmes \ref{lemma-adjoint} et \ref{lemma-adjoint-kernel-parasitic}
et la définition \ref{definition-parasitic}).
Ainsi, $Q(\mathcal{F}) \otimes_{\mathcal{O}_\mathcal{X}}
Q(\mathcal{G}) \to
Q(\mathcal{F} \otimes_{\mathcal{O}_\mathcal{X}} \mathcal{G})$
est un isomorphisme ; voir par exemple le
lemme \ref{lemma-quasi-coherent-check-exact}.
\end{remark}

\begin{remark}
\label{remark-bousfield-colocalization}
Soit $\mathcal{X}$ un champ algébrique. Soit
$\textit{Parasitic}(\mathcal{O}_\mathcal{X}) \subset
\textit{Mod}(\mathcal{O}_\mathcal{X})$ la
sous-catégorie pleine formée des modules parasites. Les résultats des
lemmes \ref{lemma-adjoint} et \ref{lemma-adjoint-kernel-parasitic}
impliquent que
$$
\QCoh(\mathcal{O}_\mathcal{X}) =
\textit{LQCoh}^{fbc}(\mathcal{O}_\mathcal{X}) /
\textit{Parasitic}(\mathcal{O}_\mathcal{X})
\cap \textit{LQCoh}^{fbc}(\mathcal{O}_\mathcal{X})
$$
autrement dit : la catégorie des modules quasi-cohérents est la catégorie
des modules localement quasi-cohérents possédant la propriété de changement de base plat,
quotientée par la sous-catégorie de Serre formée des objets parasites.
Voir Homologie, lemme \ref{homology-lemma-serre-subcategory-is-kernel}.
L'existence du foncteur d'inclusion
$i : \QCoh(\mathcal{O}_\mathcal{X}) \to
\textit{LQCoh}^{fbc}(\mathcal{O}_\mathcal{X})$
qui est adjoint à gauche du foncteur quotient est un trait essentiel de
la situation. Dans Catégories dérivées de champs, section
\ref{stacks-perfect-section-derived}, et en particulier dans le
lemme \ref{stacks-perfect-lemma-bousfield-colocalization}
nous démontrons qu'un résultat analogue vaut au niveau des catégories dérivées.
\end{remark}

\begin{lemma}
\label{lemma-internal-hom-fp-into-qcoh}
Soit $\mathcal{X}$ un champ algébrique. Soit $\mathcal{F}$ un
$\mathcal{O}_\mathcal{X}$-module de présentation finie et soit
$\mathcal{G}$ un $\mathcal{O}_\mathcal{X}$-module quasi-cohérent.
Les Hom internes
$\SheafHom_{\mathcal{O}_\mathcal{X}}(\mathcal{F}, \mathcal{G})$
calculés dans
$\textit{Mod}(\mathcal{X}_\etale, \mathcal{O}_\mathcal{X})$ ou dans
$\textit{Mod}(\mathcal{O}_\mathcal{X})$ coïncident, et leur valeur commune
est un objet de $\textit{LQCoh}^{fbc}(\mathcal{O}_\mathcal{X})$.
Le module quasi-cohérent
$
hom(\mathcal{F}, \mathcal{G}) =
Q(\SheafHom_{\mathcal{O}_\mathcal{X}}(\mathcal{F}, \mathcal{G}))
$
possède la propriété universelle suivante
$$
\Hom_\mathcal{X}(\mathcal{H}, hom(\mathcal{F}, \mathcal{G})) =
\Hom_\mathcal{X}(\mathcal{H} \otimes_{\mathcal{O}_\mathcal{X}} \mathcal{F},
\mathcal{G})
$$
pour $\mathcal{H}$ dans $\QCoh(\mathcal{O}_\mathcal{X})$.
\end{lemma}

\begin{proof}
La construction de
$\SheafHom_{\mathcal{O}_\mathcal{X}}(\mathcal{F}, \mathcal{G})$ dans
Modules sur les sites, section \ref{sites-modules-section-internal-hom}
ne dépend que de $\mathcal{F}$ et $\mathcal{G}$ comme préfaisceaux
de modules ; le résultat $\SheafHom$ est un faisceau pour la topologie
fppf parce que $\mathcal{F}$ et $\mathcal{G}$ sont supposés être des
faisceaux pour la topologie fppf ; voir
Modules sur les sites, lemme \ref{sites-modules-lemma-internal-hom}.
D'après Faisceaux sur les champs, lemme \ref{stacks-sheaves-lemma-lqc-colimits},
nous voyons que $\SheafHom_{\mathcal{O}_\mathcal{X}}(\mathcal{F}, \mathcal{G})$
est localement quasi-cohérent.
D'après le lemme \ref{lemma-check-flat-comparison-on-etale-covering},
nous voyons que $\SheafHom_{\mathcal{O}_\mathcal{X}}(\mathcal{F}, \mathcal{G})$
possède la propriété de changement de base plat. Ainsi,
$\SheafHom_{\mathcal{O}_\mathcal{X}}(\mathcal{F}, \mathcal{G})$
est un objet de $\textit{LQCoh}^{fbc}(\mathcal{O}_\mathcal{X})$
et il est légitime de lui appliquer le foncteur $Q$ du lemme \ref{lemma-adjoint}.
Par la propriété universelle de $Q$, nous avons
$$
\Hom_\mathcal{X}(\mathcal{H},
Q(\SheafHom_{\mathcal{O}_\mathcal{X}}(\mathcal{F}, \mathcal{G}))) =
\Hom_\mathcal{X}(\mathcal{H},
\SheafHom_{\mathcal{O}_\mathcal{X}}(\mathcal{F}, \mathcal{G}))
$$
pour $\mathcal{H}$ quasi-cohérent ; la formule affichée
du lemme résulte donc de Modules sur les sites, lemme
\ref{sites-modules-lemma-internal-hom-adjoint-tensor}.
\end{proof}

\begin{lemma}
\label{lemma-flat-morphism-and-hom}
Soit $f : \mathcal{X} \to \mathcal{Y}$ un morphisme plat de champs algébriques.
Soit $\mathcal{F}$ un
$\mathcal{O}_\mathcal{Y}$-module de présentation finie et soit
$\mathcal{G}$ un $\mathcal{O}_\mathcal{Y}$-module quasi-cohérent.
Alors $f^*hom(\mathcal{F}, \mathcal{G}) = hom(f^*\mathcal{F}, f^*\mathcal{G})$
avec les notations du lemme \ref{lemma-internal-hom-fp-into-qcoh}.
\end{lemma}

\begin{proof}
Nous avons $f^*\SheafHom_{\mathcal{O}_\mathcal{Y}}(\mathcal{F}, \mathcal{G}) =
\SheafHom_{\mathcal{O}_\mathcal{X}}(f^*\mathcal{F}, f^*\mathcal{G})$
d'après Modules sur les sites, lemme \ref{sites-modules-lemma-pullback-internal-hom}.
(Observons que ce n'est pas à cette étape qu'est utilisée la platitude de $f$,
car le morphisme de topos annelés associé à $f$ est toujours plat ; voir
Faisceaux sur les champs, remarque \ref{stacks-sheaves-remark-flat}.)
Appliquons alors le lemme \ref{lemma-flat-morphism-and-Q} (et c'est ici que nous utilisons
la platitude de $f$).
\end{proof}








\section{Image directe des modules quasi-cohérents}
\label{section-pushforward-quasi-coherent}

\noindent
Soit $f : \mathcal{X} \to \mathcal{Y}$ un morphisme de champs algébriques.
Considérons l'image directe
$$
f_* :
\textit{Mod}(\mathcal{O}_\mathcal{X})
\longrightarrow
\textit{Mod}(\mathcal{O}_\mathcal{Y})
$$
Il s'avère que ce foncteur ne préserve presque jamais les sous-catégories
de faisceaux quasi-cohérents. Considérons par exemple le morphisme de schémas
$$
j : X = \mathbf{A}^2_k \setminus \{0\} \longrightarrow \mathbf{A}^2_k = Y.
$$
Il lui est associé le morphisme correspondant de champs algébriques
$$
f = j_{big} : \mathcal{X} = (\Sch/X)_{fppf} \to
(\Sch/Y)_{fppf} = \mathcal{Y}
$$
L'image directe $f_*\mathcal{O}_\mathcal{X}$ du faisceau structural a pour
sections globales $k[x, y]$. Ainsi, si $f_*\mathcal{O}_\mathcal{X}$ était
quasi-cohérent sur $\mathcal{Y}$, nous aurions
$f_*\mathcal{O}_\mathcal{X} = \mathcal{O}_\mathcal{Y}$. Cependant,
considérons $T = \Spec(k) \to \mathbf{A}^2_k = Y$ envoyé sur $0$.
Alors $\Gamma(T, f_*\mathcal{O}_\mathcal{X}) = 0$ parce que
$X \times_Y T = \emptyset$, tandis que $\Gamma(T, \mathcal{O}_\mathcal{Y}) = k$.
En revanche, pour tout morphisme plat $T \to Y$, nous avons l'égalité
$\Gamma(T, f_*\mathcal{O}_\mathcal{X}) = \Gamma(T, \mathcal{O}_\mathcal{Y})$
qui résulte de
Cohomologie des schémas, lemme \ref{coherent-lemma-flat-base-change-cohomology}
en utilisant que $j$ est quasi-compact et quasi-séparé.

\medskip\noindent
Soit $f : \mathcal{X} \to \mathcal{Y}$ un morphisme quasi-compact et
quasi-séparé de champs algébriques. Nous contournons le problème
mentionné ci-dessus à l'aide des trois observations suivantes :
\begin{enumerate}
\item $f_*$ préserve bien les modules localement quasi-cohérents
(lemme \ref{lemma-pushforward-locally-quasi-coherent}) ;
\item $f_*$ transforme un faisceau quasi-cohérent en un faisceau localement quasi-cohérent
dont les morphismes de comparaison plats sont des isomorphismes
(lemme \ref{lemma-flat-comparison}) ; et
\item les $\mathcal{O}_\mathcal{Y}$-modules localement quasi-cohérents
possédant la propriété de changement de base plat donnent des modules quasi-cohérents
sur une présentation de $\mathcal{Y}$, et donc des modules quasi-cohérents
sur $\mathcal{Y}$ ; voir
Faisceaux sur les champs, section
\ref{stacks-sheaves-section-quasi-coherent-algebraic-stacks}.
\end{enumerate}
Nous obtenons ainsi un foncteur
$$
f_{\QCoh, *} :
\QCoh(\mathcal{O}_\mathcal{X})
\longrightarrow
\QCoh(\mathcal{O}_\mathcal{Y})
$$
qui est adjoint à droite de
$f^* : \QCoh(\mathcal{O}_\mathcal{Y}) \to
\QCoh(\mathcal{O}_\mathcal{X})$
et tel que, de plus,
$$
\Gamma(y, f_*\mathcal{F}) = \Gamma(y, f_{\QCoh, *}\mathcal{F})
$$
pour tout $y \in \Ob(\mathcal{Y})$ tel que le $1$-morphisme
associé $y : V \to \mathcal{Y}$ soit plat ; voir
lemme \ref{lemma-direct-image-quasi-coherent-over-flat}.
De plus, une construction semblable produira des foncteurs
$R^if_{\QCoh, *}$.
Cependant, ces résultats ne suffiront pas à produire un foncteur
image directe totale (de complexes à faisceaux de cohomologie
quasi-cohérents).

\begin{proposition}
\label{proposition-direct-image-quasi-coherent}
Soit $f : \mathcal{X} \to \mathcal{Y}$ un morphisme quasi-compact et quasi-séparé
de champs algébriques. Le foncteur
$f^* : \QCoh(\mathcal{O}_\mathcal{Y}) \to
\QCoh(\mathcal{O}_\mathcal{X})$
admet un adjoint à droite
$$
f_{\QCoh, *} :
\QCoh(\mathcal{O}_\mathcal{X})
\longrightarrow
\QCoh(\mathcal{O}_\mathcal{Y})
$$
qui peut être défini comme le composé
$$
\QCoh(\mathcal{O}_\mathcal{X}) \to
\textit{LQCoh}^{fbc}(\mathcal{O}_\mathcal{X})
\xrightarrow{f_*} \textit{LQCoh}^{fbc}(\mathcal{O}_\mathcal{Y})
\xrightarrow{Q} \QCoh(\mathcal{O}_\mathcal{Y})
$$
où les foncteurs $f_*$ et $Q$ sont ceux de la
proposition \ref{proposition-loc-qcoh-flat-base-change}
et du
lemme \ref{lemma-adjoint}.
De plus, si nous définissons $R^if_{\QCoh, *}$ comme le composé
$$
\QCoh(\mathcal{O}_\mathcal{X}) \to
\textit{LQCoh}^{fbc}(\mathcal{O}_\mathcal{X})
\xrightarrow{R^if_*} \textit{LQCoh}^{fbc}(\mathcal{O}_\mathcal{Y})
\xrightarrow{Q} \QCoh(\mathcal{O}_\mathcal{Y})
$$
alors la suite de foncteurs $\{R^if_{\QCoh, *}\}_{i \geq 0}$
forme un $\delta$-foncteur cohomologique.
\end{proposition}

\begin{proof}
C'est une combinaison des résultats mentionnés dans l'énoncé.
L'adjonction se démontre comme suit. Soit $\mathcal{F}$
un $\mathcal{O}_\mathcal{X}$-module quasi-cohérent et soit
$\mathcal{G}$ un $\mathcal{O}_\mathcal{Y}$-module quasi-cohérent.
Alors nous avons
\begin{align*}
\Mor_{\QCoh(\mathcal{O}_\mathcal{X})}(f^*\mathcal{G}, \mathcal{F})
& =
\Mor_{\textit{LQCoh}^{fbc}(\mathcal{O}_\mathcal{Y})}
(\mathcal{G}, f_*\mathcal{F}) \\
& =
\Mor_{\QCoh(\mathcal{O}_\mathcal{Y})}(\mathcal{G}, Q(f_*\mathcal{F}))
\\
& =
\Mor_{\QCoh(\mathcal{O}_\mathcal{Y})}(\mathcal{G},
f_{\QCoh, *}\mathcal{F})
\end{align*}
la première égalité résultant de l'adjonction de $f_*$ et $f^*$ (pour des faisceaux
de modules arbitraires). D'après la
proposition \ref{proposition-loc-qcoh-flat-base-change}
nous voyons que $f_*\mathcal{F}$ est un objet de
$\textit{LQCoh}^{fbc}(\mathcal{O}_\mathcal{Y})$
(et peut être calculé pour la topologie fppf comme pour la topologie étale), et nous
obtenons la deuxième égalité d'après le lemme \ref{lemma-adjoint}. La troisième
égalité est la définition de $f_{\QCoh, *}$.

\medskip\noindent
Pour voir que $\{R^if_{\QCoh, *}\}_{i \geq 0}$ est un
$\delta$-foncteur cohomologique au sens de
Homologie, définition \ref{homology-definition-cohomological-delta-functor}
soit
$$
0 \to \mathcal{F}_1 \to \mathcal{F}_2 \to \mathcal{F}_3 \to 0
$$
une suite exacte courte de $\QCoh(\mathcal{O}_\mathcal{X})$.
Cette suite n'est pas nécessairement exacte dans
$\textit{Mod}(\mathcal{O}_\mathcal{X})$, mais nous savons qu'elle l'est
à des modules parasites près ; voir
lemme \ref{lemma-exact-sequence-quasi-coherent-parasitic-cohomology}.
Nous pouvons donc décomposer la suite en suites exactes courtes
$$
\begin{matrix}
0 \to \mathcal{P}_1 \to \mathcal{F}_1 \to \mathcal{I}_2 \to 0 \\
0 \to \mathcal{I}_2 \to \mathcal{F}_2 \to \mathcal{Q}_2 \to 0 \\
0 \to \mathcal{P}_2 \to \mathcal{Q}_2 \to \mathcal{I}_3 \to 0 \\
0 \to \mathcal{I}_3 \to \mathcal{F}_3 \to \mathcal{P}_3 \to 0
\end{matrix}
$$
de $\textit{Mod}(\mathcal{O}_\mathcal{X})$, les $\mathcal{P}_i$ étant parasites.
Notons que chacun des faisceaux
$\mathcal{P}_j$, $\mathcal{I}_j$, $\mathcal{Q}_j$ est un objet de
$\textit{LQCoh}^{fbc}(\mathcal{O}_\mathcal{X})$ ; voir
proposition \ref{proposition-loc-qcoh-flat-base-change}.
En appliquant $R^if_*$, nous obtenons des suites exactes longues
$$
\begin{matrix}
0 \to f_*\mathcal{P}_1 \to f_*\mathcal{F}_1 \to f_*\mathcal{I}_2 \to
R^1f_*\mathcal{P}_1 \to \ldots \\
0 \to f_*\mathcal{I}_2 \to f_*\mathcal{F}_2 \to f_*\mathcal{Q}_2 \to
R^1f_*\mathcal{I}_2 \to \ldots \\
0 \to f_*\mathcal{P}_2 \to f_*\mathcal{Q}_2 \to f_*\mathcal{I}_3 \to
R^1f_*\mathcal{P}_2 \to \ldots \\
0 \to f_*\mathcal{I}_3 \to f_*\mathcal{F}_3 \to f_*\mathcal{P}_3 \to
R^1f_*\mathcal{I}_3 \to \ldots
\end{matrix}
$$
où tous les termes sont des objets de
$\textit{LQCoh}^{fbc}(\mathcal{O}_\mathcal{Y})$ d'après la
proposition \ref{proposition-loc-qcoh-flat-base-change}.
D'après le
lemme \ref{lemma-pushforward-parasitic},
les faisceaux $R^if_*\mathcal{P}_j$ sont parasites et s'annulent donc après application
du foncteur $Q$ ; voir le
lemme \ref{lemma-adjoint-kernel-parasitic}.
Puisque $Q$ est exact, les morphismes
$$
Q(R^if_*\mathcal{F}_3)
\cong
Q(R^if_*\mathcal{I}_3)
\cong
Q(R^if_*\mathcal{Q}_2)
\rightarrow
Q(R^{i + 1}f_*\mathcal{I}_2)
\cong
Q(R^{i + 1}f_*\mathcal{F}_1)
$$
peuvent servir de morphisme de connexion qui fait de la famille de foncteurs
$\{R^if_{\QCoh, *}\}_{i \geq 0}$
un $\delta$-foncteur cohomologique.
\end{proof}

\begin{lemma}
\label{lemma-direct-image-quasi-coherent-over-flat}
Soit $f : \mathcal{X} \to \mathcal{Y}$ un morphisme quasi-compact et quasi-séparé
de champs algébriques. Soit $y : V \to \mathcal{Y}$ dans $\Ob(\mathcal{Y})$,
avec $y$ plat. Soit $\mathcal{F}$ dans
$\QCoh(\mathcal{O}_\mathcal{X})$.
Alors $(f_*\mathcal{F})(y) = (f_{\QCoh, *}\mathcal{F})(y)$
et $(R^if_*\mathcal{F})(y) = (R^if_{\QCoh, *}\mathcal{F})(y)$
pour tout $i \in \mathbf{Z}$.
\end{lemma}

\begin{proof}
Cela résulte de la construction des foncteurs $R^if_{\QCoh, *}$ dans la
proposition \ref{proposition-direct-image-quasi-coherent},
de la définition des modules parasites dans définition \ref{definition-parasitic}
et du lemme \ref{lemma-adjoint-kernel-parasitic}, partie (2).
\end{proof}

\begin{remark}
\label{remark-direct-image-quasi-coherent-tensor}
Soit $f : \mathcal{X} \to \mathcal{Y}$ un morphisme quasi-compact et quasi-séparé
de champs algébriques. Soient $\mathcal{F}$ et $\mathcal{G}$ dans
$\QCoh(\mathcal{O}_\mathcal{X})$. Il existe alors un diagramme commutatif
canonique
$$
\xymatrix{
f_{\QCoh, *}\mathcal{F}
\otimes_{\mathcal{O}_\mathcal{Y}}
f_{\QCoh, *}\mathcal{G} \ar[r] \ar[d] &
f_*\mathcal{F}
\otimes_{\mathcal{O}_\mathcal{Y}}
f_*\mathcal{G} \ar[d]^c \\
f_{\QCoh, *}(\mathcal{F}
\otimes_{\mathcal{O}_\mathcal{X}}
\mathcal{G}) \ar[r] &
f_*(\mathcal{F}
\otimes_{\mathcal{O}_\mathcal{X}}
\mathcal{G})
}
$$
La flèche verticale $c$ de droite est le cup-produit relatif naïf
(en degré $0$) ;
voir Cohomologie sur les sites, section \ref{sites-cohomology-section-cup-product}.
La source et le but de $c$ appartiennent à
$\textit{LQCoh}^{fbc}(\mathcal{O}_\mathcal{X})$ ; voir la
proposition \ref{proposition-loc-qcoh-flat-base-change}.
En appliquant $Q$ à $c$, nous obtenons la flèche verticale de gauche,
car $Q$ commute aux produits tensoriels ; voir la
remarque \ref{remark-QCoh-tensor}.
Cette construction est fonctorielle en $\mathcal{F}$ et $\mathcal{G}$.
\end{remark}

\begin{lemma}
\label{lemma-leray}
Soit $f : \mathcal{X} \to \mathcal{Y}$
un morphisme quasi-compact et quasi-séparé de champs algébriques.
Soit $\mathcal{F}$ un faisceau quasi-cohérent sur $\mathcal{X}$. Alors
il existe une suite spectrale dont la page $E_2$ est
$$
E_2^{p, q} = H^p(\mathcal{Y}, R^qf_{\QCoh, *}\mathcal{F})
$$
et qui converge vers $H^{p + q}(\mathcal{X}, \mathcal{F})$.
\end{lemma}

\begin{proof}
D'après Cohomologie sur les sites, lemme \ref{sites-cohomology-lemma-Leray},
la suite spectrale de Leray dont la page est
$$
E_2^{p, q} = H^p(\mathcal{Y}, R^qf_*\mathcal{F})
$$
converge vers $H^{p + q}(\mathcal{X}, \mathcal{F})$. Le noyau et le conoyau
du morphisme d'adjonction
$$
R^qf_{\QCoh, *}\mathcal{F} \longrightarrow R^qf_*\mathcal{F}
$$
sont des modules parasites sur $\mathcal{Y}$
(lemme \ref{lemma-adjoint-kernel-parasitic})
et ont donc une cohomologie nulle
(lemme \ref{lemma-pushforward-parasitic}).
Il en résulte formellement que
$H^p(\mathcal{Y}, R^qf_{\QCoh, *}\mathcal{F}) =
H^p(\mathcal{Y}, R^qf_*\mathcal{F})$, ce qui conclut.
\end{proof}

\begin{lemma}
\label{lemma-relative-leray}
Soient $f : \mathcal{X} \to \mathcal{Y}$ et $g : \mathcal{Y} \to \mathcal{Z}$
des morphismes quasi-compacts et quasi-séparés de champs algébriques.
Soit $\mathcal{F}$ un faisceau quasi-cohérent sur $\mathcal{X}$. Alors
il existe une suite spectrale dont la page $E_2$ est
$$
E_2^{p, q} = R^pg_{\QCoh, *}(R^qf_{\QCoh, *}\mathcal{F})
$$
et qui converge vers $R^{p + q}(g \circ f)_{\QCoh, *}\mathcal{F}$.
\end{lemma}

\begin{proof}
D'après Cohomologie sur les sites, lemme \ref{sites-cohomology-lemma-relative-Leray},
la suite spectrale de Leray dont la page est
$$
E_2^{p, q} = R^pg_*(R^qf_*\mathcal{F})
$$
converge vers $R^{p + q}(g \circ f)_*\mathcal{F}$. D'après les résultats de la
proposition \ref{proposition-loc-qcoh-flat-base-change}
les termes de cette suite spectrale sont tous des objets de
$\textit{LQCoh}^{fbc}(\mathcal{O}_\mathcal{Z})$. En appliquant le foncteur exact
$Q_\mathcal{Z} : \textit{LQCoh}^{fbc}(\mathcal{O}_\mathcal{Z}) \to
\QCoh(\mathcal{O}_\mathcal{Z})$, nous obtenons une suite spectrale dans
$\QCoh(\mathcal{O}_\mathcal{Z})$ convergeant vers
$R^{p + q}(g \circ f)_{\QCoh, *}\mathcal{F}$. Le résultat
en découle donc si nous pouvons montrer que
$$
Q_\mathcal{Z}(R^pg_*(R^qf_*\mathcal{F})) =
Q_\mathcal{Z}(R^pg_*(Q_\mathcal{X}(R^qf_*\mathcal{F}))
$$
Cela résulte du fait que le noyau et le conoyau du morphisme
$$
Q_\mathcal{X}(R^qf_*\mathcal{F}) \longrightarrow R^qf_*\mathcal{F}
$$
sont parasites (lemme \ref{lemma-adjoint-kernel-parasitic}) et que
$R^pg_*$ transforme les modules parasites en modules parasites
(lemme \ref{lemma-pushforward-parasitic}).
\end{proof}

\noindent
Pour terminer cette section, explicitons les suites spectrales
associées à un recouvrement lisse par un schéma. Comparer avec
Faisceaux sur les champs, sections \ref{stacks-sheaves-section-cohomology} et
\ref{stacks-sheaves-section-higher-direct-images}.

\begin{proposition}
\label{proposition-smooth-covering-compute-cohomology}
Soit $f : \mathcal{U} \to \mathcal{X}$ un morphisme de champs algébriques.
Supposons que $f$ soit représentable par des espaces algébriques, surjectif, plat et
localement de présentation finie. Soit $\mathcal{F}$ un
$\mathcal{O}_\mathcal{X}$-module quasi-cohérent. Il existe alors une suite spectrale
$$
E_2^{p, q} = H^q(\mathcal{U}_p, f_p^*\mathcal{F})
\Rightarrow
H^{p + q}(\mathcal{X}, \mathcal{F})
$$
où $f_p$ est le morphisme
$\mathcal{U} \times_\mathcal{X} \ldots \times_\mathcal{X} \mathcal{U} \to
\mathcal{X}$ (avec $p + 1$ facteurs).
\end{proposition}

\begin{proof}
C'est un cas particulier de
Faisceaux sur les champs, proposition
\ref{stacks-sheaves-proposition-smooth-covering-compute-cohomology}.
\end{proof}

\begin{proposition}
\label{proposition-smooth-covering-compute-direct-image}
Soient $f : \mathcal{U} \to \mathcal{X}$ et $g : \mathcal{X} \to \mathcal{Y}$
des morphismes composables de champs algébriques.
Supposons que
\begin{enumerate}
\item $f$ soit représentable par des espaces algébriques, surjectif,
plat, localement de présentation finie, quasi-compact et quasi-séparé, et
\item $g$ soit quasi-compact et quasi-séparé.
\end{enumerate}
Si $\mathcal{F}$ appartient à $\QCoh(\mathcal{O}_\mathcal{X})$, alors
il existe une suite spectrale
$$
E_2^{p, q} = R^q(g \circ f_p)_{\QCoh, *}f_p^*\mathcal{F}
\Rightarrow
R^{p + q}g_{\QCoh, *}\mathcal{F}
$$
dans $\QCoh(\mathcal{O}_\mathcal{Y})$.
\end{proposition}

\begin{proof}
Notons que chacun des morphismes
$f_p : \mathcal{U} \times_\mathcal{X} \ldots \times_\mathcal{X} \mathcal{U} \to
\mathcal{X}$ est quasi-compact et quasi-séparé ; ainsi, $g \circ f_p$
est quasi-compact et quasi-séparé, et l'assertion a donc un sens
(c'est-à-dire que les foncteurs $R^q(g \circ f_p)_{\QCoh, *}$ sont définis).
Il existe une suite spectrale
$$
E_2^{p, q} = R^q(g \circ f_p)_*f_p^{-1}\mathcal{F}
\Rightarrow
R^{p + q}g_*\mathcal{F}
$$
d'après Faisceaux sur les champs, proposition
\ref{stacks-sheaves-proposition-smooth-covering-compute-direct-image}.
En appliquant le foncteur exact
$Q_\mathcal{Y} : \textit{LQCoh}^{fbc}(\mathcal{O}_\mathcal{Y}) \to
\QCoh(\mathcal{O}_\mathcal{Y})$, on obtient la suite spectrale voulue dans
$\QCoh(\mathcal{O}_\mathcal{Y})$.
\end{proof}







\section{Remarques supplémentaires sur les modules quasi-cohérents}
\label{section-further-remarks}

\noindent
Dans cette section, nous réunissons quelques résultats qui aident à comprendre comment
utiliser les modules quasi-cohérents sur les champs algébriques.

\medskip\noindent
Soit $f : \mathcal{U} \to \mathcal{X}$ un morphisme de champs algébriques.
Supposons que $\mathcal{U}$ soit représenté par l'espace algébrique $U$.
Considérons le foncteur
$$
a :
\textit{Mod}(\mathcal{X}_\etale, \mathcal{O}_\mathcal{X})
\longrightarrow
\textit{Mod}(U_\etale, \mathcal{O}_U),\quad
\mathcal{F}
\longmapsto
f^*\mathcal{F}|_{U_\etale}
$$
obtenu en faisant suivre l'image inverse (Faisceaux sur les champs, section
\ref{stacks-sheaves-section-modules}) de la restriction
(Faisceaux sur les champs, section
\ref{stacks-sheaves-section-restriction-algebraic-spaces}).
En appliquant ce foncteur aux modules localement quasi-cohérents, on obtient un foncteur
$$
b : \textit{LQCoh}(\mathcal{O}_\mathcal{X})
\longrightarrow
\QCoh(U_\etale, \mathcal{O}_U)
$$
Voir Faisceaux sur les champs, lemmes \ref{stacks-sheaves-lemma-pullback-lqc} et
\ref{stacks-sheaves-lemma-compare-locally-quasi-coherent}.
On peut encore restreindre notre foncteur à des sous-catégories plus petites et obtenir
$$
c :
\textit{LQCoh}^{fbc}(\mathcal{O}_\mathcal{X})
\longrightarrow
\QCoh(U_\etale, \mathcal{O}_U)
$$
ainsi que
$$
d :
\QCoh(\mathcal{O}_\mathcal{X})
\longrightarrow
\QCoh(U_\etale, \mathcal{O}_U)
$$
Ces foncteurs possèdent les propriétés suivantes :\footnote{Nous suggérons
d'établir pourquoi ces assertions sont vraies sur un coin de table plutôt que de
suivre les références indiquées.}
\begin{enumerate}
\item Le foncteur $a$ est exact. En effet, l'image inverse $f^* = f^{-1}$ est exacte
(Faisceaux sur les champs, section \ref{stacks-sheaves-section-modules}),
et la restriction à $U_\etale$ est exacte ; voir
Faisceaux sur les champs, équation (\ref{stacks-sheaves-equation-restrict}).
\item Le foncteur $b$ est exact. En effet, d'après
Faisceaux sur les champs, lemme \ref{stacks-sheaves-lemma-lqc-colimits}
le foncteur d'inclusion $\textit{LQCoh}(\mathcal{O}_\mathcal{X}) \to
\textit{Mod}(\mathcal{X}_\etale, \mathcal{O}_\mathcal{X})$ est exact.
\item Le foncteur $c$ est exact. En effet, d'après la
proposition \ref{proposition-loc-qcoh-flat-base-change}
le foncteur d'inclusion $\textit{LQCoh}^{fbc}(\mathcal{O}_\mathcal{X}) \to
\textit{Mod}(\mathcal{X}_\etale, \mathcal{O}_\mathcal{X})$ est exact.
\item Le foncteur $d$ est exact à droite, mais n'est pas exact en général.
En effet, d'après Faisceaux sur les champs, lemme
\ref{stacks-sheaves-lemma-qc-colimits}
le foncteur d'inclusion $\QCoh(\mathcal{O}_\mathcal{X}) \to
\textit{Mod}(\mathcal{X}_\etale, \mathcal{O}_\mathcal{X})$ est exact à
droite. Nous omettons de donner un exemple montrant qu'il n'est pas exact.
\item Si $f$ est plat, alors $d$ est exact. Cela résulte de la combinaison du
Lemme \ref{lemma-flat-pullback-quasi-coherent} et de
Faisceaux sur les champs, lemme \ref{stacks-sheaves-lemma-compare-quasi-coherent}.
\item Si $f$ est plat, alors $c$ annule les objets parasites.
En effet, $f^*$ préserve les objets parasites d'après le Lemme \ref{lemma-parasitic}.
Pour tout schéma $V$ étale sur $U$, et donc plat sur $\mathcal{X}$,
on voit alors que
$0 = f^*\mathcal{F}|_{V_\etale} = c(\mathcal{F})|_{V_\etale}$
par compatibilité de la restriction à la localisation étale
(Faisceaux sur les champs, remarque \ref{stacks-sheaves-remark-compare}).
Ainsi, manifestement, $c(\mathcal{F}) = 0$.
\item Si $f$ est plat, alors $c = d \circ Q$. En effet, le noyau et le
conoyau de $Q(\mathcal{F}) \to \mathcal{F}$ sont parasites d'après le
Lemme \ref{lemma-adjoint-kernel-parasitic}. Puisque
$c$ est exact (3) et annule les objets parasites (6),
on voit que l'image par $c$ du morphisme $Q(\mathcal{F}) \to \mathcal{F}$
est un isomorphisme.
\item Les foncteurs $a, b, c, d$ commutent aux limites inductives
et aux sommes directes arbitraires. Cela vaut pour $f^*$ et la restriction,
qui sont des adjoints à gauche, et vaut donc pour $a$. Le résultat pour
$b$, $c$, $d$ découle alors des références données ci-dessus.
\item Les foncteurs $a, b, c, d$ commutent aux produits tensoriels.
\item Si $f$ est plat et surjectif, si $\mathcal{F}$ appartient à
$\textit{LQCoh}^{fbc}(\mathcal{O}_\mathcal{X})$ et si $c(\mathcal{F}) = 0$,
alors $\mathcal{F}$ est parasite. En effet, (7) donne $d(Q(\mathcal{F})) = 0$.
On peut supposer que $U$ est un schéma, par compatibilité de la restriction
avec la localisation étale (voir la référence ci-dessus). Le
Lemme \ref{lemma-quasi-coherent-check-exact}, appliqué à $0 \to Q(\mathcal{F})$
et au morphisme $f : U \to \mathcal{X}$, montre alors que $Q(\mathcal{F}) = 0$.
Ainsi, $\mathcal{F}$ est parasite d'après le Lemme \ref{lemma-adjoint-kernel-parasitic}.
\item Si $f$ est plat et surjectif, alors le foncteur $d$
reflète l'exactitude. Plus précisément, soit $\mathcal{F}^\bullet$ un complexe de
$\QCoh(\mathcal{O}_\mathcal{X})$. Alors $\mathcal{F}^\bullet$ est exact
dans $\QCoh(\mathcal{O}_\mathcal{X})$ si et seulement si $d(\mathcal{F}^\bullet)$
est exact. Nous avons vu une implication en (5). Réciproquement,
supposons que $H^i(d(\mathcal{F}^\bullet)) = 0$. Alors
$\mathcal{G} = H^i(\mathcal{F}^\bullet)$ est un objet de
$\QCoh(\mathcal{O}_\mathcal{X})$ tel que $d(\mathcal{G}) = 0$.
Ainsi, $\mathcal{G}$ est à la fois quasi-cohérent et parasite d'après (10),
donc égal à $0$, par exemple d'après la Remarque \ref{remark-bousfield-colocalization}.
\item
\label{item-hom-restriction}
Si $f$ est plat,
$\mathcal{F}, \mathcal{G} \in \Ob(\QCoh(\mathcal{O}_\mathcal{X}))$,
et si $\mathcal{F}$ est de présentation finie,
alors on a
$$
d(hom(\mathcal{F}, \mathcal{G})) =
\SheafHom_{\mathcal{O}_U}(d(\mathcal{F}), d(\mathcal{G}))
$$
avec les notations du Lemme \ref{lemma-internal-hom-fp-into-qcoh}.
La manière la plus simple de le voir est peut-être la suivante :
\begin{align*}
d(hom(\mathcal{F}, \mathcal{G}))
& =
d(Q(\SheafHom_{\mathcal{O}_\mathcal{X}}(\mathcal{F}, \mathcal{G}))) \\
& =
c(\SheafHom_{\mathcal{O}_\mathcal{X}}(\mathcal{F}, \mathcal{G})) \\
& =
f^*\SheafHom_{\mathcal{O}_\mathcal{X}}(\mathcal{F}, \mathcal{G})|_{U_\etale} \\
& =
\SheafHom_{\mathcal{O}_\mathcal{U}}(f^*\mathcal{F},
f^*\mathcal{G})|_{U_\etale} \\
& =
\SheafHom_{\mathcal{O}_U}(f^*\mathcal{F}|_{U_\etale},
f^*\mathcal{G}|_{U_\etale})
\end{align*}
La première égalité résulte de la construction de $hom$. La deuxième
résulte de (7). La troisième résulte de la définition de $c$. La quatrième résulte de
Modules sur les sites, lemme \ref{sites-modules-lemma-pullback-internal-hom}.
La dernière égalité résulte de la même référence appliquée au morphisme plat
de topos annelés $i_U (U_\etale, \mathcal{O}_U) \to
(\mathcal{U}_\etale, \mathcal{O}_\mathcal{U})$ de
Faisceaux sur les champs, lemme \ref{stacks-sheaves-lemma-compare}.
\item Ajouter ici d'autres éléments.
\end{enumerate}






\section{Limites inductives et cohomologie}
\label{section-colimits}

\noindent
Le lemme suivant s'applique notamment aux diagrammes de faisceaux
quasi-cohérents.

\begin{lemma}
\label{lemma-colimits}
Soit $\mathcal{X}$ un champ algébrique quasi-compact et quasi-séparé.
Alors
$$
\colim_i H^p(\mathcal{X}, \mathcal{F}_i)
\longrightarrow
H^p(\mathcal{X}, \colim_i \mathcal{F}_i)
$$
est un isomorphisme pour tout diagramme filtrant de faisceaux abéliens sur
$\mathcal{X}$. Il en va de même pour les faisceaux abéliens sur $\mathcal{X}_\etale$
lorsque la cohomologie est prise pour la topologie étale.
\end{lemma}

\begin{proof}
Soit $\tau = fppf$, resp.\ $\tau = \etale$. Le lemme résulte de
Cohomologie sur les sites, lemme \ref{sites-cohomology-lemma-colim-global},
appliqué au site $\mathcal{X}_\tau$. Pour en vérifier les hypothèses, nous
utilisons Cohomologie sur les sites, remarque \ref{sites-cohomology-remark-colim-global}.
Soit en effet $\mathcal{B} \subset \Ob(\mathcal{X}_\tau)$ l'ensemble
des objets au-dessus de schémas affines. Autrement dit, un élément
de $\mathcal{B}$ est un morphisme $x : U \to \mathcal{X}$ où $U$ est affine.
Vérifions successivement chacune des conditions (1) -- (4) de la remarque :
\begin{enumerate}
\item Puisque $\mathcal{X}$ est quasi-compact, il existe un morphisme
lisse et surjectif $x : U \to \mathcal{X}$ où $U$ est affine
(Propriétés des champs, lemme
\ref{stacks-properties-lemma-quasi-compact-stack}).
Alors $h_x^\# \to *$ est un morphisme surjectif de faisceaux sur $\mathcal{X}_\tau$.
\item Puisque les recouvrements dans $\mathcal{X}_\tau$ sont des recouvrements
fppf, resp.\ étales, tout recouvrement de $U \in \mathcal{B}$ est raffiné
par un recouvrement fppf affine fini ; voir
Topologies, lemme \ref{topologies-lemma-fppf-affine},
resp.\ lemme \ref{topologies-lemma-etale-affine}.
\item Soient $x : U \to \mathcal{X}$ et $x' : U' \to \mathcal{X}$ dans
$\mathcal{B}$. Le produit $h_x^\# \times h_{x'}^\#$ dans
$\Sh(\mathcal{X}_\tau)$ est égal au faisceau sur $\mathcal{X}_\tau$
déterminé par l'espace algébrique $W = U \times_{x, \mathcal{X}, x'} U'$ au-dessus de
$\mathcal{X}$ : pour un objet $y : V \to \mathcal{X}$ de
$\mathcal{X}_\tau$, on a
$(h_x^\# \times h_{x'}^\#)(y) = \{f : V \to W \mid y =
x \circ \text{pr}_1 \circ f = x' \circ \text{pr}_2 \circ f\}$.
L'espace algébrique $W$ est quasi-compact, car $\mathcal{X}$
est quasi-séparé ; voir, par exemple, Morphismes de champs, lemme
\ref{stacks-morphisms-lemma-quasi-compact-quasi-separated-permanence}
On peut donc choisir un schéma affine $U''$ et
un morphisme étale surjectif $U'' \to W$. Notons $x'' : U'' \to \mathcal{X}$
le composé de $U'' \to W$ et $W \to \mathcal{X}$. Alors
$h_{x''}^\# \to h_x^\# \times h_{x'}^\#$ est surjectif, comme voulu.
\item Soient $x : U \to \mathcal{X}$ et $x' : U' \to \mathcal{X}$ dans
$\mathcal{B}$. Soient $a, b : U \to U'$ des morphismes au-dessus de $\mathcal{X}$,
c'est-à-dire $a, b  : x \to x'$ des morphismes dans $\mathcal{X}_\tau$.
Alors l'égalisateur de $h_a$ et $h_b$ est représenté par
l'égalisateur de $a, b : U \to U'$, qui est un
schéma affine au-dessus de $\mathcal{X}$ et appartient donc à $\mathcal{B}$.
\end{enumerate}
Cela achève la démonstration.
\end{proof}

\begin{lemma}
\label{lemma-colimit-cohomology}
Soit $f : \mathcal{X} \to \mathcal{Y}$ un morphisme quasi-compact et quasi-séparé
de champs algébriques. Soit $\mathcal{F} = \colim \mathcal{F}_i$
une limite inductive filtrante de faisceaux abéliens sur $\mathcal{X}$.
Alors, pour tout $p \geq 0$, on a
$$
R^pf_*\mathcal{F} = \colim R^pf_*\mathcal{F}_i.
$$
Il en va de même pour les faisceaux abéliens sur $\mathcal{X}_\etale$
lorsque les images directes supérieures sont prises pour la topologie étale.
\end{lemma}

\begin{proof}
Nous le démontrons pour la topologie fppf ; la démonstration pour la topologie
étale est la même. Rappelons que $R^if_*\mathcal{F}$ est le faisceau sur
$\mathcal{Y}_{fppf}$ associé au préfaisceau
$$
(y : V \to \mathcal{Y})
\longmapsto
H^i(V \times_{y, \mathcal{Y}} \mathcal{X}, \text{pr}^{-1}\mathcal{F})
$$
Voir
Faisceaux sur les champs, lemme \ref{stacks-sheaves-lemma-pushforward-restriction}.
Rappelons que la limite inductive est le faisceau associé à la limite inductive des préfaisceaux.
Lorsque $V$ est affine, le produit fibré $V \times_\mathcal{Y} \mathcal{X}$
est quasi-compact et quasi-séparé. On peut donc appliquer le
Lemme \ref{lemma-colimits} à $H^p(V \times_\mathcal{Y} \mathcal{X}, -)$
lorsque $V$ est affine. Puisque tout $V$ possède un recouvrement fppf par des
objets affines, cela démontre le lemme. Nous omettons certains détails.
\end{proof}

\begin{lemma}
\label{lemma-quasi-coherent-pushforward-direct-sums}
Soit $f : \mathcal{X} \to \mathcal{Y}$ un morphisme quasi-compact et quasi-séparé
de champs algébriques. Le foncteur $f_{\QCoh, *}$
et les foncteurs $R^if_{\QCoh, *}$ commutent aux sommes directes
et aux limites inductives filtrantes.
\end{lemma}

\begin{proof}
Les foncteurs $f_*$ et $R^if_*$ commutent aux sommes directes et aux limites
inductives filtrantes sur tous les modules d'après le Lemme \ref{lemma-colimit-cohomology}.
Le lemme en résulte, car $f_{\QCoh, *} = Q \circ f_*$ et
$R^if_{\QCoh, *} = Q \circ R^if_*$, et $Q$ commute à
toutes les limites inductives ; voir le Lemme \ref{lemma-adjoint-kernel-parasitic}.
\end{proof}

\begin{lemma}
\label{lemma-quasi-coherent-pushforward-affine}
Soit $f : \mathcal{X} \to \mathcal{Y}$ un morphisme affine de champs
algébriques. Les foncteurs $R^if_{\QCoh, *}$, $i > 0$, sont nuls, et le foncteur
$f_{\QCoh, *}$ est exact et commute aux sommes directes et à toutes les limites inductives.
\end{lemma}

\begin{proof}
Puisque $R^if_{\QCoh, *} = Q \circ R^if_*$, l'annulation résulte du
Lemme \ref{lemma-loc-qcoh-fbc-affine-direct-image}.
Cette annulation entraîne que $f_{\QCoh, *}$ est exact, car
$\{R^if_{\QCoh, *}\}_{i \geq 0}$ forme un $\delta$-foncteur ; voir la
proposition \ref{proposition-direct-image-quasi-coherent}.
Puis $f_{\QCoh, *}$ commute aux sommes directes, par exemple d'après le
Lemme \ref{lemma-quasi-coherent-pushforward-direct-sums}.
Un foncteur exact qui commute aux sommes directes commute
à toutes les limites inductives.
\end{proof}

\noindent
Le lemme suivant affirme que les modules de présentation finie se comportent
comme prévu sur les champs algébriques quasi-compacts et quasi-séparés.

\begin{lemma}
\label{lemma-finite-presentation-quasi-compact-colimit}
Soit $\mathcal{X}$ un champ algébrique quasi-compact et quasi-séparé.
Soit $I$ un ensemble ordonné filtrant et soit $(\mathcal{F}_i, \varphi_{ii'})$ un
système indexé par $I$ de $\mathcal{O}_\mathcal{X}$-modules. Soit $\mathcal{G}$ un
$\mathcal{O}_\mathcal{X}$-module de présentation finie. Alors on a
$$
\colim_i \Hom_\mathcal{X}(\mathcal{G}, \mathcal{F}_i)
=
\Hom_\mathcal{X}(\mathcal{G}, \colim_i \mathcal{F}_i).
$$
En particulier, $\Hom_\mathcal{X}(\mathcal{G}, -)$ commute aux limites
inductives filtrantes dans $\QCoh(\mathcal{O}_\mathcal{X})$.
\end{lemma}

\begin{proof}
L'égalité affichée est un cas particulier de Modules sur les sites, lemme
\ref{sites-modules-lemma-finite-presentation-quasi-compact-colimit}.
Pour l'appliquer, il faut vérifier les hypothèses de la partie (4) de
Sites, lemme \ref{sites-lemma-directed-colimits-global-sections},
pour le site $\mathcal{X}_{fppf}$. À cette fin, nous vérifions
les hypothèses (2)(a), (2)(b), (2)(c) de
Sites, remarque \ref{sites-remark-stronger-conditions}.
Soit en effet $\mathcal{B} \subset \Ob(\mathcal{X}_{fppf})$ l'ensemble
des objets au-dessus de schémas affines. Autrement dit, un élément
de $\mathcal{B}$ est un morphisme $x : U \to \mathcal{X}$ où $U$ est affine.
Vérifions successivement chacune des conditions (2)(a), (2)(b) et (2)(c)
de la remarque :
\begin{enumerate}
\item Puisque $\mathcal{X}$ est quasi-compact, il existe un morphisme
lisse et surjectif $x : U \to \mathcal{X}$ où $U$ est affine
(Propriétés des champs, lemme
\ref{stacks-properties-lemma-quasi-compact-stack}).
Alors $h_x^\# \to *$ est un morphisme surjectif de faisceaux sur $\mathcal{X}_{fppf}$.
\item Puisque les recouvrements dans $\mathcal{X}_{fppf}$ sont des recouvrements
fppf, tout recouvrement de $U \in \mathcal{B}$ est raffiné
par un recouvrement fppf affine fini ; voir
Topologies, lemme \ref{topologies-lemma-fppf-affine}.
\item Soient $x : U \to \mathcal{X}$ et $x' : U' \to \mathcal{X}$ dans
$\mathcal{B}$. Le produit $h_x^\# \times h_{x'}^\#$ dans
$\Sh(\mathcal{X}_{fppf})$ est égal au faisceau sur $\mathcal{X}_{fppf}$
déterminé par l'espace algébrique $W = U \times_{x, \mathcal{X}, x'} U'$ au-dessus de
$\mathcal{X}$ : pour un objet $y : V \to \mathcal{X}$ de
$\mathcal{X}_{fppf}$, on a
$(h_x^\# \times h_{x'}^\#)(y) = \{f : V \to W \mid y =
x \circ \text{pr}_1 \circ f = x' \circ \text{pr}_2 \circ f\}$.
L'espace algébrique $W$ est quasi-compact, car $\mathcal{X}$
est quasi-séparé ; voir, par exemple, Morphismes de champs, lemme
\ref{stacks-morphisms-lemma-quasi-compact-quasi-separated-permanence}
On peut donc choisir un schéma affine $U''$ et
un morphisme étale surjectif $U'' \to W$. Notons $x'' : U'' \to \mathcal{X}$
le composé de $U'' \to W$ et $W \to \mathcal{X}$. Alors
$h_{x''}^\# \to h_x^\# \times h_{x'}^\#$ est surjectif, comme voulu.
\end{enumerate}
Pour la dernière assertion, observons que le foncteur d'inclusion
$\QCoh(\mathcal{O}_X) \to \textit{Mod}(\mathcal{O}_X)$
commute aux limites inductives et que les modules de présentation finie
sont quasi-cohérents. Voir Faisceaux sur les champs, lemme
\ref{stacks-sheaves-lemma-quasi-coherent-algebraic-stack}.
\end{proof}



\section{Les sites lisse-étale et plat-fppf}
\label{section-lisse-etale}

\noindent
Dans l'ouvrage \cite{LM-B}, beaucoup des résultats ci-dessus sont démontrés à l'aide du
site lisse-étale d'un champ algébrique. Nous définissons ici ce site.
Dans Exemples, section \ref{examples-section-lisse-etale-not-functorial},
nous montrons que le site lisse-étale n'est pas fonctoriel.
Nous définissons aussi son analogue, le site plat-fppf, qui est mieux adapté
au développement des champs algébriques adopté dans le projet Champs
(car nous prenons la topologie fppf comme topologie de base). Bien entendu, le
site plat-fppf n'est pas non plus fonctoriel.

\begin{definition}
\label{definition-lisse-etale}
Soit $\mathcal{X}$ un champ algébrique.
\begin{enumerate}
\item Le {\it site lisse-étale} de $\mathcal{X}$ est la sous-catégorie pleine
$\mathcal{X}_{lisse,\etale}$\footnote{Dans la littérature, le
site est noté $\text{Lis-\'et}(\mathcal{X})$ ou
$\text{Lis-Et}(\mathcal{X})$, et le topos associé est noté
$\mathcal{X}_{\text{lis-\'e}t}$ ou $\mathcal{X}_{\text{lis-et}}$.
Dans le projet Champs, notre convention consiste à nommer le site et à
noter le topos correspondant $\Sh(\mathcal{C})$.} de $\mathcal{X}$
dont les objets sont les $x \in \Ob(\mathcal{X})$ au-dessus d'un schéma $U$
tels que $x : U \to \mathcal{X}$ soit lisse. Un recouvrement de
$\mathcal{X}_{lisse,\etale}$ est une famille de morphismes
$\{x_i \to x\}_{i \in I}$ de $\mathcal{X}_{lisse,\etale}$
qui forme un recouvrement de $\mathcal{X}_\etale$.
\item Le {\it site plat-fppf} de $\mathcal{X}$ est la sous-catégorie pleine
$\mathcal{X}_{flat,fppf}$ de $\mathcal{X}$
dont les objets sont les $x \in \Ob(\mathcal{X})$ au-dessus d'un schéma $U$
tels que $x : U \to \mathcal{X}$ soit plat. Un recouvrement de
$\mathcal{X}_{flat,fppf}$ est une famille de morphismes
$\{x_i \to x\}_{i \in I}$ de $\mathcal{X}_{flat,fppf}$
qui forme un recouvrement de $\mathcal{X}_{fppf}$.
\end{enumerate}
\end{definition}

\noindent
Nous notons $\mathcal{O}_{\mathcal{X}_{lisse,\etale}}$
la restriction de $\mathcal{O}_\mathcal{X}$ au site lisse-étale,
et de même pour $\mathcal{O}_{\mathcal{X}_{flat,fppf}}$.
La relation entre le site lisse-étale et le site étale est
la suivante (nous nous en tenons surtout aux propriétés « topologiques » dans ce lemme).

\begin{lemma}
\label{lemma-lisse-etale}
Soit $\mathcal{X}$ un champ algébrique.
\begin{enumerate}
\item Le foncteur d'inclusion
$\mathcal{X}_{lisse,\etale} \to \mathcal{X}_\etale$
est pleinement fidèle, continu et cocontinu. Il s'ensuit que
\begin{enumerate}
\item il existe un morphisme de topos
$$
g :
\Sh(\mathcal{X}_{lisse,\etale})
\longrightarrow
\Sh(\mathcal{X}_\etale)
$$
où $g^{-1}$ est donné par la restriction,
\item le foncteur $g^{-1}$ admet un adjoint à gauche $g_!^{Sh}$ sur les faisceaux d'ensembles,
\item les morphismes d'adjonction $g^{-1}g_* \to \text{id}$ et
$\text{id} \to g^{-1}g_!^{Sh}$ sont des isomorphismes,
\item le foncteur $g^{-1}$ admet un adjoint à gauche $g_!$ sur les faisceaux abéliens,
\item le morphisme d'adjonction $\text{id} \to g^{-1}g_!$ est un isomorphisme, et
\item on a $g^{-1}\mathcal{O}_\mathcal{X} =
\mathcal{O}_{\mathcal{X}_{lisse,\etale}}$ ; ainsi, $g$ induit un morphisme plat
de topos annelés tel que $g^{-1} = g^*$.
\end{enumerate}
\item Le foncteur d'inclusion
$\mathcal{X}_{flat,fppf} \to \mathcal{X}_{fppf}$
est pleinement fidèle, continu et cocontinu. Il s'ensuit que
\begin{enumerate}
\item il existe un morphisme de topos
$$
g :
\Sh(\mathcal{X}_{flat,fppf})
\longrightarrow
\Sh(\mathcal{X}_{fppf})
$$
où $g^{-1}$ est donné par la restriction,
\item le foncteur $g^{-1}$ admet un adjoint à gauche $g_!^{Sh}$ sur les faisceaux d'ensembles,
\item les morphismes d'adjonction $g^{-1}g_* \to \text{id}$ et
$\text{id} \to g^{-1}g_!^{Sh}$ sont des isomorphismes,
\item le foncteur $g^{-1}$ admet un adjoint à gauche $g_!$ sur les faisceaux abéliens,
\item le morphisme d'adjonction $\text{id} \to g^{-1}g_!$ est un isomorphisme, et
\item on a $g^{-1}\mathcal{O}_\mathcal{X} =
\mathcal{O}_{\mathcal{X}_{flat,fppf}}$ ; ainsi, $g$ induit un morphisme plat
de topos annelés tel que $g^{-1} = g^*$.
\end{enumerate}
\end{enumerate}
\end{lemma}

\begin{proof}
Dans les deux cas, il est immédiat que le foncteur est pleinement fidèle,
continu et cocontinu (voir
Sites, définitions \ref{sites-definition-continuous} et
\ref{sites-definition-cocontinuous}).
Les propriétés (a), (b), (c) résultent donc de
Sites, lemmes \ref{sites-lemma-when-shriek} et
\ref{sites-lemma-back-and-forth}.
Les parties (d), (e) résultent de
Modules sur les sites, lemmes \ref{sites-modules-lemma-g-shriek-adjoint} et
\ref{sites-modules-lemma-back-and-forth}.
La partie (f) est immédiate.
\end{proof}

\begin{lemma}
\label{lemma-lisse-etale-cohomology}
Soit $\mathcal{X}$ un champ algébrique. On reprend les notations du
Lemme \ref{lemma-lisse-etale}.
\begin{enumerate}
\item Pour un faisceau abélien $\mathcal{F}$ sur $\mathcal{X}_\etale$, on a
\begin{enumerate}
\item $H^p(\mathcal{X}_\etale, \mathcal{F}) =
H^p(\mathcal{X}_{lisse,\etale}, g^{-1}\mathcal{F})$, et
\item $H^p(x, \mathcal{F}) =
H^p(\mathcal{X}_{lisse,\etale}/x, g^{-1}\mathcal{F})$
pour tout objet $x$ de $\mathcal{X}_{lisse,\etale}$.
\end{enumerate}
Il en va de même pour les faisceaux de modules.
\item Pour un faisceau abélien $\mathcal{F}$ sur $\mathcal{X}_{fppf}$, on a
\begin{enumerate}
\item $H^p(\mathcal{X}_{fppf}, \mathcal{F}) =
H^p(\mathcal{X}_{flat,fppf}, g^{-1}\mathcal{F})$, et
\item $H^p(x, \mathcal{F}) =
H^p(\mathcal{X}_{flat,fppf}/x, g^{-1}\mathcal{F})$
pour tout objet $x$ de $\mathcal{X}_{flat,fppf}$.
\end{enumerate}
Il en va de même pour les faisceaux de modules.
\end{enumerate}
\end{lemma}

\begin{proof}
La partie (1)(a) résulte de Faisceaux sur les champs, lemme
\ref{stacks-sheaves-lemma-cohomology-on-subcategory}, appliqué au foncteur
d'inclusion $\mathcal{X}_{lisse,\etale} \to \mathcal{X}_\etale$.
La partie (1)(b) résulte de la partie (1)(a). En effet, si $x$ est au-dessus du
schéma $U$, alors le site $\mathcal{X}_\etale/x$ est équivalent
à $(\Sch/U)_\etale$, et $\mathcal{X}_{lisse,\etale}$ est équivalent
à $U_{lisse,\etale}$. La partie (2) se démontre de la même manière.
\end{proof}

\begin{lemma}
\label{lemma-lisse-etale-modules}
Soit $\mathcal{X}$ un champ algébrique. On reprend les notations du
Lemme \ref{lemma-lisse-etale}.
\begin{enumerate}
\item Il existe un foncteur
$$
g_! :
\textit{Mod}(\mathcal{X}_{lisse,\etale},
\mathcal{O}_{\mathcal{X}_{lisse,\etale}})
\longrightarrow
\textit{Mod}(\mathcal{X}_\etale, \mathcal{O}_{\mathcal{X}})
$$
qui est adjoint à gauche de $g^*$. En outre, il coïncide avec le foncteur $g_!$
sur les faisceaux abéliens, et $g^*g_! = \text{id}$.
\item Il existe un foncteur
$$
g_! :
\textit{Mod}(\mathcal{X}_{flat,fppf},
\mathcal{O}_{\mathcal{X}_{flat,fppf}})
\longrightarrow
\textit{Mod}(\mathcal{X}_{fppf}, \mathcal{O}_{\mathcal{X}})
$$
qui est adjoint à gauche de $g^*$. En outre, il coïncide avec le foncteur $g_!$
sur les faisceaux abéliens, et $g^*g_! = \text{id}$.
\end{enumerate}
\end{lemma}

\begin{proof}
Dans les deux cas, l'existence du foncteur $g_!$ résulte de
Modules sur les sites, lemme \ref{sites-modules-lemma-lower-shriek-modules}.
Pour voir que $g_!$ coïncide avec le foncteur sur les faisceaux abéliens, nous
montrons que les morphismes de Modules sur les sites, équation
(\ref{sites-modules-equation-compare-on-localizations})
sont des isomorphismes.

\medskip\noindent
Cas lisse-étale. Soit $x \in \Ob(\mathcal{X}_{lisse,\etale})$
au-dessus d'un schéma $U$, avec $x : U \to \mathcal{X}$ lisse.
Considérons le foncteur pleinement fidèle induit
$$
g' :
\mathcal{X}_{lisse,\etale}/x
\longrightarrow
\mathcal{X}_\etale/x
$$
Le membre de droite s'identifie à $(\Sch/U)_\etale$, et celui de
gauche à la sous-catégorie pleine des schémas $U'/U$ tels que
le composé $U' \to U \to \mathcal{X}$ soit lisse. Ainsi,
Cohomologie étale, lemme
\ref{etale-cohomology-lemma-compare-structure-sheaves}
s'applique.

\medskip\noindent
Cas plat-fppf. Soit $x \in \Ob(\mathcal{X}_{flat,fppf})$
au-dessus d'un schéma $U$, avec $x : U \to \mathcal{X}$ plat.
Considérons le foncteur pleinement fidèle induit
$$
g' :
\mathcal{X}_{flat,fppf}/x
\longrightarrow
\mathcal{X}_{fppf}/x
$$
Le membre de droite s'identifie à $(\Sch/U)_{fppf}$, et celui de
gauche à la sous-catégorie pleine des schémas $U'/U$ tels que
le composé $U' \to U \to \mathcal{X}$ soit plat. Ainsi,
Cohomologie étale, lemme
\ref{etale-cohomology-lemma-compare-structure-sheaves}
s'applique.

\medskip\noindent
Dans les deux cas, l'égalité $g^*g_! = \text{id}$ résulte de
$g^* = g^{-1}$ et de
l'égalité pour les faisceaux abéliens du Lemme \ref{lemma-lisse-etale}.
\end{proof}

\begin{lemma}
\label{lemma-lisse-etale-structure-sheaf}
Soit $\mathcal{X}$ un champ algébrique. On reprend les notations des
lemmes \ref{lemma-lisse-etale} et \ref{lemma-lisse-etale-modules}.
\begin{enumerate}
\item On a $g_!\mathcal{O}_{\mathcal{X}_{lisse,\etale}} =
\mathcal{O}_\mathcal{X}$.
\item On a $g_!\mathcal{O}_{\mathcal{X}_{flat, fppf}} =
\mathcal{O}_\mathcal{X}$.
\end{enumerate}
\end{lemma}

\begin{proof}
Dans cette démonstration, nous posons
$\mathcal{C} = \mathcal{X}_\etale$
(resp.\ $\mathcal{C} = \mathcal{X}_{fppf}$)
et nous notons
$\mathcal{C}' = \mathcal{X}_{lisse,\etale}$
(resp.\ $\mathcal{C}' = \mathcal{X}_{flat, fppf}$).
Alors $\mathcal{C}'$ est une sous-catégorie pleine de $\mathcal{C}$.
Dans cette démonstration, nous regardons les objets $V$ de $\mathcal{C}$
comme des schémas au-dessus de $\mathcal{X}$ et les objets $U$ de $\mathcal{C}'$
comme des schémas lisses (resp.\ plats) au-dessus de $\mathcal{X}$.
Enfin, nous posons $\mathcal{O} = \mathcal{O}_\mathcal{X}$
et $\mathcal{O}' = \mathcal{O}_{\mathcal{X}_{lisse,\etale}}$
(resp.\ $\mathcal{O}' = \mathcal{O}_{\mathcal{X}_{flat,fppf}}$).
Avec ces notations, on a $\mathcal{O}(V) = \Gamma(V, \mathcal{O}_V)$
et $\mathcal{O}'(U) = \Gamma(U, \mathcal{O}_U)$.
Considérons l'homomorphisme de $\mathcal{O}$-modules
$g_!\mathcal{O}' \to \mathcal{O}$
adjoint à l'identification $\mathcal{O}' = g^{-1}\mathcal{O}$.

\medskip\noindent
Rappelons que $g_!\mathcal{O}'$ est le faisceau associé au préfaisceau
$g_{p!}\mathcal{O}'$ défini par la règle
$$
V \longmapsto \colim_{V \to U} \mathcal{O}'(U)
$$
où la limite inductive est prise dans la catégorie des groupes abéliens
(Modules sur les sites, définition \ref{sites-modules-definition-g-shriek}).
Nous utiliserons fréquemment ci-dessous le fait que, si
$$
V \to U \to U'
$$
sont des morphismes et si $f' \in \mathcal{O}'(U')$ a pour restriction
$f \in \mathcal{O}'(U)$, alors $(V \to U, f)$ et $(V \to U', f')$
définissent le même élément de la limite inductive. De plus,
$g_!\mathcal{O}' \to \mathcal{O}$ envoie simplement l'élément
$(V \to U, f)$ sur l'image réciproque de $f$ dans $V$.

\medskip\noindent
Démontrons que $g_!\mathcal{O}' \to \mathcal{O}$ est surjectif.
Soit $h \in \mathcal{O}(V)$ pour un objet $V$ de $\mathcal{C}$.
Il suffit de montrer que $h$ appartient localement à l'image. Choisissons un objet
$U$ de $\mathcal{C}'$ correspondant à un morphisme lisse surjectif
$U \to \mathcal{X}$. Puisque $U \times_\mathcal{X} V \to V$ est lisse et
surjectif, après avoir remplacé $V$ par les membres d'un recouvrement étale de $V$,
nous pouvons supposer qu'il existe un morphisme $V \to U$ ; voir
Topologies sur les espaces algébriques, lemme
\ref{spaces-topologies-lemma-etale-dominates-smooth}. À l'aide de $h$, nous obtenons
un morphisme $V \to U \times \mathbf{A}^1$ tel que, si l'on écrit
$\mathbf{A}^1 = \Spec(\mathbf{Z}[t])$, l'élément
$t \in \mathcal{O}(U \times \mathbf{A}^1)$ ait pour image réciproque $h$.
Comme $U \times \mathbf{A}^1$ est un objet de $\mathcal{C}'$,
$(V \to U \times \mathbf{A}^1, t)$ est un élément de la limite inductive
ci-dessus qui s'envoie sur $h \in \mathcal{O}(V)$,
comme voulu.

\medskip\noindent
Supposons que $s \in g_!\mathcal{O}'(V)$ soit une section
dont l'image dans $\mathcal{O}(V)$ est nulle. Pour achever la démonstration,
il reste à montrer que $s$ est nul. Après avoir remplacé $V$ par les membres
d'un recouvrement, nous pouvons supposer que $s$ est un élément de la limite inductive
$$
\colim_{V \to U} \mathcal{O}'(U)
$$
Écrivons $s = \sum (\varphi_i, s_i)$ comme une somme finie, où
$\varphi_i : V \to U_i$, les $U_i$ sont lisses (resp.\ plats) sur $\mathcal{X}$ et
$s_i \in \Gamma(U_i, \mathcal{O}_{U_i})$. Choisissons un schéma $W$ étale
et surjectif sur l'espace algébrique
$U = U_1 \times_\mathcal{X} \ldots \times_\mathcal{X} U_n$.
Notons que $W$ est encore lisse (resp.\ plat) sur $\mathcal{X}$, c'est-à-dire
qu'il définit un objet de $\mathcal{C}'$. Le produit fibré
$$
V' = V \times_{(\varphi_1, \ldots, \varphi_n), U} W
$$
est étale et surjectif sur $V$ ; il suffit donc de montrer que l'image de $s$
est nulle dans $g_!\mathcal{O}'(V')$. La restriction
$\sum (\varphi_i, s_i)|_{V'}$ correspond à la somme des images réciproques
des fonctions $s_i$ sur $W$. Autrement dit, nous sommes ramenés au cas
d'un couple $(\varphi, s)$, où $\varphi : V \to U$ est un morphisme avec $U$
dans $\mathcal{C}'$ et où la restriction de $s \in \mathcal{O}'(U)$ dans
$\mathcal{O}(V)$ est nulle. D'après le diagramme commutatif
$$
\xymatrix{
V \ar[rr]_-{(\varphi, 0)} \ar[rrd]_\varphi & & U \times \mathbf{A}^1 \\
& & U \ar[u]_{(\text{id}, 0)}
}
$$
on voit que
$((\varphi, 0) : V \to U \times \mathbf{A}^1, \text{pr}_2^*x)$
représente zéro dans la limite inductive ci-dessus. Nous pouvons donc
remplacer $U$ par $U \times \mathbf{A}^1$, $\varphi$ par $(\varphi, 0)$
et $s$ par $\text{pr}_1^*s + \text{pr}_2^*x$. Nous pouvons ainsi supposer que
le lieu $Z : s = 0$ dans $U$ défini par $s$ est lisse (resp.\ plat)
sur $\mathcal{X}$. On voit alors que $(V \to Z, 0)$ et $(\varphi, s)$
ont la même valeur dans la limite inductive, c'est-à-dire que l'élément $s$
est nul, comme voulu.
\end{proof}

\noindent
Les sites lisse-étale et plat-fppf permettent de caractériser
les modules parasites comme suit.

\begin{lemma}
\label{lemma-parasitic-in-terms-flat-fppf}
Soit $\mathcal{X}$ un champ algébrique.
\begin{enumerate}
\item Soit $\mathcal{F}$ un $\mathcal{O}_\mathcal{X}$-module
sur $\mathcal{X}_\etale$ ayant la propriété de changement de base plat.
Les conditions suivantes sont équivalentes :
\begin{enumerate}
\item $\mathcal{F}$ est parasite ;
\item $g^*\mathcal{F} = 0$, où
$g : \Sh(\mathcal{X}_{lisse,\etale}) \to
\Sh(\mathcal{X}_\etale)$ est défini comme dans le lemme \ref{lemma-lisse-etale}.
\end{enumerate}
\item Soit $\mathcal{F}$ un $\mathcal{O}_\mathcal{X}$-module sur
$\mathcal{X}_{fppf}$. Les conditions suivantes sont équivalentes :
\begin{enumerate}
\item $\mathcal{F}$ est parasite ;
\item $g^*\mathcal{F} = 0$, où
$g :  \Sh(\mathcal{X}_{flat,fppf}) \to \Sh(\mathcal{X}_{fppf})$
est défini comme dans le lemme \ref{lemma-lisse-etale}.
\end{enumerate}
\end{enumerate}
\end{lemma}

\begin{proof}
La partie (2) résulte immédiatement des définitions (c'est l'un des avantages
du site plat-fppf sur le site lisse-étale). L'implication
(1)(a) $\Rightarrow$ (1)(b) est elle aussi immédiate. Pour établir (1)(b)
$\Rightarrow$ (1)(a), soient $U$ un schéma et $x : U \to \mathcal{X}$
un morphisme lisse surjectif. Alors $x$ est un objet du
site lisse-étale de $\mathcal{X}$. La condition (1)(b)
implique donc que $\mathcal{F}|_{U_\etale} = 0$. Soit $V \to \mathcal{X}$
un morphisme plat, où $V$ est un schéma. Posons $W = U \times_\mathcal{X} V$
et considérons le diagramme
$$
\xymatrix{
W \ar[d]_p \ar[r]_q & V \ar[d] \\
U \ar[r] & \mathcal{X}
}
$$
La projection $p : W \to U$ est plate, tandis que la projection
$q : W \to V$ est lisse et surjective. Il en résulte que $q_{small}^*$
est un foncteur fidèle sur les modules quasi-cohérents. Par hypothèse,
$\mathcal{F}$ possède la propriété de changement de base plat, d'où un isomorphisme
$p_{small}^*\mathcal{F}|_{U_\etale} \cong
q_{small}^*\mathcal{F}|_{V_\etale}$. Ainsi, si $\mathcal{F}$
appartient au noyau de $g^*$, alors $\mathcal{F}|_{V_\etale} = 0$,
comme voulu.
\end{proof}






\section{Fonctorialité des sites lisse-étale et plat-fppf}
\label{section-lisse-etale-functorial}

\noindent
Le site lisse-étale est fonctoriel pour les morphismes lisses de champs
algébriques, et le site plat-fppf l'est pour les morphismes plats de champs
algébriques. Avertissons le lecteur que les topos lisse-étale et plat-fppf ne
sont pas fonctoriels par rapport à tous les morphismes de champs algébriques ; voir
Exemples, section \ref{examples-section-lisse-etale-not-functorial}.

\begin{lemma}
\label{lemma-lisse-etale-functorial}
Soit $f : \mathcal{X} \to \mathcal{Y}$ un morphisme de champs algébriques.
\begin{enumerate}
\item Si $f$ est lisse, alors $f$ se restreint en un foncteur continu et
cocontinu
$\mathcal{X}_{lisse,\etale} \to \mathcal{Y}_{lisse,\etale}$
qui donne un morphisme de topos annelés s'insérant dans le diagramme
commutatif suivant
$$
\xymatrix{
\Sh(\mathcal{X}_{lisse,\etale}) \ar[r]_{g'} \ar[d]_{f'} &
\Sh(\mathcal{X}_\etale) \ar[d]^f \\
\Sh(\mathcal{Y}_{lisse,\etale}) \ar[r]^g &
\Sh(\mathcal{Y}_\etale)
}
$$
On a $f'_*(g')^{-1} = g^{-1}f_*$ et $g'_!(f')^{-1} = f^{-1}g_!$.
\item Si $f$ est plat, alors $f$ se restreint en un foncteur continu et
cocontinu
$\mathcal{X}_{flat,fppf} \to \mathcal{Y}_{flat,fppf}$
qui donne un morphisme de topos annelés s'insérant dans le diagramme
commutatif suivant
$$
\xymatrix{
\Sh(\mathcal{X}_{flat,fppf}) \ar[r]_{g'} \ar[d]_{f'} &
\Sh(\mathcal{X}_{fppf}) \ar[d]^f \\
\Sh(\mathcal{Y}_{flat,fppf}) \ar[r]^g &
\Sh(\mathcal{Y}_{fppf})
}
$$
On a $f'_*(g')^{-1} = g^{-1}f_*$ et $g'_!(f')^{-1} = f^{-1}g_!$.
\end{enumerate}
\end{lemma}

\begin{proof}
La première assertion résulte du fait que, si $x \in \Ob(\mathcal{X})$
est au-dessus d'un schéma $U$, si $x : U \to \mathcal{X}$ est lisse
(resp.\ plat) et si $f$ est lisse (resp.\ plat), alors
$f(x) : U \to \mathcal{Y}$ est lisse (resp.\ plat) ; voir
Morphismes de champs, lemmes \ref{stacks-morphisms-lemma-composition-smooth} et
\ref{stacks-morphisms-lemma-composition-flat}. Le foncteur induit
$\mathcal{X}_{lisse,\etale} \to \mathcal{Y}_{lisse,\etale}$
(resp.\ $\mathcal{X}_{flat,fppf} \to \mathcal{Y}_{flat,fppf}$) est
continu et cocontinu par notre définition des recouvrements dans ces
catégories. Enfin, la commutativité du diagramme résulte du fait que les
morphismes horizontaux sont donnés par les foncteurs d'inclusion (voir le
lemme \ref{lemma-lisse-etale}) et du
lemme \ref{sites-lemma-composition-cocontinuous} de Sites.

\medskip\noindent
Pour montrer que $f'_*(g')^{-1} = g^{-1}f_*$, soit $\mathcal{F}$ un faisceau
sur $\mathcal{X}_\etale$ (resp.\ $\mathcal{X}_{fppf}$).
Il existe un morphisme canonique d'image inverse
$$
g^{-1}f_*\mathcal{F} \longrightarrow f'_*(g')^{-1}\mathcal{F}
$$
voir Sites, section \ref{sites-section-pullback}.
Nous affirmons que ce morphisme est un isomorphisme.
Pour le montrer, prenons un objet $y$ de $\mathcal{Y}_{lisse,\etale}$
(resp.\ $\mathcal{Y}_{flat,fppf}$). Supposons que $y$ soit au-dessus du schéma
$V$ et que $y : V \to \mathcal{Y}$ soit lisse (resp.\ plat). Comme
$g^{-1}$ est le foncteur de restriction, on obtient
$$
\left(g^{-1}f_*\mathcal{F}\right)(y) =
\Gamma(V \times_{y, \mathcal{Y}} \mathcal{X},\ \text{pr}^{-1}\mathcal{F})
$$
d'après Faisceaux sur les champs, équation
(\ref{stacks-sheaves-equation-pushforward}). Soit
$(V \times_{y, \mathcal{Y}} \mathcal{X})' \subset
V \times_{y, \mathcal{Y}} \mathcal{X}$
la sous-catégorie pleine formée des objets
$z : W \to V \times_{y, \mathcal{Y}} \mathcal{X}$ tels que le morphisme
induit $W \to \mathcal{X}$ soit lisse (resp.\ plat). Notons
$$
\text{pr}' :
(V \times_{y, \mathcal{Y}} \mathcal{X})'
\longrightarrow
\mathcal{X}_{lisse,\etale}
\ (\text{resp. }\mathcal{X}_{flat,fppf})
$$
la restriction du foncteur $\text{pr}$ utilisé dans la formule ci-dessus.
Exactement le même argument que celui qui démontre
Faisceaux sur les champs, équation (\ref{stacks-sheaves-equation-pushforward}),
montre que, pour tout faisceau $\mathcal{H}$ sur $\mathcal{X}_{lisse,\etale}$
(resp.\ $\mathcal{X}_{flat,fppf}$), on a
\begin{equation}
\label{equation-pushforward-lisse-etale}
f'_*\mathcal{H}(y) =
\Gamma((V \times_{y, \mathcal{Y}} \mathcal{X})',
\ (\text{pr}')^{-1}\mathcal{H})
\end{equation}
Comme $(g')^{-1}$ est le foncteur de restriction, on voit que
$$
\left(f'_*(g')^{-1}\mathcal{F}\right)(y) =
\Gamma((V \times_{y, \mathcal{Y}} \mathcal{X})',
\ \text{pr}^{-1}\mathcal{F}|_{(V \times_{y, \mathcal{Y}} \mathcal{X})'})
$$
D'après le
lemme \ref{stacks-sheaves-lemma-cohomology-on-subcategory} de Faisceaux sur les champs,
on voit que
$$
\Gamma((V \times_{y, \mathcal{Y}} \mathcal{X})',
\ \text{pr}^{-1}\mathcal{F}|_{(V \times_{y, \mathcal{Y}} \mathcal{X})'})
=
\Gamma(V \times_{y, \mathcal{Y}} \mathcal{X},\ \text{pr}^{-1}\mathcal{F})
$$
sont égaux, comme voulu ; bien que nous omettions la vérification des hypothèses
du lemme, notons que le fait que $V \to \mathcal{Y}$ soit lisse
(resp.\ plat) sert à vérifier la seconde condition.

\medskip\noindent
Enfin, l'égalité $g'_!(f')^{-1} = f^{-1}g_!$ découle formellement de
l'égalité $f'_*(g')^{-1} = g^{-1}f_*$ par les adjonctions entre
$f^{-1}$ et $f_*$, entre $g_!$ et $g^{-1}$, et par leurs versions
« primées ».
\end{proof}

\begin{lemma}
\label{lemma-lisse-etale-functorial-pushforward}
Sous les hypothèses et avec les notations du lemme
\ref{lemma-lisse-etale-functorial}, soit $\mathcal{H}$ un faisceau abélien sur
$\mathcal{X}_{lisse,\etale}$ (resp.\ $\mathcal{X}_{flat,fppf}$). Alors
\begin{equation}
\label{equation-higher-direct-image-lisse-etale}
R^pf'_*\mathcal{H} =
\text{faisceau associé au préfaisceau }y \longmapsto
H^p((V \times_{y, \mathcal{Y}} \mathcal{X})', (\text{pr}')^{-1}\mathcal{H})
\end{equation}
Ici, $y$ est un objet de $\mathcal{Y}_{lisse,\etale}$
(resp.\ $\mathcal{Y}_{flat,fppf}$) au-dessus du schéma $V$, et les notations
$(V \times_{y, \mathcal{Y}} \mathcal{X})'$ et $\text{pr}'$ sont
expliquées dans la démonstration.
\end{lemma}

\begin{proof}
Comme dans la démonstration du lemme \ref{lemma-lisse-etale-functorial},
soit $(V \times_{y, \mathcal{Y}} \mathcal{X})' \subset
V \times_{y, \mathcal{Y}} \mathcal{X}$ la sous-catégorie pleine
formée des objets $(x, \varphi)$, où $x$ est un objet de
$\mathcal{X}_{lisse,\etale}$ (resp.\ $\mathcal{X}_{flat,fppf}$)
et $\varphi : f(x) \to y$ est un morphisme dans $\mathcal{Y}$.
D'après l'équation (\ref{equation-pushforward-lisse-etale}),
on a
$$
f'_*\mathcal{H}(y) =
\Gamma((V \times_{y, \mathcal{Y}} \mathcal{X})',
\ (\text{pr}')^{-1}\mathcal{H})
$$
où $\text{pr}'$ est la projection. Pour un objet $(x, \varphi)$
de $(V \times_{y, \mathcal{Y}} \mathcal{X})'$,
on peut considérer $\varphi$ comme une section de $(f')^{-1}h_y$ au-dessus de $x$.
Ainsi, $(V \times_\mathcal{Y} \mathcal{X})'$ est la localisation
du site $\mathcal{X}_{lisse,\etale}$
(resp. $\mathcal{X}_{flat,fppf}$) en le faisceau d'ensembles $(f')^{-1}h_y$ ; voir
Sites, lemme \ref{sites-lemma-localize-topos-site}. Le morphisme
$$
\text{pr}' : (V \times_{y, \mathcal{Y}} \mathcal{X})'
\to \mathcal{X}_{lisse,\etale}
\ (\text{resp. }
\text{pr}' : (V \times_{y, \mathcal{Y}} \mathcal{X})'
\to \mathcal{X}_{flat,fppf})
$$
est le morphisme de localisation.
En particulier, le foncteur image inverse $(\text{pr}')^{-1}$ préserve
les faisceaux abéliens injectifs ; voir
Cohomologie sur les sites, lemme
\ref{sites-cohomology-lemma-cohomology-on-sheaf-sets}.

\medskip\noindent
Choisissons une résolution injective $\mathcal{H} \to \mathcal{I}^\bullet$
sur $\mathcal{X}_{lisse,\etale}$ (resp.\ $\mathcal{X}_{flat,fppf}$).
La formule de l'image directe montre que $R^if'_*\mathcal{H}$ est le
faisceau associé au préfaisceau qui, à $y$, associe la cohomologie
du complexe
$$
\begin{matrix}
\Gamma\Big((V \times_{y, \mathcal{Y}} \mathcal{X})',
(\text{pr}')^{-1}\mathcal{I}^{i - 1}\Big) \\
\downarrow \\
\Gamma\Big((V \times_{y, \mathcal{Y}} \mathcal{X})',
(\text{pr}')^{-1}\mathcal{I}^i\Big) \\
\downarrow \\
\Gamma\Big((V \times_{y, \mathcal{Y}} \mathcal{X})',
(\text{pr}')^{-1}\mathcal{I}^{i + 1}\Big)
\end{matrix}
$$
Comme $(\text{pr}')^{-1}$ est exact et préserve les objets injectifs, le
complexe $(\text{pr}')^{-1}\mathcal{I}^\bullet$ est une
résolution injective de $(\text{pr}')^{-1}\mathcal{H}$,
ce qui achève la démonstration.
\end{proof}

\begin{lemma}
\label{lemma-lisse-etale-functorial-cohomology}
Sous les hypothèses et avec les notations du lemme
\ref{lemma-lisse-etale-functorial}, le morphisme canonique de changement de base
$$
g^{-1}Rf_*\mathcal{F} \longrightarrow Rf'_*(g')^{-1}\mathcal{F}
$$
est un isomorphisme pour tout faisceau abélien $\mathcal{F}$
sur $\mathcal{X}_\etale$ (resp.\ $\mathcal{X}_{fppf}$).
\end{lemma}

\begin{proof}
En comparant les formules pour $g^{-1}R^pf_*\mathcal{F}$
et $R^pf'_*(g')^{-1}\mathcal{F}$ données par le
lemme \ref{stacks-sheaves-lemma-pushforward-restriction} de Faisceaux sur les champs
et le lemme \ref{lemma-lisse-etale-functorial-pushforward},
on voit qu'il suffit de montrer
$$
H^p((V \times_{y, \mathcal{Y}} \mathcal{X})',
\ \text{pr}^{-1}\mathcal{F}|_{(V \times_{y, \mathcal{Y}} \mathcal{X})'})
=
H^p_\tau(V \times_{y, \mathcal{Y}} \mathcal{X},\ \text{pr}^{-1}\mathcal{F})
$$
où $\tau = \etale$ (resp.\ $\tau = fppf$). Ici, $y$ est un objet
de $\mathcal{Y}$ au-dessus d'un schéma $V$ tel que le morphisme
$y : V \to \mathcal{Y}$ soit lisse (resp.\ plat).
Cette égalité résulte du
lemme \ref{stacks-sheaves-lemma-cohomology-on-subcategory} de Faisceaux sur les champs.
Bien que nous omettions la vérification des hypothèses
du lemme, notons que le fait que $V \to \mathcal{Y}$ soit lisse
(resp.\ plat) sert à vérifier la seconde condition.
\end{proof}






\section{Modules quasi-cohérents et sites lisse-étale et plat-fppf}
\label{section-quasi-coherent-modules-II}

\noindent
Dans cette section, nous expliquons comment décrire les modules quasi-cohérents
sur un champ algébrique au moyen de son site lisse-étale ou plat-fppf.

\begin{lemma}
\label{lemma-check-qc-on-etale-covering}
Soit $\mathcal{X}$ un champ algébrique.
\begin{enumerate}
\item Soit $f_j : \mathcal{X}_j \to \mathcal{X}$ une famille de morphismes
lisses de champs algébriques telle que
$|\mathcal{X}| =\bigcup |f_j|(|\mathcal{X}_j|)$.
Soit $\mathcal{F}$ un faisceau de $\mathcal{O}_\mathcal{X}$-modules
sur $\mathcal{X}_\etale$. Si chaque $f_j^{-1}\mathcal{F}$ est
quasi-cohérent, alors $\mathcal{F}$ l'est aussi.
\item Soit $f_j : \mathcal{X}_j \to \mathcal{X}$ une famille de morphismes
plats et localement de présentation finie de champs algébriques telle que
$|\mathcal{X}| =\bigcup |f_j|(|\mathcal{X}_j|)$.
Soit $\mathcal{F}$ un faisceau de $\mathcal{O}_\mathcal{X}$-modules
sur $\mathcal{X}_{fppf}$. Si chaque $f_j^{-1}\mathcal{F}$ est
quasi-cohérent, alors $\mathcal{F}$ l'est aussi.
\end{enumerate}
\end{lemma}

\begin{proof}
Preuve de (1). Nous pouvons remplacer chacun des champs algébriques $\mathcal{X}_j$
par un schéma $U_j$ (puisque tout champ algébrique admet un recouvrement lisse par
un schéma et que les composés de morphismes lisses sont lisses ; voir
Morphismes de champs, lemme \ref{stacks-morphisms-lemma-composition-smooth}).
L'image inverse de $\mathcal{F}$ sur $(\Sch/U_j)_\etale$ est encore
quasi-cohérente ; voir
Modules sur les sites, lemme \ref{sites-modules-lemma-local-pullback}.
Alors $f = \coprod f_j : U = \coprod U_j \to \mathcal{X}$ est un morphisme lisse
surjectif. Soit $x : V \to \mathcal{X}$ un objet de $\mathcal{X}$. D'après
Faisceaux sur les champs, lemme
\ref{stacks-sheaves-lemma-surjective-flat-locally-finite-presentation}
il existe un recouvrement étale $\{x_i \to x\}_{i \in I}$ tel que chaque
$x_i$ se relève en un objet $u_i$ de $(\Sch/U)_\etale$.
Cela signifie simplement que $x_i$ est au-dessus d'un schéma $V_i$, que
$\{V_i \to V\}$ est un recouvrement étale et que $x_i$ provient
d'un morphisme $u_i : V_i \to U$. Alors
$x_i^*\mathcal{F} = u_i^*f^*\mathcal{F}$ est quasi-cohérent.
Il s'ensuit que $x^*\mathcal{F}$ sur $(\Sch/V)_\etale$
est quasi-cohérent, par exemple d'après
Modules sur les sites, lemme \ref{sites-modules-lemma-local-final-object}.
D'après Faisceaux sur les champs, lemme
\ref{stacks-sheaves-lemma-characterize-quasi-coherent-bis}
on voit que $x^*\mathcal{F}$ est un faisceau fppf et, puisque $x$
était arbitraire, que $\mathcal{F}$ est un faisceau pour la
topologie fppf. En appliquant Faisceaux sur les champs, lemme
\ref{stacks-sheaves-lemma-characterize-quasi-coherent}
on voit que $\mathcal{F}$ est quasi-cohérent.

\medskip\noindent
Preuve de (2). La démonstration emploie exactement le même argument, que nous
détaillons ici. Nous pouvons remplacer chacun des champs algébriques $\mathcal{X}_j$
par un schéma $U_j$ (puisque tout champ algébrique admet un recouvrement lisse par
un schéma et que les morphismes plats et localement de présentation finie sont
stables par composition ; voir Morphismes de champs, lemmes
\ref{stacks-morphisms-lemma-composition-flat} et
\ref{stacks-morphisms-lemma-composition-finite-presentation}).
L'image inverse de $\mathcal{F}$ sur $(\Sch/U_j)_\etale$ est encore
localement quasi-cohérente ; voir
Faisceaux sur les champs, lemme \ref{stacks-sheaves-lemma-pullback-quasi-coherent}.
Alors $f = \coprod f_j : U = \coprod U_j \to \mathcal{X}$ est un morphisme
surjectif, plat et localement de présentation finie. Soit
$x : V \to \mathcal{X}$ un objet de $\mathcal{X}$. D'après
Faisceaux sur les champs, lemme
\ref{stacks-sheaves-lemma-surjective-flat-locally-finite-presentation}
il existe un recouvrement fppf $\{x_i \to x\}_{i \in I}$ tel que chaque
$x_i$ se relève en un objet $u_i$ de $(\Sch/U)_\etale$.
Cela signifie simplement que $x_i$ est au-dessus d'un schéma $V_i$, que
$\{V_i \to V\}$ est un recouvrement fppf et que $x_i$ provient
d'un morphisme $u_i : V_i \to U$. Alors
$x_i^*\mathcal{F} = u_i^*f^*\mathcal{F}$ est quasi-cohérent.
Il s'ensuit que $x^*\mathcal{F}$ sur $(\Sch/V)_\etale$
est quasi-cohérent, par exemple d'après
Modules sur les sites, lemme \ref{sites-modules-lemma-local-final-object}.
D'après Faisceaux sur les champs, lemme
\ref{stacks-sheaves-lemma-characterize-quasi-coherent}
on voit que $\mathcal{F}$ est quasi-cohérent.
\end{proof}

\noindent
Rappelons que nous avons défini la notion de module quasi-cohérent sur
un topos annelé quelconque dans
Modules sur les sites, section \ref{sites-modules-section-local}.

\begin{lemma}
\label{lemma-shriek-quasi-coherent}
Soit $\mathcal{X}$ un champ algébrique. Les notations sont celles du
lemme \ref{lemma-lisse-etale}.
\begin{enumerate}
\item Soit $\mathcal{H}$ un
$\mathcal{O}_{\mathcal{X}_{lisse,\etale}}$-module quasi-cohérent
sur le site lisse-étale de $\mathcal{X}$. Alors $g_!\mathcal{H}$ est un
module quasi-cohérent sur $\mathcal{X}$.
\item Soit $\mathcal{H}$ un
$\mathcal{O}_{\mathcal{X}_{flat,fppf}}$-module quasi-cohérent
sur le site plat-fppf de $\mathcal{X}$. Alors $g_!\mathcal{H}$ est un
module quasi-cohérent sur $\mathcal{X}$.
\end{enumerate}
\end{lemma}

\begin{proof}
Choisissons un schéma $U$ et un morphisme lisse surjectif $x : U \to \mathcal{X}$.
D'après
Modules sur les sites, définition \ref{sites-modules-definition-site-local},
il existe un recouvrement étale (resp.\ fppf)
$\{U_i \to U\}_{i \in I}$ tel que chaque image inverse $f_i^{-1}\mathcal{H}$
admette une présentation globale (voir
Modules sur les sites, définition \ref{sites-modules-definition-global}).
Ici, $f_i : U_i \to \mathcal{X}$ est le composé
$U_i \to U \to \mathcal{X}$, qui est un morphisme de champs algébriques.
(Rappelons que l'image inverse « est » la restriction à $\mathcal{X}/f_i$ ; voir
Faisceaux sur les champs, définition \ref{stacks-sheaves-definition-pullback}
et la discussion qui suit.) Comme chaque $f_i$ est lisse (resp.\ plat), le
lemme \ref{lemma-lisse-etale-functorial}
donne $f_i^{-1}g_!\mathcal{H} = g_{i, !}(f'_i)^{-1}\mathcal{H}$.
Le lemme \ref{lemma-check-qc-on-etale-covering}
ramène l'énoncé au cas où $\mathcal{H}$
admet une présentation globale. Écrivons
$$
\bigoplus\nolimits_{j \in J} \mathcal{O} \longrightarrow
\bigoplus\nolimits_{i \in I} \mathcal{O} \longrightarrow
\mathcal{H} \longrightarrow 0
$$
de $\mathcal{O}$-modules, où
$\mathcal{O} = \mathcal{O}_{\mathcal{X}_{lisse,\etale}}$
(resp.\ $\mathcal{O} = \mathcal{O}_{\mathcal{X}_{flat,fppf}}$).
Comme $g_!$ commute aux limites inductives arbitraires (en tant que foncteur
adjoint à gauche ; voir le lemme \ref{lemma-lisse-etale-modules} et
Catégories, lemme \ref{categories-lemma-adjoint-exact}),
on en déduit l'existence d'une suite exacte
$$
\bigoplus\nolimits_{j \in J} g_!\mathcal{O} \longrightarrow
\bigoplus\nolimits_{i \in I} g_!\mathcal{O} \longrightarrow
g_!\mathcal{H} \longrightarrow 0
$$
Le lemme \ref{lemma-lisse-etale-structure-sheaf}
montre que $g_!\mathcal{O} = \mathcal{O}_\mathcal{X}$.
Dans le cas (2), la démonstration est terminée. Dans le cas (1), on applique
Faisceaux sur les champs, lemme
\ref{stacks-sheaves-lemma-characterize-quasi-coherent-bis}
pour conclure.
\end{proof}

\begin{lemma}
\label{lemma-quasi-coherent}
Soit $\mathcal{X}$ un champ algébrique.
\begin{enumerate}
\item Pour le site lisse-étale, avec $g$ comme dans le
lemme \ref{lemma-lisse-etale}, on a :
\begin{enumerate}
\item les foncteurs $g^{-1}$ et $g_!$ définissent des foncteurs quasi-inverses
$$
\xymatrix{
\QCoh(\mathcal{O}_\mathcal{X}) \ar@<1ex>[r]^-{g^{-1}} &
\QCoh(\mathcal{X}_{lisse,\etale},
\mathcal{O}_{\mathcal{X}_{lisse,\etale}}) \ar@<1ex>[l]^-{g_!}
}
$$
\item si $\mathcal{F}$ appartient à $\textit{LQCoh}^{fbc}(\mathcal{O}_\mathcal{X})$,
alors $g^{-1}\mathcal{F}$ appartient à
$\QCoh(\mathcal{O}_{\mathcal{X}_{lisse,\etale}})$ ;
\item $Q(\mathcal{F}) = g_!g^{-1}\mathcal{F}$, où $Q$ est celui du
lemme \ref{lemma-adjoint}.
\end{enumerate}
\item Pour le site plat-fppf, avec $g$ comme dans le
lemme \ref{lemma-lisse-etale}, on a :
\begin{enumerate}
\item les foncteurs $g^{-1}$ et $g_!$ définissent des foncteurs quasi-inverses
$$
\xymatrix{
\QCoh(\mathcal{O}_\mathcal{X}) \ar@<1ex>[r]^-{g^{-1}} &
\QCoh(\mathcal{X}_{flat,fppf},
\mathcal{O}_{\mathcal{X}_{flat,fppf}}) \ar@<1ex>[l]^-{g_!}
}
$$
\item si $\mathcal{F}$ appartient à $\textit{LQCoh}^{fbc}(\mathcal{O}_\mathcal{X})$,
alors $g^{-1}\mathcal{F}$ appartient à
$\QCoh(\mathcal{O}_{\mathcal{X}_{flat,fppf}})$
et
\item $Q(\mathcal{F}) = g_!g^{-1}\mathcal{F}$, où $Q$ est celui du
lemme \ref{lemma-adjoint}.
\end{enumerate}
\end{enumerate}
\end{lemma}

\begin{proof}
L'image inverse par tout morphisme de topos annelés préserve les catégories de
modules quasi-cohérents ; voir
Modules sur les sites, lemme \ref{sites-modules-lemma-local-pullback}.
Par conséquent, $g^{-1}$ préserve les catégories de modules quasi-cohérents ;
nous utilisons ici l'égalité
$\QCoh(\mathcal{O}_\mathcal{X}) =
\QCoh(\mathcal{X}_\etale, \mathcal{O}_\mathcal{X})$
fournie par Faisceaux sur les champs, lemme
\ref{stacks-sheaves-lemma-characterize-quasi-coherent-bis}.
Il en va de même pour $g_!$ d'après le
lemme \ref{lemma-shriek-quasi-coherent}.
Le lemme \ref{lemma-lisse-etale} montre que
$\mathcal{H} \to g^{-1}g_!\mathcal{H}$ est un isomorphisme.
Réciproquement, si $\mathcal{F}$ appartient à $\QCoh(\mathcal{O}_\mathcal{X})$,
alors le morphisme $g_!g^{-1}\mathcal{F} \to \mathcal{F}$ est un morphisme de
modules quasi-cohérents sur $\mathcal{X}$ dont la restriction à tout schéma lisse sur
$\mathcal{X}$ est un isomorphisme. La discussion de
Faisceaux sur les champs, sections
\ref{stacks-sheaves-section-quasi-coherent-presentation} et
\ref{stacks-sheaves-section-quasi-coherent-algebraic-stacks}
(par comparaison avec les modules quasi-cohérents sur des présentations),
montre donc qu'il s'agit d'un isomorphisme. Cela démontre (1)(a) et (2)(a).

\medskip\noindent
Soit $\mathcal{F}$ un objet de
$\textit{LQCoh}^{fbc}(\mathcal{O}_\mathcal{X})$. D'après le
lemme \ref{lemma-adjoint-kernel-parasitic},
le noyau et le conoyau du morphisme
$Q(\mathcal{F}) \to \mathcal{F}$ sont parasites. Par conséquent, d'après le
lemme \ref{lemma-parasitic-in-terms-flat-fppf}
et puisque $g^* = g^{-1}$ est exact, on conclut que
$g^*Q(\mathcal{F}) \to g^*\mathcal{F}$ est un isomorphisme. Ainsi,
$g^*\mathcal{F}$ est quasi-cohérent. Cela démontre (1)(b) et (2)(b).
Enfin, (1)(c) et (2)(c) résultent du fait que
$g_!g^*Q(\mathcal{F}) \to Q(\mathcal{F})$ est un isomorphisme d'après
les arguments qui précèdent.
\end{proof}

\begin{lemma}
\label{lemma-quasi-coherent-weak-serre}
Soit $\mathcal{X}$ un champ algébrique.
\begin{enumerate}
\item $\QCoh(\mathcal{O}_{\mathcal{X}_{lisse,\etale}})$
est une sous-catégorie faible de Serre de
$\textit{Mod}(\mathcal{O}_{\mathcal{X}_{lisse,\etale}})$.
\item $\QCoh(\mathcal{O}_{\mathcal{X}_{flat,fppf}})$
est une sous-catégorie faible de Serre de
$\textit{Mod}(\mathcal{O}_{\mathcal{X}_{flat,fppf}})$.
\end{enumerate}
\end{lemma}

\begin{proof}
Nous allons vérifier les conditions (1), (2), (3) et (4) du
lemme \ref{homology-lemma-characterize-weak-serre-subcategory} d'Homologie.

\medskip\noindent
Comme $0$ est un module quasi-cohérent sur tout site annelé, la condition (1)
est satisfaite.

\medskip\noindent
Par définition, $\QCoh(\mathcal{O})$
est une sous-catégorie strictement pleine de $\textit{Mod}(\mathcal{O})$ ; donc (2) est satisfaite.

\medskip\noindent
Soit $\varphi : \mathcal{G} \to \mathcal{F}$ un morphisme de modules
quasi-cohérents sur $\mathcal{X}_{lisse,\etale}$ ou $\mathcal{X}_{flat,fppf}$.
On a $g^*g_!\mathcal{F} = \mathcal{F}$ et de même pour
$\mathcal{G}$ et $\varphi$ ; voir le lemme \ref{lemma-lisse-etale-modules}.
D'après le lemme \ref{lemma-shriek-quasi-coherent},
$g_!\mathcal{F}$ et $g_!\mathcal{G}$ sont des
$\mathcal{O}_\mathcal{X}$-modules quasi-cohérents. D'après Faisceaux sur les champs, lemme
\ref{stacks-sheaves-lemma-quasi-coherent-algebraic-stack}
le module $\Coker(g_!\varphi)$ est quasi-cohérent
sur $\mathcal{X}$ (et c'est le conoyau dans la catégorie
des modules quasi-cohérents sur $\mathcal{X}$).
Comme $g^*$ est exact (voir le lemme \ref{lemma-lisse-etale}),
$g^*\Coker(g_!\varphi) = \Coker(g^*g_!\varphi) = \Coker(\varphi)$
est également quasi-cohérent (voir le lemme \ref{lemma-quasi-coherent}).
D'après la proposition \ref{proposition-loc-qcoh-flat-base-change},
le noyau $\Ker(g_!\varphi)$ appartient à
$\textit{LQCoh}^{fbc}(\mathcal{O}_\mathcal{X})$.
Comme $g^*$ est exact, on a
$g^*\Ker(g_!\varphi) = \Ker(g^*g_!\varphi) = \Ker(\varphi)$.
Puisque $g^*$ envoie les objets de
$\textit{LQCoh}^{fbc}(\mathcal{O}_\mathcal{X})$ sur des modules
quasi-cohérents d'après le lemme \ref{lemma-quasi-coherent}, on en déduit que
$\Ker(\varphi)$ est également quasi-cohérent. Cela démontre (3).

\medskip\noindent
Enfin, supposons que
$$
0 \to \mathcal{F} \to \mathcal{E} \to \mathcal{G} \to 0
$$
soit une extension de $\mathcal{O}_{\mathcal{X}_{lisse,\etale}}$-modules
(resp.\ de $\mathcal{O}_{\mathcal{X}_{flat,fppf}}$-modules), où $\mathcal{F}$
et $\mathcal{G}$ sont quasi-cohérents. Pour démontrer (4) et achever la preuve,
il nous faut montrer que
$\mathcal{E}$ est quasi-cohérent sur $\mathcal{X}_{lisse,\etale}$
(resp.\ $\mathcal{X}_{flat,fppf}$). Soit $U$
un objet de $\mathcal{X}_{lisse,\etale}$
(resp.\ $\mathcal{X}_{flat,fppf}$) ; nous considérons $U$ comme un schéma lisse
(resp.\ plat) sur $\mathcal{X}$. Il faut montrer que la restriction
de $\mathcal{E}$ à $U_{lisse,\etale}$ (resp.\ $=U_{flat,fppf}$)
est quasi-cohérente. Nous pouvons donc supposer
que $\mathcal{X} = U$ est un schéma. Puisque $\mathcal{G}$
est quasi-cohérent sur $U_{lisse,\etale}$ (resp.\ $U_{flat,fppf}$),
nous pouvons supposer, après avoir remplacé $U$ par les membres
d'un recouvrement étale (resp.\ fppf), que $\mathcal{G}$ admet une
présentation
$$
\bigoplus\nolimits_{j \in J} \mathcal{O} \longrightarrow
\bigoplus\nolimits_{i \in I} \mathcal{O} \longrightarrow
\mathcal{G} \longrightarrow 0
$$
sur $U_{lisse,\etale}$ (resp.\ $U_{flat,fppf}$), où $\mathcal{O}$
est le faisceau structural du site. Nous pouvons aussi supposer $U$ affine.
Comme $\mathcal{F}$ est quasi-cohérent, on a
$$
H^1(U_{lisse,\etale}, \mathcal{F}) = 0,
\quad\text{resp.}\quad
H^1(U_{flat,fppf}, \mathcal{F}) = 0
$$
En effet, $\mathcal{F}$ est l'image inverse d'un module quasi-cohérent
$\mathcal{F}'$ sur le grand site de $U$
(d'après le lemme \ref{lemma-quasi-coherent}), les cohomologies
de $\mathcal{F}$ et $\mathcal{F}'$ coïncident (d'après le
lemme \ref{lemma-lisse-etale-cohomology}), et l'on sait que
la cohomologie de $\mathcal{F}'$ sur le grand site du schéma affine $U$
est nulle (pour l'obtenir dans la situation présente, il faut
combiner Descentee, propositions
\ref{descent-proposition-equivalence-quasi-coherent} et
\ref{descent-proposition-same-cohomology-quasi-coherent} avec
Cohomologie des schémas, lemme
\ref{coherent-lemma-quasi-coherent-affine-cohomology-zero}).
Nous pouvons donc relever le morphisme
$\bigoplus_{i \in I} \mathcal{O} \to \mathcal{G}$
à $\mathcal{E}$. Une chasse au diagramme donne alors
une suite exacte
$$
\bigoplus\nolimits_{j \in J} \mathcal{O} \to
\mathcal{F} \oplus \bigoplus\nolimits_{i \in I} \mathcal{O}
\to
\mathcal{E} \to 0
$$
D'après (3), démontré ci-dessus, on en conclut que $\mathcal{E}$ est
quasi-cohérent, comme voulu.
\end{proof}









\section{Faisceaux cohérents sur les champs localement noethériens}
\sectionmark{Faisceaux cohérents sur les champs noethériens}
\label{section-coherent-sheaves}

\noindent
Cette section est l'analogue de
Cohomologie des espaces, section
\ref{spaces-cohomology-section-coherent}.
Nous avons défini la notion de module cohérent sur tout topos annelé dans
Modules sur les sites, section \ref{sites-modules-section-local}.
Cependant, pour tout champ algébrique $\mathcal{X}$, la catégorie des
$\mathcal{O}_\mathcal{X}$-modules cohérents est nulle, essentiellement parce que
le site $\mathcal{X}$ contient trop d'objets non noethériens (même si
$\mathcal{X}$ est lui-même localement noethérien). Nous définirons plutôt les
modules cohérents au moyen du lemme suivant.

\begin{lemma}
\label{lemma-coherent-Noetherian}
Soit $\mathcal{X}$ un champ algébrique localement noethérien.
Soit $\mathcal{F}$ un $\mathcal{O}_\mathcal{X}$-module.
Les conditions suivantes sont équivalentes :
\begin{enumerate}
\item $\mathcal{F}$ est un $\mathcal{O}_\mathcal{X}$-module
quasi-cohérent de type fini ;
\item $\mathcal{F}$ est un $\mathcal{O}_\mathcal{X}$-module
de présentation finie ;
\item $\mathcal{F}$ est quasi-cohérent et, pour tout morphisme
$f : U \to \mathcal{X}$ où $U$ est un espace algébrique localement noethérien,
l'image inverse $f^*\mathcal{F}|_{U_\etale}$ est cohérente ;
\item $\mathcal{F}$ est quasi-cohérent et il existe un espace algébrique
$U$ et un morphisme $f : U \to \mathcal{X}$ localement de type fini,
plat et surjectif, tels que l'image inverse $f^*\mathcal{F}|_{U_\etale}$
soit cohérente.
\end{enumerate}
\end{lemma}

\begin{proof}
Soit $f : U \to \mathcal{X}$ comme dans (4).
Alors $U$ est localement noethérien (Morphismes de champs, lemme
\ref{stacks-morphisms-lemma-locally-finite-type-locally-noetherian})
et l'on voit que l'énoncé du lemme a un sens. De plus,
$f$ est localement de présentation finie d'après Morphismes de champs, lemme
\ref{stacks-morphisms-lemma-noetherian-finite-type-finite-presentation}.
Soit $x$ un objet de $\mathcal{X}$ au-dessus du schéma $V$.
Pour démontrer (2), il faut montrer qu'après avoir remplacé $V$
par les membres d'un recouvrement fppf de $V$, la restriction $x^*\mathcal{F}$
admet une présentation globale finie sur $\mathcal{X}/x \cong (\Sch/V)_{fppf}$.
La projection $W = U \times_\mathcal{X} V \to V$ est
localement de présentation finie, plate et surjective.
On peut donc remplacer $V$ par les membres d'un recouvrement étale de $W$
par des schémas et supposer qu'il existe un morphisme $h : V \to U$ tel que $f \circ h = x$.
Comme $\mathcal{F}$ est quasi-cohérent, on voit que
la restriction $x^*\mathcal{F}$
est l'image inverse de $h_{small}^*(f^*\mathcal{F})|_{U_\etale}$
par $\pi_V$ ; voir Faisceaux sur les champs, lemme
\ref{stacks-sheaves-lemma-compare-quasi-coherent}.
Puisque, par hypothèse, $f^*\mathcal{F}|_{U_\etale}$ admet localement pour la
topologie étale une présentation finie, on en conclut
(4) $\Rightarrow$ (2).

\medskip\noindent
La condition (2) implique (1) pour tout topos annelé (cela résulte aussitôt de la définition).
Les propriétés « de type fini » et « quasi-cohérent »
sont préservées par l'image inverse de tout morphisme de topos annelés ; voir
Modules sur les sites, lemme \ref{sites-modules-lemma-local-pullback}.
Par conséquent, (1) implique (3) ; voir Cohomologie des espaces, lemme
\ref{spaces-cohomology-lemma-coherent-Noetherian}.
Enfin, (3) implique trivialement (4).
\end{proof}

\begin{definition}
\label{definition-coherent}
Soit $\mathcal{X}$ un champ algébrique localement noethérien.
Un $\mathcal{O}_\mathcal{X}$-module $\mathcal{F}$ est dit {\it cohérent}
si $\mathcal{F}$ satisfait l'une (et donc chacune) des conditions
équivalentes du lemme \ref{lemma-coherent-Noetherian}.
La catégorie des $\mathcal{O}_\mathcal{X}$-modules cohérents est
notée $\textit{Coh}(\mathcal{O}_\mathcal{X})$.
\end{definition}

\begin{lemma}
\label{lemma-elementary-coherent}
Soit $\mathcal{X}$ un champ algébrique localement noethérien.
Le module $\mathcal{O}_\mathcal{X}$ est cohérent, tout
$\mathcal{O}_\mathcal{X}$-module inversible est cohérent et, plus généralement, tout
$\mathcal{O}_\mathcal{X}$-module localement libre de type fini est cohérent.
\end{lemma}

\begin{proof}
Cela résulte de la définition et de
Cohomologie des espaces, lemme
\ref{spaces-cohomology-lemma-coherent-Noetherian}.
\end{proof}

\begin{lemma}
\label{lemma-pullback-coherent}
Soit $f : \mathcal{X} \to \mathcal{Y}$ un morphisme de champs algébriques
localement noethériens. Alors $f^*$ envoie les modules cohérents
sur $\mathcal{Y}$ sur des modules cohérents sur $\mathcal{X}$.
\end{lemma}

\begin{proof}
Cela résulte aussitôt de la définition et du fait que l'image inverse par tout
morphisme de topos annelés préserve les modules de présentation finie ; voir
Modules sur les sites, lemme \ref{sites-modules-lemma-local-pullback}.
\end{proof}

\begin{lemma}
\label{lemma-coherent-abelian-Noetherian}
Soit $\mathcal{X}$ un champ algébrique localement noethérien.
La catégorie des $\mathcal{O}_\mathcal{X}$-modules cohérents est abélienne.
Si $\varphi : \mathcal{F} \to \mathcal{G}$ est un morphisme
de $\mathcal{O}_\mathcal{X}$-modules cohérents, alors :
\begin{enumerate}
\item le conoyau $\Coker(\varphi)$ calculé dans
$\textit{Mod}(\mathcal{O}_\mathcal{X})$ est un
$\mathcal{O}_\mathcal{X}$-module cohérent ;
\item l'image $\Im(\varphi)$ calculée dans
$\textit{Mod}(\mathcal{O}_\mathcal{X})$ est un
$\mathcal{O}_\mathcal{X}$-module cohérent ;
\item le noyau $\Ker(\varphi)$ calculé dans
$\textit{Mod}(\mathcal{O}_\mathcal{X})$
peut ne pas être cohérent, mais il appartient à
$\textit{LQCoh}^{fbc}(\mathcal{O}_\mathcal{X})$ et $Q(\Ker(\varphi))$
est cohérent et constitue le noyau de $\varphi$ dans
$\textit{Coh}(\mathcal{O}_\mathcal{X})$.
\end{enumerate}
Le foncteur d'inclusion $\textit{Coh}(\mathcal{O}_\mathcal{X}) \to
\QCoh(\mathcal{O}_\mathcal{X})$ est exact.
\end{lemma}

\begin{proof}
Les règles données pour former les noyaux, images et conoyaux dans
$\textit{Coh}(\mathcal{O}_\mathcal{X})$ coïncident avec la prescription
pour les modules quasi-cohérents de la remarque \ref{remark-QCoh-abelian}.
Le lemme résultera donc du fait que les modules quasi-cohérents
$\Coker(\varphi)$, $\Im(\varphi)$ et $Q(\Ker(\varphi))$ sont cohérents.
D'après le lemme \ref{lemma-coherent-Noetherian}, il suffit de le montrer après
restriction à $U_\etale$ pour un morphisme lisse surjectif
$f : U \to \mathcal{X}$. Le foncteur
$\mathcal{F} \mapsto f^*\mathcal{F}|_{U_\etale}$
est exact. Par conséquent, $f^*\Coker(\varphi)$ et $f^*\Im(\varphi)$ sont
le conoyau et l'image d'un morphisme entre $\mathcal{O}_U$-modules cohérents,
donc sont cohérents comme voulu. Le foncteur
$\mathcal{F} \mapsto f^*\mathcal{F}|_{U_\etale}$
annule les modules parasites d'après le lemme \ref{lemma-parasitic}.
Ainsi, $f^*Q(\Ker(\varphi))|_{U_\etale} = f^*\Ker(\varphi)|_{U_\etale}$
d'après la partie (2) du lemme \ref{lemma-adjoint-kernel-parasitic}.
On conclut donc de même que $Q(\Ker(\varphi))$ est cohérent.
\end{proof}

\begin{lemma}
\label{lemma-extension-coherent}
Soit $\mathcal{X}$ un champ algébrique localement noethérien.
Étant donnée une suite exacte courte
$0 \to \mathcal{F}_1 \to \mathcal{F}_2 \to \mathcal{F}_3 \to 0$
dans $\textit{Mod}(\mathcal{O}_\mathcal{X})$,
si $\mathcal{F}_1$ et $\mathcal{F}_3$ sont cohérents, alors
$\mathcal{F}_2$ est cohérent.
\end{lemma}

\begin{proof}
D'après Faisceaux sur les champs, lemme
\ref{stacks-sheaves-lemma-quasi-coherent-algebraic-stack}, partie (7),
$\mathcal{F}_2$ est quasi-cohérent. On peut alors vérifier que
$\mathcal{F}_2$ est cohérent par restriction à $U_\etale$
pour un morphisme $U \to \mathcal{X}$ lisse et surjectif.
Cela résulte de Cohomologie des espaces, lemme
\ref{spaces-cohomology-lemma-coherent-abelian-Noetherian}.
Quelques détails sont omis.
\end{proof}

\noindent
Les modules cohérents forment une sous-catégorie de Serre de la
catégorie des $\mathcal{O}_\mathcal{X}$-modules quasi-cohérents.
Cela n'est pas vrai pour les modules sur un topos annelé général.

\begin{lemma}
\label{lemma-coherent-Noetherian-quasi-coherent-sub-quotient}
Soit $\mathcal{X}$ un champ algébrique localement noethérien.
Alors $\textit{Coh}(\mathcal{O}_\mathcal{X})$ est une sous-catégorie de Serre de
$\QCoh(\mathcal{O}_\mathcal{X})$. Soit $\varphi : \mathcal{F} \to \mathcal{G}$
un morphisme de $\mathcal{O}_\mathcal{X}$-modules quasi-cohérents. On a :
\begin{enumerate}
\item si $\mathcal{F}$ est cohérent et $\varphi$ est surjectif,
alors $\mathcal{G}$ est cohérent ;
\item si $\mathcal{F}$ est cohérent, alors $\Im(\varphi)$ est cohérente ;
\item si $\mathcal{G}$ est cohérent et $\Ker(\varphi)$ est parasite, alors
$\mathcal{F}$ est cohérent.
\end{enumerate}
\end{lemma}

\begin{proof}
Choisissons un schéma $U$ et un morphisme lisse surjectif $f : U \to \mathcal{X}$.
Alors le foncteur
$f^* : \QCoh(\mathcal{O}_\mathcal{X}) \to \QCoh(\mathcal{O}_U)$
est exact (lemme \ref{lemma-flat-pullback-quasi-coherent})
et, en outre, $\textit{Coh}(\mathcal{O}_\mathcal{X})$ est par définition
la sous-catégorie pleine de $\QCoh(\mathcal{O}_\mathcal{X})$ formée
des objets $\mathcal{F}$ tels que $f^*\mathcal{F}$ appartienne à
$\textit{Coh}(\mathcal{O}_U)$. Le fait que
$\textit{Coh}(\mathcal{O}_\mathcal{X})$ soit une sous-catégorie de Serre de
$\QCoh(\mathcal{O}_\mathcal{X})$ résulte aussitôt de ceci
et du fait correspondant pour $U$ ; voir
Cohomologie des espaces, lemmes
\ref{spaces-cohomology-lemma-coherent-abelian-Noetherian} et
\ref{spaces-cohomology-lemma-coherent-Noetherian-quasi-coherent-sub-quotient}.
Nous omettons la preuve de (1), (2) et (3). Indication : comparer
avec la démonstration du lemme \ref{lemma-coherent-abelian-Noetherian}.
\end{proof}

\noindent
Soit $\mathcal{X}$ un champ algébrique localement noethérien.
Soit $U$ un espace algébrique et soit $f : U \to \mathcal{X}$
un morphisme surjectif, localement de présentation finie et plat.
Observons que $U$ est localement noethérien
(Morphismes de champs, lemme
\ref{stacks-morphisms-lemma-locally-finite-type-locally-noetherian}).
Soit $(U, R, s, t, c)$ le groupoïde en espaces algébriques
et $f_{can} : [U/R] \to \mathcal{X}$ l'isomorphisme
construit dans
Champs algébriques, lemme \ref{algebraic-lemma-map-space-into-stack} et
remarque \ref{algebraic-remark-flat-fp-presentation}.
Comme dans Faisceaux sur les champs, section
\ref{stacks-sheaves-section-quasi-coherent-algebraic-stacks}
on obtient des équivalences
$$
\QCoh(\mathcal{O}_\mathcal{X})
\cong
\QCoh(\mathcal{O}_{[U/R]})
\cong
\QCoh(U, R, s, t, c)
$$
où la seconde équivalence est celle de
Faisceaux sur les champs, proposition \ref{stacks-sheaves-proposition-quasi-coherent}.
Rappelons que, dans Groupoïdes dans les espaces, section
\ref{spaces-groupoids-section-colimits}
nous avons défini la sous-catégorie pleine
$$
\textit{Coh}(U, R, s, t, c) \subset \QCoh(U, R, s, t, c)
$$
des {\it modules cohérents} comme celle des $(\mathcal{G}, \alpha)$ tels que
$\mathcal{G}$ soit un $\mathcal{O}_U$-module cohérent.

\begin{lemma}
\label{lemma-coherent-presentation}
Dans la situation décrite ci-dessus, l'équivalence
$\QCoh(\mathcal{O}_\mathcal{X}) \cong \QCoh(U, R, s, t, c)$
envoie les modules cohérents sur les modules cohérents et réciproquement, c'est-à-dire
qu'elle induit une équivalence
$\textit{Coh}(\mathcal{O}_\mathcal{X}) \cong \textit{Coh}(U, R, s, t, c)$.
\end{lemma}

\begin{proof}
Cela résulte aussitôt de la définition des
$\mathcal{O}_\mathcal{X}$-modules cohérents.
Pour expliciter les références, ce qui précède utilise
Morphismes de champs, lemme
\ref{stacks-morphisms-lemma-locally-finite-type-locally-noetherian},
Champs algébriques, lemme \ref{algebraic-lemma-map-space-into-stack} et
remarque \ref{algebraic-remark-flat-fp-presentation},
Faisceaux sur les champs, section
\ref{stacks-sheaves-section-quasi-coherent-algebraic-stacks},
Faisceaux sur les champs, proposition \ref{stacks-sheaves-proposition-quasi-coherent},
et Groupoïdes dans les espaces, section
\ref{spaces-groupoids-section-colimits}.
\end{proof}

\begin{lemma}
\label{lemma-coherent-hom}
Soit $\mathcal{X}$ un champ algébrique localement noethérien. Soient $\mathcal{F}$
et $\mathcal{G}$ des $\mathcal{O}_\mathcal{X}$-modules cohérents. Alors
le Hom interne $hom(\mathcal{F}, \mathcal{G})$
construit dans le lemme \ref{lemma-internal-hom-fp-into-qcoh}
est un $\mathcal{O}_\mathcal{X}$-module cohérent.
\end{lemma}

\begin{proof}
Soit $U \to \mathcal{X}$ un morphisme lisse surjectif issu d'un schéma.
D'après le point (\ref{item-hom-restriction}) de la section \ref{section-further-remarks},
la restriction de
$hom(\mathcal{F}, \mathcal{G})$ à $U$ est le faisceau Hom des restrictions.
Ce lemme résulte donc du cas des espaces algébriques ; voir
Cohomologie des espaces, lemme \ref{spaces-cohomology-lemma-tensor-hom-coherent}.
\end{proof}






\section{Faisceaux cohérents sur les champs noethériens}
\label{section-coherent-on-noetherian}

\noindent
Cette section est l'analogue de
Cohomologie des espaces, section
\ref{spaces-cohomology-section-coherent-quasi-compact}.

\begin{lemma}
\label{lemma-directed-colimit-coherent}
Soit $\mathcal{X}$ un champ algébrique noethérien. Tout
$\mathcal{O}_\mathcal{X}$-module quasi-cohérent est la colimite filtrante
de ses sous-modules cohérents.
\end{lemma}

\begin{proof}
Soit $\mathcal{F}$ un $\mathcal{O}_\mathcal{X}$-module quasi-cohérent.
Si $\mathcal{G}, \mathcal{H} \subset \mathcal{F}$ sont des sous-modules
$\mathcal{O}_\mathcal{X}$-cohérents, alors l'image de
$\mathcal{G} \oplus \mathcal{H} \to \mathcal{F}$ est encore un
sous-module $\mathcal{O}_\mathcal{X}$-cohérent qui les contient tous deux,
voir le lemme \ref{lemma-coherent-Noetherian-quasi-coherent-sub-quotient}.
Le système est donc filtrant.
Il suffit dès lors de montrer que $\mathcal{F}$ s'écrit comme
colimite filtrante de modules cohérents : on pourra alors prendre les
images de ces modules dans $\mathcal{F}$ et conclure qu'il y en a
suffisamment.

\medskip\noindent
Soient $U$ un schéma affine et $U \to \mathcal{X}$
un morphisme lisse surjectif (Propriétés des champs, lemme
\ref{stacks-properties-lemma-quasi-compact-stack}).
Posons $R = U \times_\mathcal{X} U$, de sorte que $\mathcal{X} = [U/R]$
comme dans Champs algébriques, lemme \ref{algebraic-lemma-stack-presentation}.
D'après le lemme \ref{lemma-coherent-presentation}, on a
$\QCoh(\mathcal{O}_X) = \QCoh(U, R, s, t, c)$ et
$\textit{Coh}(\mathcal{O}_X) = \textit{Coh}(U, R, s, t, c)$.
On est ainsi ramené à établir l'énoncé correspondant pour
$\QCoh(U, R, s, t, c)$. C'est le contenu de
Groupoïdes dans les espaces, lemme \ref{spaces-groupoids-lemma-colimit-coherent} ;
nous en vérifions les hypothèses au paragraphe suivant.

\medskip\noindent
Nous invitons le lecteur à passer le reste de la démonstration.
Le schéma affine $U$ est noethérien ; cela résulte de notre définition
d'un champ $\mathcal{X}$ localement noethérien, voir
Propriétés des champs, définition
\ref{stacks-properties-definition-type-property} et
remarque \ref{stacks-properties-remark-list-properties-local-smooth-topology}.
Les morphismes de projection $s, t : R \to U$ sont lisses (voir la référence
ci-dessus),
quasi-séparés et quasi-compacts (Morphismes de champs, lemme
\ref{stacks-morphisms-lemma-quasi-compact-quasi-separated-permanence}).
En particulier, $R$ est un espace algébrique quasi-compact et quasi-séparé,
lisse sur $U$, donc noethérien
(Morphismes d'espaces, lemme
\ref{spaces-morphisms-lemma-finite-presentation-noetherian}).
\end{proof}







\begin{multicols}{2}[\section{Autres chapitres}]
\noindent
Préliminaires
\begin{enumerate}
\item \hyperref[introduction-section-phantom]{Introduction}
\item \hyperref[conventions-section-phantom]{Conventions}
\item \hyperref[sets-section-phantom]{Théorie des ensembles}
\item \hyperref[categories-section-phantom]{Catégories}
\item \hyperref[topology-section-phantom]{Topologie}
\item \hyperref[sheaves-section-phantom]{Faisceaux sur les espaces}
\item \hyperref[sites-section-phantom]{Sites et faisceaux}
\item \hyperref[stacks-section-phantom]{Champs}
\item \hyperref[fields-section-phantom]{Corps}
\item \hyperref[algebra-section-phantom]{Algèbre commutative}
\item \hyperref[brauer-section-phantom]{Groupes de Brauer}
\item \hyperref[homology-section-phantom]{Algèbre homologique}
\item \hyperref[derived-section-phantom]{Catégories dérivées}
\item \hyperref[simplicial-section-phantom]{Méthodes simpliciales}
\item \hyperref[more-algebra-section-phantom]{Compléments d'algèbre}
\item \hyperref[smoothing-section-phantom]{Lissage des morphismes d'anneaux}
\item \hyperref[modules-section-phantom]{Faisceaux de modules}
\item \hyperref[sites-modules-section-phantom]{Modules sur les sites}
\item \hyperref[injectives-section-phantom]{Objets injectifs}
\item \hyperref[cohomology-section-phantom]{Cohomologie des faisceaux}
\item \hyperref[sites-cohomology-section-phantom]{Cohomologie sur les sites}
\item \hyperref[dga-section-phantom]{Algèbre différentielle graduée}
\item \hyperref[dpa-section-phantom]{Algèbre à puissances divisées}
\item \hyperref[sdga-section-phantom]{Faisceaux différentiels gradués}
\item \hyperref[hypercovering-section-phantom]{Hyperrecouvrements}
\end{enumerate}
Schémas
\begin{enumerate}
\setcounter{enumi}{25}
\item \hyperref[schemes-section-phantom]{Schémas}
\item \hyperref[constructions-section-phantom]{Constructions de schémas}
\item \hyperref[properties-section-phantom]{Propriétés des schémas}
\item \hyperref[morphisms-section-phantom]{Morphismes de schémas}
\item \hyperref[coherent-section-phantom]{Cohomologie des schémas}
\item \hyperref[divisors-section-phantom]{Diviseurs}
\item \hyperref[limits-section-phantom]{Limites de schémas}
\item \hyperref[varieties-section-phantom]{Variétés}
\item \hyperref[topologies-section-phantom]{Topologies sur les schémas}
\item \hyperref[descent-section-phantom]{Descente}
\item \hyperref[perfect-section-phantom]{Catégories dérivées de schémas}
\item \hyperref[more-morphisms-section-phantom]{Compléments sur les morphismes}
\item \hyperref[flat-section-phantom]{Compléments sur la platitude}
\item \hyperref[groupoids-section-phantom]{Schémas en groupoïdes}
\item \hyperref[more-groupoids-section-phantom]{Compléments sur les schémas en groupoïdes}
\item \hyperref[etale-section-phantom]{Morphismes étales de schémas}
\end{enumerate}
Thèmes de théorie des schémas
\begin{enumerate}
\setcounter{enumi}{41}
\item \hyperref[chow-section-phantom]{Homologie de Chow}
\item \hyperref[intersection-section-phantom]{Théorie de l'intersection}
\item \hyperref[pic-section-phantom]{Schémas de Picard des courbes}
\item \hyperref[weil-section-phantom]{Théories cohomologiques de Weil}
\item \hyperref[adequate-section-phantom]{Modules adéquats}
\item \hyperref[dualizing-section-phantom]{Complexes dualisants}
\item \hyperref[duality-section-phantom]{Dualité pour les schémas}
\item \hyperref[discriminant-section-phantom]{Discriminants et différentes}
\item \hyperref[derham-section-phantom]{Cohomologie de de Rham}
\item \hyperref[local-cohomology-section-phantom]{Cohomologie locale}
\item \hyperref[algebraization-section-phantom]{Géométrie algébrique et formelle}
\item \hyperref[curves-section-phantom]{Courbes algébriques}
\item \hyperref[resolve-section-phantom]{Résolution des surfaces}
\item \hyperref[models-section-phantom]{Réduction semi-stable}
\item \hyperref[functors-section-phantom]{Foncteurs et morphismes}
\item \hyperref[equiv-section-phantom]{Catégories dérivées de variétés}
\item \hyperref[pione-section-phantom]{Groupes fondamentaux de schémas}
\item \hyperref[etale-cohomology-section-phantom]{Cohomologie étale}
\item \hyperref[crystalline-section-phantom]{Cohomologie cristalline}
\item \hyperref[proetale-section-phantom]{Cohomologie pro-étale}
\item \hyperref[relative-cycles-section-phantom]{Cycles relatifs}
\item \hyperref[more-etale-section-phantom]{Compléments de cohomologie étale}
\item \hyperref[trace-section-phantom]{La formule des traces}
\end{enumerate}
Espaces algébriques
\begin{enumerate}
\setcounter{enumi}{64}
\item \hyperref[spaces-section-phantom]{Espaces algébriques}
\item \hyperref[spaces-properties-section-phantom]{Propriétés des espaces algébriques}
\item \hyperref[spaces-morphisms-section-phantom]{Morphismes d'espaces algébriques}
\item \hyperref[decent-spaces-section-phantom]{Espaces algébriques décents}
\item \hyperref[spaces-cohomology-section-phantom]{Cohomologie des espaces algébriques}
\item \hyperref[spaces-limits-section-phantom]{Limites d'espaces algébriques}
\item \hyperref[spaces-divisors-section-phantom]{Diviseurs sur les espaces algébriques}
\item \hyperref[spaces-over-fields-section-phantom]{Espaces algébriques sur des corps}
\item \hyperref[spaces-topologies-section-phantom]{Topologies sur les espaces algébriques}
\item \hyperref[spaces-descent-section-phantom]{Descente et espaces algébriques}
\item \hyperref[spaces-perfect-section-phantom]{Catégories dérivées d'espaces}
\item \hyperref[spaces-more-morphisms-section-phantom]{Compléments sur les morphismes d'espaces}
\item \hyperref[spaces-flat-section-phantom]{Platitude sur les espaces algébriques}
\item \hyperref[spaces-groupoids-section-phantom]{Groupoïdes dans les espaces algébriques}
\item \hyperref[spaces-more-groupoids-section-phantom]{Compléments sur les groupoïdes dans les espaces}
\item \hyperref[bootstrap-section-phantom]{Amorçage}
\item \hyperref[spaces-pushouts-section-phantom]{Sommes amalgamées d'espaces algébriques}
\end{enumerate}
Thèmes de géométrie
\begin{enumerate}
\setcounter{enumi}{81}
\item \hyperref[spaces-chow-section-phantom]{Groupes de Chow des espaces}
\item \hyperref[groupoids-quotients-section-phantom]{Quotients de groupoïdes}
\item \hyperref[spaces-more-cohomology-section-phantom]{Compléments sur la cohomologie des espaces}
\item \hyperref[spaces-simplicial-section-phantom]{Espaces simpliciaux}
\item \hyperref[spaces-duality-section-phantom]{Dualité pour les espaces}
\item \hyperref[formal-spaces-section-phantom]{Espaces algébriques formels}
\item \hyperref[restricted-section-phantom]{Algébrisation des espaces formels}
\item \hyperref[spaces-resolve-section-phantom]{Résolution des surfaces revisitée}
\end{enumerate}
Théorie des déformations
\begin{enumerate}
\setcounter{enumi}{89}
\item \hyperref[formal-defos-section-phantom]{Théorie formelle des déformations}
\item \hyperref[defos-section-phantom]{Théorie des déformations}
\item \hyperref[cotangent-section-phantom]{Le complexe cotangent}
\item \hyperref[examples-defos-section-phantom]{Problèmes de déformation}
\end{enumerate}
Champs algébriques
\begin{enumerate}
\setcounter{enumi}{93}
\item \hyperref[algebraic-section-phantom]{Champs algébriques}
\item \hyperref[examples-stacks-section-phantom]{Exemples de champs}
\item \hyperref[stacks-sheaves-section-phantom]{Faisceaux sur les champs algébriques}
\item \hyperref[criteria-section-phantom]{Critères de représentabilité}
\item \hyperref[artin-section-phantom]{Axiomes d'Artin}
\item \hyperref[quot-section-phantom]{Espaces de Quot et de Hilbert}
\item \hyperref[stacks-properties-section-phantom]{Propriétés des champs algébriques}
\item \hyperref[stacks-morphisms-section-phantom]{Morphismes de champs algébriques}
\item \hyperref[stacks-limits-section-phantom]{Limites de champs algébriques}
\item \hyperref[stacks-cohomology-section-phantom]{Cohomologie des champs algébriques}
\item \hyperref[stacks-perfect-section-phantom]{Catégories dérivées de champs}
\item \hyperref[stacks-introduction-section-phantom]{Introduction aux champs algébriques}
\item \hyperref[stacks-more-morphisms-section-phantom]{Compléments sur les morphismes de champs}
\item \hyperref[stacks-geometry-section-phantom]{La géométrie des champs}
\end{enumerate}
Thèmes de théorie des espaces de modules
\begin{enumerate}
\setcounter{enumi}{107}
\item \hyperref[moduli-section-phantom]{Champs de modules}
\item \hyperref[moduli-curves-section-phantom]{Espaces de modules de courbes}
\end{enumerate}
Divers
\begin{enumerate}
\setcounter{enumi}{109}
\item \hyperref[examples-section-phantom]{Exemples}
\item \hyperref[exercises-section-phantom]{Exercices}
\item \hyperref[guide-section-phantom]{Guide bibliographique}
\item \hyperref[desirables-section-phantom]{Desiderata}
\item \hyperref[coding-section-phantom]{Style de programmation}
\item \hyperref[obsolete-section-phantom]{Obsolète}
\item \hyperref[fdl-section-phantom]{Licence de documentation libre GNU}
\item \hyperref[index-section-phantom]{Index généré automatiquement}
\end{enumerate}
\end{multicols}


\bibliography{my}
\bibliographystyle{amsalpha}

\end{document}
